% Generated by tools/make_body.py; edit generator, not this file.
% UNPROOFREAD OCR TRANSCRIPTION: page-image verification remains required.
\chapter{Complex Integration}\label{chap:4}

% Original PDF physical page 117
\sourcepage{117}
\setcounter{section}{0}\section{Fundamental Theorems}\label{sec:4.1}
Many important properties of analytic functions are very difficult to prove without use of complex integration\index{Complex integration|(}\index{Integration|(}\index{Complex integration}\index{Integration}. For instance, it is only recently that it became possible to prove, without resorting to complex integrals\index{Integral|(}\index{Integral} or equivalent tools, that the derivative of an analytic function is continuous, or that the higher derivatives exist. At present the integration-free proofs are, to say the least, much more difficult than the classical proofs.\hyperlink{fn:117:dagger}{\textsuperscript{\ensuremath{\dagger}}}

As in the real case we distinguish between definite and indefinite integrals. An indefinite integral is a function whose derivative equals a given analytic function in a region; in many elementary cases indefinite integrals can be found by inversion of known derivation formulas. The definite integrals\index{Definite integral} are taken over differentiable or piecewise differentiable arcs and are not limited to analytic functions. They can be defined by a limit process which mimics the definition of a real definite integral. Actually, we shall prefer to define complex definite integrals in terms of real integrals. This will save us from repeating existence proofs which are essentially the same as in the real case. Naturally, the reader must be thoroughly familiar with the theory of definite integrals of real continuous functions.

\setcounter{section}{1}\setcounter{subsection}{0}\subsection{Line Integrals}\index{Line integral|(}\index{Line integral}\label{sec:4.1.1}
The most immediate generalization of a real integral is to the definite integral of a complex function over a real interval. If $f(t) = u(t) + iv(t)$ is a continuous function,
\begingroup\renewcommand{\thefootnote}{}\footnotetext{\hypertarget{fn:117:dagger}{}\textsuperscript{\ensuremath{\dagger}} Without use of integration R. L. Plunkett\index{Plunkett R L@Plunkett, R. L.} proved the continuity of the derivative (\emph{Bull. Am. Math. Soc.} 65, 1959). E. H. Connell\index{Connell E H@Connell, E. H.} and P. Porcelli\index{Porcelli P@Porcelli, P.|hyperpagenote} proved the existence of all derivatives (\emph{Bull. Am. Math. Soc.} 67, 1961). Both proofs lean on a topological theorem due to G. T. Whyburn\index{Whyburn G T@Whyburn, G. T.|hyperpagenote}.}\endgroup
% Original PDF physical page 118
\sourcepage{118}
defined in an interval $(a,b)$, we set by definition
\begin{equation*}
\int_ {a} ^ {b} f (t) d t = \int_ {a} ^ {b} u (t) d t + i \int_ {a} ^ {b} v (t) d t.\tag{1}\label{eq:4.1}
\end{equation*}
This integral\index{Integral} has most of the properties of the real integral. In particular, if $c = \alpha + i\beta$ is a complex constant we obtain
\begin{equation*}
\int_ {a} ^ {b} c f (t) d t = c \int_ {a} ^ {b} f (t) d t,\tag{2}\label{eq:4.2}
\end{equation*}
for both members are equal to
\begin{equation*}
\int_ {a} ^ {b} (\alpha u - \beta v) d t + i \int_ {a} ^ {b} (\alpha v + \beta u) d t.
\end{equation*}
When $a \leq b$, the fundamental inequality
\begin{equation*}
\left| \int_ {a} ^ {b} f (t) d t \right| \leq \int_ {a} ^ {b} | f (t) | d t\tag{3}\label{eq:4.3}
\end{equation*}
holds for arbitrary complex $f(t)$. To see this we choose $c = e^{-i\theta}$ with a real $\theta$ in \hyperref[eq:4.2]{(2)} and find
\begin{equation*}
\operatorname{Re} \left[ e ^ {- i \theta} \int_ {a} ^ {b} f (t) d t \right] = \int_ {a} ^ {b} \operatorname{Re} \left[ e ^ {- i \theta} f (t) \right] d t \leq \int_ {a} ^ {b} | f (t) | d t.
\end{equation*}
For $\theta = \arg \int_{a}^{b} f(t) dt$ the expression on the left reduces to the absolute value of the integral, and \hyperref[eq:4.3]{(3)} results.\hyperlink{fn:118:dagger}{\textsuperscript{\ensuremath{\dagger}}}

We consider now a piecewise differentiable arc $\gamma$ with the equation $z = z(t)$, $a \leq t \leq b$. If the function $f(z)$ is defined and continuous on $\gamma$, then $f(z(t))$ is also continuous and we can set
\begin{equation*}
\int_ {\gamma} f (z) d z = \int_ {a} ^ {b} f (z (t)) z ^ {\prime} (t) d t.\tag{4}\label{eq:4.4}
\end{equation*}
This is our definition of the complex line integral\index{Line integral} of $f(z)$ extended over the arc $\gamma$. In the right-hand member of \hyperref[eq:4.4]{(4)}, if $z'(t)$ is not continuous throughout, the interval of integration\index{Integration} has to be subdivided in the obvious manner. Whenever a line integral over an arc $\gamma$ is considered, let it be tacitly understood that $\gamma$ is piecewise differentiable.

The most important property of the integral \hyperref[eq:4.4]{(4)} is its invariance under a change of parameter. A change of parameter is determined by an increasing function $t = t(\tau)$ which maps an interval $\alpha \leq \tau \leq \beta$ onto $a \leq t \leq b$; we assume that $t(\tau)$ is piecewise differentiable. By the rule

\begingroup\renewcommand{\thefootnote}{}\footnotetext{\hypertarget{fn:118:dagger}{}\textsuperscript{\ensuremath{\dagger}} $\theta$ is not defined if $\int_{a}^{b} f dt = 0$, but then there is nothing to prove.}\endgroup

% Original PDF physical page 119
\sourcepage{119}
for changing the variable of integration\index{Integration} we have
\begin{equation*}
\int_ {a} ^ {b} f (z (t)) z ^ {\prime} (t) d t = \int_ {\alpha} ^ {\beta} f (z (t (\tau))) z ^ {\prime} (t (\tau)) t ^ {\prime} (\tau) d \tau .
\end{equation*}
But $z'(t(\tau))t'(\tau)$ is the derivative of $z(t(\tau))$ with respect to $\tau$, and hence the integral\index{Integral} \hyperref[eq:4.4]{(4)} has the same value whether $\gamma$ be represented by the equation $z = z(t)$ or by the equation $z = z(t(\tau))$.

In Chap.~\hyperref[chap:3]{3}, Sec.~\hyperref[sec:3.2.1]{2.1}, we defined the opposite arc $-\gamma$ by the equation $z = z(-t)$, $-b \leq t \leq -a$. We have thus
\begin{equation*}
\int_ {- \gamma} f (z) d z = \int_ {- b} ^ {- a} f (z (- t)) (- z ^ {\prime} (- t)) d t,
\end{equation*}
and by a change of variable the last integral can be brought to the form
\begin{equation*}
\int_ {b} ^ {a} f (z (t)) z ^ {\prime} (t) d t.
\end{equation*}
We conclude that
\begin{equation*}
\int_ {- \gamma} f (z) d z = - \int_ {\gamma} f (z) d z.\tag{5}\label{eq:4.5}
\end{equation*}
The integral \hyperref[eq:4.4]{(4)} has also a very obvious additive property. It is quite clear what is meant by subdividing an arc $\gamma$ into a finite number of subarcs. A subdivision can be indicated by a symbolic equation
\begin{equation*}
\gamma = \gamma_ {1} + \gamma_ {2} + \dots + \gamma_ {n},
\end{equation*}
and the corresponding integrals satisfy the relation
\begin{equation*}
\int_ {\gamma_ {1} + \gamma_ {2} + \dots + \gamma_ {n}} f d z = \int_ {\gamma_ {1}} f d z + \int_ {\gamma_ {2}} f d z + \dots + \int_ {\gamma_ {n}} f d z.\tag{6}\label{eq:4.6}
\end{equation*}
Finally, the integral over a closed curve is also invariant under a shift of parameter. The old and the new initial point determine two subarcs $\gamma_{1}, \gamma_{2}$, and the invariance follows from the fact that the integral over $\gamma_{1} + \gamma_{2}$ is equal to the integral over $\gamma_{2} + \gamma_{1}$.

In addition to integrals of the form \hyperref[eq:4.4]{(4)} we can also consider line integrals\index{Line integral} with respect to $\bar{z}$. The most convenient definition is by double conjugation
\begin{equation*}
\int_ {\gamma} f \overline {{d z}} = \overline {{\int_ {\gamma} \bar {f} d z}}.
\end{equation*}
Using this notation, line integrals with respect to $x$ or $y$ can be introduced by
\begin{equation*}
\int_ {\gamma} f d x = \frac {1}{2} \left(\int_ {\gamma} f d z + \int_ {\gamma} f \overline {{d z}}\right)
\end{equation*}
\begin{equation*}
\int_ {\gamma} f d y = \frac {1}{2 i} \Big (\int_ {\gamma} f d z - \int_ {\gamma} f \overline {{d z}} \Big).
\end{equation*}
% Original PDF physical page 120
\sourcepage{120}
With $f = u + iv$ we find that the integral\index{Integral} \hyperref[eq:4.4]{(4)} can be written in the form
\begin{equation*}
\int_ {\gamma} (u d x - v d y) + i \int_ {\gamma} (u d y + v d x)\tag{7}\label{eq:4.7}
\end{equation*}
which separates the real and imaginary part.

Of course we could just as well have started by defining integrals of the form
\begin{equation*}
\int_ {\gamma} p d x + q d y,
\end{equation*}
in which case formula \hyperref[eq:4.7]{(7)} would serve as definition of the integral \hyperref[eq:4.4]{(4)}. It is a matter of taste which one prefers.

An essentially different line integral\index{Line integral} is obtained by integration\index{Integration} with respect to arc length\index{Arc length}\index{Length}. Two notations are in common use, and the definition is
\begin{equation*}
\int_ {\gamma} f d s = \int_ {\gamma} f | d z | = \int_ {\gamma} f (z (t)) | z ^ {\prime} (t) | d t.\tag{8}\label{eq:4.8}
\end{equation*}
This integral is again independent of the choice of parameter. In contrast to \hyperref[eq:4.5]{(5)} we have now
\begin{equation*}
\int_ {- \gamma} f | d z | = \int_ {\gamma} f | d z |
\end{equation*}
while \hyperref[eq:4.6]{(6)} remains valid in the same form. The inequality
\begin{equation*}
\left| \int_ {\gamma} f d z \right| \leq \int_ {\gamma} | f | \cdot | d z |\tag{9}\label{eq:4.9}
\end{equation*}
is a consequence of \hyperref[eq:4.3]{(3)}.

For $f = 1$ the integral \hyperref[eq:4.8]{(8)} reduces to $\int_{\gamma} |dz|$ which is by definition the length of $\gamma$. As an example we compute the length of a circle. From the parametric equation $z = z(t) = a + \rho e^{it}$, $0 \leq t \leq 2\pi$, of a full circle we obtain $z'(t) = i\rho e^{it}$ and hence
\begin{equation*}
\int_ {0} ^ {2 \pi} | z ^ {\prime} (t) | d t = \int_ {0} ^ {2 \pi} \rho d t = 2 \pi \rho
\end{equation*}
as expected.

\setcounter{section}{1}\setcounter{subsection}{1}\subsection{Rectifiable Arcs}\index{Arc!rectifiable|(}\index{Rectifiable arc|(}\index{Arc!rectifiable}\index{Rectifiable arc}\label{sec:4.1.2}
The length of an arc can also be defined as the least upper bound of all sums
\begin{equation*}
\left| z \left(t _ {1}\right) - z \left(t _ {0}\right) \right| + \left| z \left(t _ {2}\right) - z \left(t _ {1}\right) \right| + \dots + \left| z \left(t _ {n}\right) - z \left(t _ {n - 1}\right) \right|\tag{10}\label{eq:4.10}
\end{equation*}
where $a = t_{0} < t_{1} < \cdots < t_{n} = b$. If this least upper bound is finite we say that the arc is rectifiable. It is quite easy to show that piecewise differentiable arcs are rectifiable, and that the two definitions of length coincide.


% Original PDF physical page 121
\index{Integral|)}\index{Line integral|)}\sourcepage{121}
Because $|x(t_k) - x(t_{k-1})| \leq |z(t_k) - z(t_{k-1})|, \quad |y(t_k) - y(t_{k-1})| \leq |z(t_k) - z(t_{k-1})|$ and $|z(t_k) - z(t_{k-1})| \leq |x(t_k) - x(t_{k-1})| + |y(t_k) - y(t_{k-1})|$ it is clear that the sums \hyperref[eq:4.10]{(10)} and the corresponding sums
\begin{equation*}
\begin{array}{l} | x (t _ {1}) - x (t _ {0}) | + \cdot \cdot \cdot + | x (t _ {n}) - x (t _ {n - 1}) | \\ | y (t _ {1}) - y (t _ {0}) | + \cdot \cdot \cdot + | y (t _ {n}) - y (t _ {n - 1}) | \end{array}
\end{equation*}
are bounded at the same time. When the latter sums are bounded, one says that the functions $x(t)$ and $y(t)$ are of bounded variation\index{Bounded variation}. An arc $z = z(t)$ is rectifiable if and only if the real and imaginary parts of $z(t)$ are of bounded variation.

If $\gamma$ is rectifiable and $f(z)$ continuous on $\gamma$ it is possible to define integrals of type \hyperref[eq:4.8]{(8)} as a limit
\begin{equation*}
\int_ {\gamma} f d s = \lim \sum_ {k = 1} ^ {n} f (z (t _ {k})) | z (t _ {k}) - z (t _ {k - 1}) |.
\end{equation*}
Here the limit is of the same kind as that encountered in the definition of a definite integral.

In the elementary theory of analytic functions it is seldom necessary to consider arcs which are rectifiable, but not piecewise differentiable. However, the notion of rectifiable arc\index{Arc!rectifiable}\index{Rectifiable arc} is one that every mathematician should know.

\setcounter{section}{1}\setcounter{subsection}{2}\subsection{Line Integrals as Functions of Arcs}\label{sec:4.1.3}
General line integrals of the form $\int_{\gamma} p \, dx + q \, dy$ are often studied as functions (or functionals) of the arc $\gamma$. It is then assumed that p and q are defined and continuous in a region $\Omega$ and that $\gamma$ is free to vary in $\Omega$. An important class of integrals is characterized by the property that the integral over an arc depends only on its end points. In other words, if $\gamma_{1}$ and $\gamma_{2}$ have the same initial point and the same end point, we require that $\int_{\gamma_{1}} p \, dx + q \, dy = \int_{\gamma_{2}} p \, dx + q \, dy$. To say that an integral depends only on the end points is equivalent to saying that the integral over any closed curve is zero. Indeed, if $\gamma$ is a closed curve, then $\gamma$ and $-\gamma$ have the same end points, and if the integral depends only on the end points, we obtain
\begin{equation*}
\int_ {\gamma} = \int_ {- \gamma} = - \int_ {\gamma}
\end{equation*}
and consequently $\int_{\gamma} = 0$. Conversely, if $\gamma_1$ and $\gamma_2$ have the same end points, then $\gamma_1 - \gamma_2$ is a closed curve, and if the integral over any closed curve vanishes, it follows that $\int_{\gamma_1} = \int_{\gamma_2}$.


% Original PDF physical page 122
\index{Arc!rectifiable|)}\index{Rectifiable arc|)}\sourcepage{122}
\bookfigure{122}{1}{FIG. 4-1}
The following theorem gives a necessary and sufficient condition under which a line integral depends only on the end points.

\begin{originalstatement}{Theorem}{1}\label{thm:4.1}
The line integral $\int_{\gamma} p \, dx + q \, dy$, defined in $\Omega$, depends only on the end points of $\gamma$ if and only if there exists a function $U(x,y)$ in $\Omega$ with the partial derivatives $\partial U / \partial x = p$, $\partial U / \partial y = q$.
\end{originalstatement}
The sufficiency follows at once, for if the condition is fulfilled we can write, with the usual notations,
\begin{equation*}
\begin{aligned}
\int_\gamma p\,dx+q\,dy
&=\int_a^b\left(\frac{\partial U}{\partial x}x'(t)+\frac{\partial U}{\partial y}y'(t)\right)dt\\
&=\int_a^b\frac{d}{dt}U(x(t),y(t))\,dt\\
&=U(x(b),y(b))-U(x(a),y(a)),
\end{aligned}
\end{equation*}
and the value of this difference depends only on the end points. To prove the necessity we choose a fixed point $(x_{0},y_{0})\in\Omega$, join it to $(x,y)$ by a polygon $\gamma$, contained in $\Omega$, whose sides are parallel to the coordinate axes (Fig.~\hyperref[fig:4-1]{4-1}) and define a function by
\begin{equation*}
U (x, y) = \int_ {\gamma} p d x + q d y.
\end{equation*}
Since the integral depends only on the end points, the function is well defined. Moreover, if we choose the last segment of $\gamma$ horizontal, we can keep y constant and let x vary without changing the other segments. On the last segment we can choose x for parameter and obtain
\begin{equation*}
U (x, y) = \int^ {x} p (x, y) d x + \text { const. },
\end{equation*}
the lower limit of the integral being irrelevant. From this expression it
% Original PDF physical page 123
\sourcepage{123}
follows at once that $\partial U/\partial x = p$. In the same way, by choosing the last segment vertical, we can show that $\partial U/\partial y = q$.

It is customary to write $dU = (\partial U/\partial x) dx + (\partial U/\partial y) dy$ and to say that an expression $p dx + q dy$ which can be written in this form is an exact differential\index{Exact differential}. Thus an integral depends only on the end points if and only if the integrand is an exact differential. Observe that p, q and U can be either real or complex. The function U, if it exists, is uniquely determined up to an additive constant, for if two functions have the same partial derivatives their difference must be constant.

When is $f(z) \, dz = f(z) \, dx + if(z) \, dy$ an exact differential? According to the definition there must exist a function $F(z)$ in $\Omega$ with the partial derivatives
\begin{equation*}
\begin{array}{l} \frac {\partial F (z)}{\partial x} = f (z) \\ \frac {\partial F (z)}{\partial y} = i f (z). \end{array}
\end{equation*}
If this is so, $F(z)$ fulfills the Cauchy-Riemann equation
\begin{equation*}
\frac {\partial F}{\partial x} = - i \frac {\partial F}{\partial y};
\end{equation*}
since $f(z)$ is by assumption continuous (otherwise $\int_{\gamma} f dz$ would not be defined) $F(z)$ is analytic with the derivative $f(z)$ (Chap.~\hyperref[chap:2]{2}, Sec.~\hyperref[sec:2.1.2]{1.2}).

The integral $\int_{\gamma} f dz$, with continuous $f$, depends only on the end points of $\gamma$ if and only if $f$ is the derivative of an analytic function in $\Omega$.

Under these circumstances we shall prove later that $f(z)$ is itself analytic.

As an immediate application of the above result we find that
\begin{equation*}
\int_ {\gamma} (z - a) ^ {n} d z = 0\tag{11}\label{eq:4.11}
\end{equation*}
for all closed curves $\gamma$, provided that the integer n is $\geq0$. In fact, $(z-a)^{n}$ is the derivative of $(z-a)^{n+1}/(n+1)$, a function which is analytic in the whole plane. If n is negative, but $\neq-1$, the same result holds for all closed curves which do not pass through a, for in the complementary region of the point a the indefinite integral is still analytic and single-valued. For n=-1, \hyperref[eq:4.11]{(11)} does not always hold. Consider a circle C with the center a, represented by the equation $z=a+\rho e^{it}$, $0\leq t\leq2\pi$. We obtain
\begin{equation*}
\int_ {C} \frac {d z}{z - a} = \int_ {0} ^ {2 \pi} i d t = 2 \pi i.
\end{equation*}
% Original PDF physical page 124
\sourcepage{124}
This result shows that it is impossible to define a single-valued branch of $\log(z-a)$ in an annulus $\rho_{1}<|z-a|<\rho_{2}$. On the other hand, if the closed curve $\gamma$ is contained in a half plane which does not contain a, the integral vanishes, for in such a half plane a single-valued and analytic branch of $\log(z-a)$ can be defined.

\subsection*{EXERCISES}
\originalnumber{1}{exercise:4.1.3.1.1}
1. Compute
\begin{equation*}
\int_ {\gamma} x d z
\end{equation*}
where $\gamma$ is the directed line segment from 0 to 1 + i.

\originalnumber{2}{exercise:4.1.3.1.2}
2. Compute
\begin{equation*}
\int_ {| z | = r} x d z,
\end{equation*}
for the positive sense of the circle, in two ways: first, by use of a parameter, and second, by observing that $x = \frac{1}{2}(z + \bar{z}) = \frac{1}{2}\left(z + \frac{r^2}{z}\right)$ on the circle.

\originalnumber{3}{exercise:4.1.3.1.3}
3. Compute
\begin{equation*}
\int_ {| z | = 2} \frac {d z}{z ^ {2} - 1}
\end{equation*}
for the positive sense of the circle.

\originalnumber{4}{exercise:4.1.3.1.4}
4. Compute
\begin{equation*}
\int_ {| z | = 1} | z - 1 | \cdot | d z |.
\end{equation*}
\originalnumber{5}{exercise:4.1.3.1.5}
5. Suppose that $f(z)$ is analytic on a closed curve $\gamma$ (i.e., f is analytic in a region that contains $\gamma$ ). Show that
\begin{equation*}
\int_ {\gamma} \overline {{{f (z)}}} f ^ {\prime} (z) d z
\end{equation*}
is purely imaginary. (The continuity of $f'(z)$ is taken for granted.)

\originalnumber{6}{exercise:4.1.3.1.6}
6. Assume that $f(z)$ is analytic and satisfies the inequality $|f(z) - 1| < 1$ in a region $\Omega$. Show that
\begin{equation*}
\int_ {\gamma} \frac {f ^ {\prime} (z)}{f (z)} d z = 0
\end{equation*}
for every closed curve in $\Omega$. (The continuity of $f'(z)$ is taken for granted.)

\originalnumber{7}{exercise:4.1.3.1.7}
7. If $P(z)$ is a polynomial and $C$ denotes the circle $|z - a| = R$, what is the value of $\int_{C} P(z) d\bar{z}$? Answer: $-2\pi iR^2P'(a)$.


% Original PDF physical page 125
\sourcepage{125}
\originalnumber{8}{exercise:4.1.3.1.8}
8.\index{Cauchy's integral theorem|(} Describe a set of circumstances under which the formula
\begin{equation*}
\int_ {\gamma} \log z d z = 0
\end{equation*}
is meaningful and true.

\setcounter{section}{1}\setcounter{subsection}{3}\subsection{Cauchy's Theorem for a Rectangle}\label{sec:4.1.4}
There are several forms of Cauchy's theorem, but they differ in their topological rather than in their analytical content. It is natural to begin with a case in which the topological considerations are trivial.

We consider, specifically, a rectangle R defined by inequalities $a \leq x \leq b$, $c \leq y \leq d$. Its perimeter can be considered as a simple closed curve consisting of four line segments whose direction we choose so that R lies to the left of the directed segments. The order of the vertices is thus $(a,c)$, $(b,c)$, $(b,d)$, $(a,d)$. We refer to this closed curve as the boundary curve or contour\index{Contour} of R, and we denote it by $\partial R.$\hyperlink{fn:125:dagger}{\textsuperscript{\ensuremath{\dagger}}}

We emphasize that R is chosen as a closed point set and, hence, is not a region. In the theorem that follows we consider a function which is analytic on the rectangle R. We recall to the reader that such a function is by definition defined and analytic in an open set which contains R.

The following is a preliminary version of Cauchy's theorem:

\begin{originalstatement}{Theorem}{2}\label{thm:4.2}
If the function $f(z)$ is analytic on $R$, then
\begin{equation*}
\int_ {\partial R} f (z) d z = 0.\tag{12}\label{eq:4.12}
\end{equation*}
\end{originalstatement}
The proof is based on the method of bisection. Let us introduce the notation
\begin{equation*}
\eta (R) = \int_ {\partial R} f (z) d z
\end{equation*}
which we will also use for any rectangle contained in the given one. If $R$ is divided into four congruent rectangles $R^{(1)}, R^{(2)}, R^{(3)}, R^{(4)}$, we find that
\begin{equation*}
\eta (R) = \eta (R ^ {(1)}) + \eta (R ^ {(2)}) + \eta (R ^ {(3)}) + \eta (R ^ {(4)}),\tag{13}\label{eq:4.13}
\end{equation*}
for the integrals over the common sides cancel each other. It is important to note that this fact can be verified explicitly and does not make illicit use of geometric intuition. Nevertheless, a reference to Fig.~\hyperref[fig:4-2]{4-2} is helpful.

\begingroup\renewcommand{\thefootnote}{}\footnotetext{\hypertarget{fn:125:dagger}{}\textsuperscript{\ensuremath{\dagger}} This is standard notation, and we shall use it repeatedly. Note that by earlier convention $\partial R$ is also the boundary of $R$ as a point set (Chap.~\hyperref[chap:3]{3}, Sec.~\hyperref[sec:3.1.2]{1.2}).}\endgroup

% Original PDF physical page 126
\sourcepage{126}
\bookfigure{126}{1}{FIG. 4-2. Bisection of rectangle.}
It follows from \hyperref[eq:4.13]{(13)} that at least one of the rectangles $R^{(k)}$, $k = 1, 2, 3, 4$, must satisfy the condition
\begin{equation*}
\left| \eta (R ^ {(k)}) \right| \geq \frac {1}{4} | \eta (R) |.
\end{equation*}
We denote this rectangle by $R_1$; if several $R^{(k)}$ have this property, the choice shall be made according to some definite rule.

This process can be repeated indefinitely, and we obtain a sequence of nested rectangles $R \supset R_{1} \supset R_{2} \supset \cdots \supset R_{n} \supset \cdots$ with the property
\begin{equation*}
\left| \eta (R _ {n}) \right| \geq \frac {1}{4} \left| \eta (R _ {n - 1}) \right|
\end{equation*}
and hence
\begin{equation*}
| \eta (R _ {n}) | \geq 4 ^ {- n} | \eta (R) |.\tag{14}\label{eq:4.14}
\end{equation*}
The rectangles $R_{n}$ converge to a point $z^{*} \in R$ in the sense that $R_{n}$ will be contained in a prescribed neighborhood $|z - z^{*}| < \delta$ as soon as n is sufficiently large. First of all, we choose $\delta$ so small that $f(z)$ is defined and analytic in $|z - z^{*}| < \delta$. Secondly, if $\varepsilon > 0$ is given, we can choose $\delta$ so that
\begin{equation*}
\left| \frac {f (z) - f (z ^ {*})}{z - z ^ {*}} - f ^ {\prime} (z ^ {*}) \right| <   \varepsilon
\end{equation*}
or
\begin{equation*}
\left| f (z) - f \left(z ^ {*}\right) - \left(z - z ^ {*}\right) f ^ {\prime} \left(z ^ {*}\right) \right| <   \varepsilon \left| z - z ^ {*} \right|\tag{15}\label{eq:4.15}
\end{equation*}
for $|z - z^{*}| < \delta$. We assume that $\delta$ satisfies both conditions and that $R_{n}$ is contained in $|z - z^{*}| < \delta$.

We make now the observation that
\begin{equation*}
\int_ {\partial R _ {n}} d z = 0
\end{equation*}
\begin{equation*}
\int_ {\partial R _ {n}} z d z = 0.
\end{equation*}
% Original PDF physical page 127
\sourcepage{127}
These trivial special cases of our theorem have already been proved in Sec.~\hyperref[sec:4.1.1]{1.1}. We recall that the proof depended on the fact that 1 and z are the derivatives of z and $z^{2}/2$, respectively.

By virtue of these equations we are able to write
\begin{equation*}
\eta (R _ {n}) = \int_ {\partial R _ {n}} [ f (z) - f (z ^ {*}) - (z - z ^ {*}) f ^ {\prime} (z ^ {*}) ] d z,
\end{equation*}
and it follows by \hyperref[eq:4.15]{(15)} that
\begin{equation*}
\left| \eta (R _ {n}) \right| \leq \varepsilon \int_ {\partial R _ {n}} | z - z ^ {*} | \cdot | d z |.\tag{16}\label{eq:4.16}
\end{equation*}
In the last integral $|z - z^{*}|$ is at most equal to the length $d_{n}$ of the diagonal of $R_{n}$. If $L_{n}$ denotes the length of the perimeter of $R_{n}$, the integral is hence $\leq d_{n}L_{n}$. But if d and L are the corresponding quantities for the original rectangle R, it is clear that $d_{n} = 2^{-n}d$ and $L_{n} = 2^{-n}L$. By \hyperref[eq:4.16]{(16)} we have hence
\begin{equation*}
\left| \eta (R _ {n}) \right| \leq 4 ^ {- n} d L \varepsilon ,
\end{equation*}
and comparison with \hyperref[eq:4.14]{(14)} yields
\begin{equation*}
| \eta (R) | \leq d L \varepsilon .
\end{equation*}
Since $\varepsilon$ is arbitrary, we can only have $\eta(R) = 0$, and the theorem is proved.

This beautiful proof, which could hardly be simpler, is due to É. Goursat\index{Goursat E@Goursat, E.} who discovered that the classical hypothesis of a continuous $f'(z)$ is redundant. At the same time the proof is simpler than the earlier proofs inasmuch as it leans neither on double integration\index{Integration} nor on differentiation under the integral sign.

The hypothesis in Theorem~\hyperref[thm:4.2]{2} can be weakened considerably. We shall prove at once the following stronger theorem which will find very important use.

\begin{originalstatement}{Theorem}{3}\label{thm:4.3}
Let $f(z)$ be analytic on the set $R'$ obtained from a rectangle $R$ by omitting a finite number of interior points $\zeta_j$. If it is true that
\begin{equation*}
\lim _ {z \rightarrow \zeta_ {j}} (z - \zeta_ {j}) f (z) = 0
\end{equation*}
for all $j$, then
\begin{equation*}
\int_ {\partial R} f (z) d z = 0.
\end{equation*}
\end{originalstatement}
It is sufficient to consider the case of a single exceptional point $\zeta$, for evidently $R$ can be divided into smaller rectangles which contain at most one $\zeta$.

We divide R into nine rectangles, as shown in Fig.~\hyperref[fig:4-3]{4-3}, and apply


% Original PDF physical page 128
\sourcepage{128}
\bookfigure{128}{1}{FIG. 4-3}
Theorem~\hyperref[thm:4.2]{2} to all but the rectangle $R_{0}$ in the center. If the corresponding equations \hyperref[eq:4.12]{(12)} are added, we obtain, after cancellations,
\begin{equation*}
\int_ {\partial R} f d z = \int_ {\partial R _ {0}} f d z.\tag{17}\label{eq:4.17}
\end{equation*}
If $\varepsilon > 0$ we can choose the rectangle $R_0$ so small that
\begin{equation*}
| f (z) | \leq \frac {\varepsilon}{| z - \zeta |}
\end{equation*}
on $\partial R_0$. By \hyperref[eq:4.17]{(17)} we have thus
\begin{equation*}
\left| \int_ {\partial R} f d z \right| \leq \varepsilon \int_ {\partial R _ {0}} \frac {| d z |}{| z - \zeta |}.
\end{equation*}
If we assume, as we may, that $R_{0}$ is a square of center $\zeta$, elementary estimates show that
\begin{equation*}
\int_ {\partial R _ {0}} \frac {| d z |}{| z - \zeta |} <   8.
\end{equation*}
Thus we obtain
\begin{equation*}
\left| \int_ {\partial R} f d z \right| <   8 \varepsilon ,
\end{equation*}
and since $\varepsilon$ is arbitrary the theorem follows.

We observe that the hypothesis of the theorem is certainly fulfilled if $f(z)$ is analytic and bounded on $R'$.

\setcounter{section}{1}\setcounter{subsection}{4}\subsection{Cauchy's Theorem in a Disk}\label{sec:4.1.5}
It is not true that the integral of an analytic function over a closed curve is always zero.


% Original PDF physical page 129
\sourcepage{129}
Indeed,\index{Function!integral|(} we have found that
\begin{equation*}
\int_ {C} \frac {d z}{z - a} = 2 \pi i
\end{equation*}
when C is a circle about a. In order to make sure that the integral vanishes, it is necessary to make a special assumption concerning the region $\Omega$ in which $f(z)$ is known to be analytic and to which the curve $\gamma$ is restricted. We are not yet in a position to formulate this condition, and for this reason we must restrict attention to a very special case. In what follows we assume that $\Omega$ is an open disk $|z - z_{0}| < \rho$ to be denoted by $\Delta$.

\begin{originalstatement}{Theorem}{4}\label{thm:4.4}
If $f(z)$ is analytic in an open disk $\Delta$, then
\begin{equation*}
\int_ {\gamma} f (z) d z = 0\tag{18}\label{eq:4.18}
\end{equation*}
for every closed curve $\gamma$ in $\Delta$.

\end{originalstatement}
The proof is a repetition of the argument used in proving the second half of Theorem~\hyperref[thm:4.1]{1}. We define a function $F(z)$ by
\begin{equation*}
F (z) = \int_ {\sigma} f d z\tag{19}\label{eq:4.19}
\end{equation*}
where $\sigma$ consists of the horizontal line segment from the center $(x_{0},y_{0})$ to $(x,y_{0})$ and the vertical segment from $(x,y_{0})$ to $(x,y)$; it is immediately seen that $\partial F/\partial y = if(z)$. On the other hand, by Theorem~\hyperref[thm:4.2]{2} $\sigma$ can be replaced by a path consisting of a vertical segment followed by a horizontal segment. This choice defines the same function $F(z)$, and we obtain $\partial F/\partial x = f(z)$. Hence $F(z)$ is analytic in $\Delta$ with the derivative $f(z)$, and $f(z) dz$ is an exact differential.

Clearly, the same proof would go through for any region which contains the rectangle with the opposite vertices $z_{0}$ and z as soon as it contains z. A rectangle, a half plane, or the inside of an ellipse all have this property, and hence Theorem~\hyperref[thm:4.4]{4} holds for any of these regions. By this method we cannot, however, reach full generality.

For the applications it is very important that the conclusion of Theorem~\hyperref[thm:4.4]{4} remains valid under the weaker condition of Theorem~\hyperref[thm:4.3]{3}. We state this as a separate theorem.

\begin{originalstatement}{Theorem}{5}\label{thm:4.5}
Let $f(z)$ be analytic in the region $\Delta'$ obtained by omitting a finite number of points $\zeta_j$ from an open disk $\Delta$. If $f(z)$ satisfies the condition $\lim_{z \to \zeta_j} (z - \zeta_j)f(z) = 0$ for all $j$, then \hyperref[eq:4.18]{(18)} holds for any closed curve $\gamma$ in $\Delta'$.
\end{originalstatement}

% Original PDF physical page 130
\index{Function!integral|)}\sourcepage{130}
\bookfigure{130}{1}{FIG. 4-4}\index{Winding number|(}
The proof must be modified, for we cannot let $\sigma$ pass through the exceptional points. Assume first that no $\zeta_{j}$ lies on the lines $x = x_{0}$ and $y = y_{0}$. It is then possible to avoid the exceptional points by letting $\sigma$ consist of three segments (Fig.~\hyperref[fig:4-4]{4-4}). By an obvious application of Theorem~\hyperref[thm:4.3]{3} we find that the value of $F(z)$ in \hyperref[eq:4.18]{(18)} is independent of the choice of the middle segment; moreover, the last segment can be either vertical or horizontal. We conclude as before that $F(z)$ is an indefinite integral of $f(z)$, and the theorem follows.

In case there are exceptional points on the lines $x = x_{0}$ and $y = y_{0}$ the reader will easily convince himself that a similar proof can be carried out, provided that we use four line segments in the place of three.

\setcounter{section}{1}\section{Cauchy's Integral Formula}\label{sec:4.2}
Through a very simple application of Cauchy's theorem it becomes possible to represent an analytic function $f(z)$ as a line integral in which the variable z enters as a parameter. This representation, known as Cauchy's integral formula\index{Cauchy's integral formula|(}\index{Cauchy's integral formula}, has numerous important applications. Above all, it enables us to study the local properties of an analytic function in great detail.

\setcounter{section}{2}\setcounter{subsection}{0}\subsection{The Index of a Point with Respect to a Closed Curve}\index{Index of a point|(}\index{Index of a point}\label{sec:4.2.1}
As a preliminary to the derivation of Cauchy's formula we must define a notion which in a precise way indicates how many times a closed curve winds around a fixed point not on the curve. If the curve is piecewise differentiable, as we shall assume without serious loss of generality, the definition can be based on the following lemma:


% Original PDF physical page 131
\sourcepage{131}
\begin{originalstatement}{Lemma}{1}\label{lem:4.1}
If the piecewise differentiable closed curve $\gamma$ does not pass through the point a, then the value of the integral
\begin{equation*}
\int_ {\gamma} \frac {d z}{z - a}
\end{equation*}
is a multiple of $2\pi i$.

\end{originalstatement}
This lemma may seem trivial, for we can write
\begin{equation*}
\int_ {\gamma} \frac {d z}{z - a} = \int_ {\gamma} d \log (z - a) = \int_ {\gamma} d \log | z - a | + i \int_ {\gamma} d \arg (z - a).
\end{equation*}
When z describes a closed curve, $\log |z - a|$ returns to its initial value and $\arg (z - a)$ increases or decreases by a multiple of $2\pi$. This would seem to imply the lemma, but more careful thought shows that the reasoning is of no value unless we define $\arg (z - a)$ in a unique way.

The simplest proof is computational. If the equation of $\gamma$ is $z = z(t)$, $\alpha \leq t \leq \beta$, let us consider the function
\begin{equation*}
h (t) = \int_ {\alpha} ^ {t} \frac {z ^ {\prime} (t)}{z (t) - a} d t.
\end{equation*}
It is defined and continuous on the closed interval $[\alpha,\beta]$, and it has the derivative
\begin{equation*}
h ^ {\prime} (t) = \frac {z ^ {\prime} (t)}{z (t) - a}
\end{equation*}
whenever $z'(t)$ is continuous. From this equation it follows that the derivative of $e^{-h(t)}(z(t) - a)$ vanishes except perhaps at a finite number of points, and since this function is continuous it must reduce to a constant. We have thus
\begin{equation*}
e ^ {h (t)} = \frac {z (t) - a}{z (\alpha) - a}.
\end{equation*}
Since $z(\beta) = z(\alpha)$ we obtain $e^{h(\beta)} = 1$, and therefore $h(\beta)$ must be a multiple of $2\pi i$. This proves the lemma.

We can now define the index of the point $a$ with respect to the curve $\pmb{\gamma}$ by the equation
\begin{equation*}
n (\gamma , a) = \frac {1}{2 \pi i} \int_ {\gamma} \frac {d z}{z - a}.
\end{equation*}
With a suggestive terminology the index is also called the winding number\index{Winding number} of $\gamma$ with respect to $a$.

It is clear that $n(-\gamma, a) = -n(\gamma, a)$.

The following property is an immediate consequence of Theorem~\hyperref[thm:4.4]{4}:


% Original PDF physical page 132
\sourcepage{132}
(i)\index{Region!determined by symbol@determined by $\gamma$|(} If $\gamma$ lies inside of a circle, then $n(\gamma, a) = 0$ for all points a outside of the same circle.

As a point set $\gamma$ is closed and bounded. Its complement is open and can be represented as a union of disjoint regions, the components of the complement. We shall say, for short, that $\gamma$ determines these regions. If the complementary regions are considered in the extended plane, there is exactly one which contains the point at infinity. Consequently, $\gamma$ determines one and only one unbounded region.

(ii) As a function of a the index $n(\gamma, a)$ is constant in each of the regions determined by $\gamma$, and zero in the unbounded region.

Any two points in the same region determined by $\gamma$ can be joined by a polygon which does not meet $\gamma$. For this reason it is sufficient to prove that $n(\gamma, a) = n(\gamma, b)$ if $\gamma$ does not meet the line segment from $a$ to $b$. Outside of this segment the function $(z - a)/(z - b)$ is never real and $\leq 0$. For this reason the principal branch of $\log [(z - a)/(z - b)]$ is analytic in the complement of the segment. Its derivative is equal to $(z - a)^{-1} - (z - b)^{-1}$, and if $\gamma$ does not meet the segment we must have
\begin{equation*}
\int_ {\gamma} \left(\frac {1}{z - a} - \frac {1}{z - b}\right) d z = 0;
\end{equation*}
hence $n(\gamma, a) = n(\gamma, b)$. If $|a|$ is sufficiently large, $\gamma$ is contained in a disk $|z| < \rho < |a|$ and we conclude by (i) that $n(\gamma, a) = 0$. This proves that $n(\gamma, a) = 0$ in the unbounded region.

We shall find the case $n(\gamma,a)=1$ particularly important, and it is desirable to formulate a geometric condition which leads to this consequence. For simplicity we take a=0.

\begin{originalstatement}{Lemma}{2}\label{lem:4.2}
Let $z_1, z_2$ be two points on a closed curve $\gamma$ which does not pass through the origin. Denote the subarc from $z_1$ to $z_2$ in the direction of the curve by $\gamma_1$, and the subarc from $z_2$ to $z_1$ by $\gamma_2$. Suppose that $z_1$ lies in the lower half plane and $z_2$ in the upper half plane. If $\gamma_1$ does not meet the negative real axis and $\gamma_2$ does not meet the positive real axis, then $n(\gamma,0) = 1$.
\end{originalstatement}
For the proof we draw the half lines $L_{1}$ and $L_{2}$ from the origin through $z_{1}$ and $z_{2}$ (Fig.~\hyperref[fig:4-5]{4-5}). Let $\zeta_{1}, \zeta_{2}$ be the points in which $L_{1}, L_{2}$ intersect a circle C about the origin. If C is described in the positive sense, the arc $C_{1}$ from $\zeta_{1}$ to $\zeta_{2}$ does not intersect the negative axis, and the arc $C_{2}$ from $\zeta_{2}$ to $\zeta_{1}$ does not intersect the positive axis. Denote the directed line segments from $z_{1}$ to $\zeta_{1}$ and from $z_{2}$ to $\zeta_{2}$ by $\delta_{1}, \delta_{2}$. Introducing the closed curves $\sigma_{1} = \gamma_{1} + \delta_{2} - C_{1} - \delta_{1}, \sigma_{2} = \gamma_{2} + \delta_{1} - C_{2} - \delta_{2}$ we find that $n(\gamma,0) = n(C,0) + n(\sigma_{1},0) + n(\sigma_{2},0)$ because of cancellations. But $\sigma_{1}$ does not meet the negative axis. Hence the origin belongs to the


% Original PDF physical page 133
\index{Region!determined by symbol@determined by $\gamma$|)}\sourcepage{133}
\bookfigure{133}{1}{FIG. 4-5}
unbounded region determined by $\sigma_{1}$, and we obtain $n(\sigma_{1},0) = 0$. For a similar reason $n(\sigma_{2},0) = 0$, and we conclude that $n(\gamma,0) = n(C,0) = 1$.
\subsection*{*EXERCISES}
These are not routine exercises. They serve to illustrate the topological use of winding numbers\index{Winding number}.

\originalnumber{1}{exercise:4.2.1.2.1}
1. Give an alternate proof of Lemma~\hyperref[lem:4.1]{1} by dividing $\gamma$ into a finite number of subarcs such that there exists a single-valued branch of arg $(z - a)$ on each subarc. Pay particular attention to the compactness argument that is needed to prove the existence of such a subdivision.

\originalnumber{2}{exercise:4.2.1.2.2}
2. It is possible to define $n(\gamma,a)$ for any continuous closed curve $\gamma$ that does not pass through a, whether piecewise differentiable or not. For this purpose $\gamma$ is divided into subarcs $\gamma_{1},\ldots,\gamma_{n}$, each contained in a disk that does not include a. Let $\sigma_{k}$ be the directed line segment from the initial to the terminal point of $\gamma_{k}$, and set $\sigma=\sigma_{1}+\cdots+\sigma_{n}$. We define $n(\gamma,a)$ to be the value of $n(\sigma,a)$.

To justify the definition, prove the following:

(a) the result is independent of the subdivision;

(b) if $\gamma$ is piecewise differentiable the new definition is equivalent to the old;

(c) the properties (i) and (ii) of the text continue to hold.


% Original PDF physical page 134
\sourcepage{134}
\bookfigure{134}{1}{FIG. 4-6. Part of the Jordan curve theorem.}
\originalnumber{3}{exercise:4.2.1.2.3}
3. The Jordan curve theorem\index{Jordan curve theorem} asserts that every Jordan curve in the plane determines exactly two regions. The notion of winding number\index{Winding number} leads to a quick proof of one part of the theorem, namely that the complement of a Jordan curve $\gamma$ has at least two components. This will be so if there exists a point a with $n(\gamma,a) \neq 0$.

We may assume that $\operatorname{Re} z > 0$ on $\gamma$, and that there are points $z_1, z_2 \in \gamma$ with $\operatorname{Im} z_1 < 0$, $\operatorname{Im} z_2 > 0$. These points may be chosen so that there are no other points of $\gamma$ on the line segments from 0 to $z_1$ and from 0 to $z_2$. Let $\gamma_1$ and $\gamma_2$ be the arcs of $\gamma$ from $z_1$ to $z_2$ (excluding the end points).

Let $\sigma_{1}$ be the closed curve that consists of the line segment from 0 to $z_{1}$ followed by $\gamma_{1}$ and the segment from $z_{2}$ to 0, and let $\sigma_{2}$ be constructed in the same way with $\gamma_{2}$ in the place of $\gamma_{1}$. Then $\sigma_{1} - \sigma_{2} = \gamma$ or $-\gamma$.

The positive real axis intersects both $\gamma_{1}$ and $\gamma_{2}$ (why?). Choose the notation so that the intersection $x_{2}$ farthest to the right is with $\gamma_{2}$ (Fig.~\hyperref[fig:4-6]{4-6}).

Prove the following:

(a) $n(\sigma_1,x_2) = 0$, hence $n(\sigma_1,z) = 0$ for $z\in \gamma_{2}$

(b) $n(\sigma_1,x) = n(\sigma_2,x) = 1$ for small $x > 0$ (Lemma~\hyperref[lem:4.2]{2});

(c) the first intersection $x_{1}$ of the positive real axis with $\gamma$ lies on $\gamma_{1}$;

(d) $n(\sigma_2, x_1) = 1$, hence $n(\sigma_2, z) = 1$ for $z \in \gamma_1$;

(e) there exists a segment of the positive real axis with one end point on $\gamma_{1}$, the other on $\gamma_{2}$, and no other points on $\gamma$. The points x between the end points satisfy $n(\gamma,x)=1$ or -1.

\setcounter{section}{2}\setcounter{subsection}{1}\subsection{The Integral Formula}\label{sec:4.2.2}
Let $f(z)$ be analytic in an open disk $\Delta$. Consider a closed curve $\gamma$ in $\Delta$ and a point $a \in \Delta$ which does not lie on $\gamma$. We apply Cauchy's theorem to the function
\begin{equation*}
F (z) = \frac {f (z) - f (a)}{z - a}.
\end{equation*}
% Original PDF physical page 135
\index{Index of a point|)}\index{Winding number|)}\sourcepage{135}
This function is analytic for $z \neq a$. For $z = a$ it is not defined, but it satisfies the condition
\begin{equation*}
\lim _ {z \rightarrow a} F (z) (z - a) = \lim _ {z \rightarrow a} (f (z) - f (a)) = 0
\end{equation*}
which is the condition of Theorem~\hyperref[thm:4.5]{5}. We conclude that
\begin{equation*}
\int_ {\gamma} \frac {f (z) - f (a)}{z - a} d z = 0.
\end{equation*}
This equation can be written in the form
\begin{equation*}
\int_ {\gamma} \frac {f (z) d z}{z - a} = f (a) \int_ {\gamma} \frac {d z}{z - a},
\end{equation*}
and we observe that the integral in the right-hand member is by definition $2\pi i \cdot n(\gamma, a)$. We have thus proved:

\begin{originalstatement}{Theorem}{6}\label{thm:4.6}
Suppose that $f(z)$ is analytic in an open disk $\Delta$, and let $\gamma$ be a closed curve in $\Delta$. For any point a not on $\gamma$
\begin{equation*}
n (\gamma , a) \cdot f (a) = \frac {1}{2 \pi i} \int_ {\gamma} \frac {f (z) d z}{z - a},\tag{20}\label{eq:4.20}
\end{equation*}
where $n(\gamma, a)$ is the index of $a$ with respect to $\gamma$.

\end{originalstatement}
In this statement we have suppressed the requirement that a be a point in $\Delta$. We have done so in view of the obvious interpretation of the formula \hyperref[eq:4.20]{(20)} for the case that a is not in $\Delta$. Indeed, in this case $n(\gamma,a)$ and the integral in the right-hand member are both zero.

It is clear that Theorem~\hyperref[thm:4.6]{6} remains valid for any region $\Omega$ to which Theorem~\hyperref[thm:4.5]{5} can be applied. The presence of exceptional points $\zeta_j$ is permitted, provided none of them coincides with a.

The most common application is to the case where $n(\gamma,a)=1$. We have then
\begin{equation*}
f (a) = \frac {1}{2 \pi i} \int_ {\gamma} \frac {f (z) d z}{z - a},\tag{21}\label{eq:4.21}
\end{equation*}
and this we interpret as a representation formula. Indeed, it permits us to compute $f(a)$ as soon as the values of $f(z)$ on $\gamma$ are given, together with the fact that $f(z)$ is analytic in $\Delta$. In \hyperref[eq:4.21]{(21)} we may let a take different values, provided that the order of a with respect to $\gamma$ remains equal to 1. We may thus treat a as a variable, and it is convenient to change the notation and rewrite \hyperref[eq:4.21]{(21)} in the form
\begin{equation*}
f (z) = \frac {1}{2 \pi i} \int_ {\gamma} \frac {f (\zeta) d \zeta}{\zeta - z}.\tag{22}\label{eq:4.22}
\end{equation*}
% Original PDF physical page 136
\sourcepage{136}
It is this formula which is usually referred to as Cauchy's integral formula\index{Cauchy's integral formula}. We must remember that it is valid only when $n(\gamma,z)=1$, and that we have proved it only when $f(z)$ is analytic in a disk.

\subsection*{EXERCISES}
\originalnumber{1}{exercise:4.2.2.3.1}
1. Compute
\begin{equation*}
\int_ {| z | = 1} \frac {e ^ {z}}{z} d z.
\end{equation*}
\originalnumber{2}{exercise:4.2.2.3.2}
2. Compute
\begin{equation*}
\int_ {| z | = 2} \frac {d z}{z ^ {2} + 1}
\end{equation*}
by decomposition of the integrand in partial fractions.

\originalnumber{3}{exercise:4.2.2.3.3}
3. Compute
\begin{equation*}
\int_ {| z | = \rho} \frac {| d z |}{| z - a | ^ {2}}
\end{equation*}
under the condition $|a| \neq \rho$. Hint: make use of the equations $z\bar{z} = \rho^2$ and
\begin{equation*}
| d z | = - i \rho \frac {d z}{z}.
\end{equation*}
\setcounter{section}{2}\setcounter{subsection}{2}\subsection{Higher Derivatives}\label{sec:4.2.3}
The representation formula \hyperref[eq:4.22]{(22)} gives us an ideal tool for the study of the local properties of analytic functions. In particular we can now show that an analytic function has derivatives of all orders, which are then also analytic.

We consider a function $f(z)$ which is analytic in an arbitrary region $\Omega$. To a point $a \in \Omega$ we determine a $\delta$ -neighborhood $\Delta$ contained in $\Omega$, and in $\Delta$ a circle C about a. Theorem~\hyperref[thm:4.6]{6} can be applied to $f(z)$ in $\Delta$. Since $n(C, a) = 1$ we have $n(C, z) = 1$ for all points z inside of C. For such z we obtain by \hyperref[eq:4.22]{(22)}
\begin{equation*}
f (z) = \frac {1}{2 \pi i} \int_ {C} \frac {f (\zeta) d \zeta}{\zeta - z}.
\end{equation*}
Provided that the integral can be differentiated under the sign of integration\index{Integration} we find
\begin{equation*}
f ^ {\prime} (z) = \frac {1}{2 \pi i} \int_ {C} \frac {f (\zeta) d \zeta}{(\zeta - z) ^ {2}}\tag{23}\label{eq:4.23}
\end{equation*}
and
\begin{equation*}
f ^ {(n)} (z) = \frac {n !}{2 \pi i} \int_ {C} \frac {f (\zeta) d \zeta}{(\zeta - z) ^ {n + 1}}.\tag{24}\label{eq:4.24}
\end{equation*}
% Original PDF physical page 137
\sourcepage{137}
If the differentiations can be justified, we shall have proved the existence of all derivatives at the points inside of C. Since every point in $\Omega$ lies inside of some such circle, the existence will be proved in the whole region $\Omega$. At the same time we shall have obtained a convenient representation formula for the derivatives.

For the justification we could either refer to corresponding theorems in the real case, or we could prove a general theorem concerning line integrals whose integrand depends analytically on a parameter. Actually, we shall prove only the following lemma which is all we need in the present case:

\begin{originalstatement}{Lemma}{3}\label{lem:4.3}
Suppose that $\varphi (\zeta)$ is continuous on the arc $\gamma$. Then the function
\begin{equation*}
F _ {n} (z) = \int_ {\gamma} \frac {\varphi (\zeta) d \zeta}{(\zeta - z) ^ {n}}
\end{equation*}
is analytic in each of the regions determined by $\gamma$, and its derivative is $F_n'(z) = nF_{n+1}(z)$.

\end{originalstatement}
We prove first that $F_{1}(z)$ is continuous. Let $z_{0}$ be a point not on $\gamma$, and choose the neighborhood $|z - z_{0}| < \delta$ so that it does not meet $\gamma$. By restricting z to the smaller neighborhood $|z - z_{0}| < \delta/2$ we attain that $|\zeta - z| > \delta/2$ for all $\zeta \in \gamma$. From
\begin{equation*}
F _ {1} (z) - F _ {1} \left(z _ {0}\right) = (z - z _ {0}) \int_ {\gamma} \frac {\varphi (\zeta) d \zeta}{(\zeta - z) (\zeta - z _ {0})}
\end{equation*}
we obtain at once
\begin{equation*}
\left| F _ {1} (z) - F _ {1} \left(z _ {0}\right) \right| <   \left| z - z _ {0} \right| \cdot \frac {2}{\delta^ {2}} \int_ {\gamma} | \varphi | | d \zeta |,
\end{equation*}
and this inequality proves the continuity of $F_{1}(z)$ at $z_{0}$.

From this part of the lemma, applied to the function $\varphi(\zeta)/(\zeta - z_0)$, we conclude that the difference quotient
\begin{equation*}
\frac {F _ {1} (z) - F _ {1} \left(z _ {0}\right)}{z - z _ {0}} = \int_ {\gamma} \frac {\varphi (\zeta) d \zeta}{(\zeta - z) (\zeta - z _ {0})}
\end{equation*}
tends to the limit $F_{2}(z_{0})$ as $z \to z_{0}$. Hence it is proved that $F_{1}'(z) = F_{2}(z)$.

The general case is proved by induction. Suppose we have shown that $F_{n-1}'(z) = (n - 1)F_n(z)$. From the identity
\begin{equation*}
\begin{aligned}
&F_n(z)-F_n(z_0)\\
&\quad=\left[\int_\gamma\frac{\varphi\,d\zeta}{(\zeta-z)^{n-1}(\zeta-z_0)}
-\int_\gamma\frac{\varphi\,d\zeta}{(\zeta-z_0)^n}\right]\\
&\qquad+(z-z_0)\int_\gamma\frac{\varphi\,d\zeta}{(\zeta-z)^n(\zeta-z_0)}
\end{aligned}
\end{equation*}
% Original PDF physical page 138
\sourcepage{138}
we can conclude that $F_{n}(z)$ is continuous. Indeed, by the induction hypothesis, applied to $\varphi(\zeta)/(\zeta - z_{0})$, the first term tends to zero for $z \to z_{0}$, and in the second term the factor of $z - z_{0}$ is bounded in a neighborhood of $z_{0}$. Now, if we divide the identity by $z - z_{0}$ and let $z$ tend to $z_{0}$, the quotient in the first term tends to a derivative which by the induction hypothesis equals $(n - 1)F_{n+1}(z_{0})$. The remaining factor in the second term is continuous, by what we have already proved, and has the limit $F_{n+1}(z_{0})$. Hence $F_{n}'(z_{0})$ exists and equals $nF_{n+1}(z_{0})$.

It is clear that Lemma~\hyperref[lem:4.3]{3} is just what is needed in order to deduce \hyperref[eq:4.23]{(23)} and \hyperref[eq:4.24]{(24)} in a rigorous way. We have thus proved that an analytic function has derivatives of all orders which are analytic and can be represented by the formula \hyperref[eq:4.24]{(24)}.

Among the consequences of this result we like to single out two classical theorems. The first is known as Morera's theorem\index{Morera's theorem}, and it can be stated as follows:

If $f(z)$ is defined and continuous in a region $\Omega$, and if $\int_{\gamma} f dz = 0$ for all closed curves $\gamma$ in $\Omega$, then $f(z)$ is analytic in $\Omega$.

The hypothesis implies, as we have already remarked in Sec.~\hyperref[sec:4.1.3]{1.3}, that $f(z)$ is the derivative of an analytic function $F(z)$. We know now that $f(z)$ is then itself analytic.

A second classical result goes under the name of Liouville's theorem\index{Liouville's theorem}: A function which is analytic and bounded in the whole plane must reduce to a constant.

For the proof we make use of a simple estimate derived from \hyperref[eq:4.24]{(24)}. Let the radius of C be r, and assume that $|f(\zeta)| \leq M$ on C. If we apply \hyperref[eq:4.24]{(24)} with z = a, we obtain at once
\begin{equation*}
\left| f ^ {(n)} (a) \right| \leq M n! r ^ {- n}.\tag{25}\label{eq:4.25}
\end{equation*}
For Liouville's theorem we need only the case $n = 1$. The hypothesis means that $|f(\zeta)| \leq M$ on all circles. Hence we can let $r$ tend to $\infty$, and \hyperref[eq:4.25]{(25)} leads to $f'(a) = 0$ for all $a$. We conclude that the function is constant.

Liouville's theorem leads to an almost trivial proof of the fundamental theorem of algebra\index{Fundamental theorem of algebra}. Suppose that $P(z)$ is a polynomial of degree $>0$. If $P(z)$ were never zero, the function $1 / P(z)$ would be analytic in the whole plane. We know that $P(z) \to \infty$ for $z \to \infty$, and therefore $1 / P(z)$ tends to zero. This implies boundedness (the absolute value is continuous on the Riemann sphere and has thus a finite maximum), and by Liouville's theorem $1 / P(z)$ would be constant. Since this is not so, the equation $P(z) = 0$ must have a root.

The inequality \hyperref[eq:4.25]{(25)} is known as Cauchy's estimate\index{Cauchy's estimate}. It shows above
% Original PDF physical page 139
\sourcepage{139}
all that the successive derivatives of an analytic function cannot be arbitrary; there must always exist an M and an r so that \hyperref[eq:4.25]{(25)} is fulfilled. In order to make the best use of the inequality it is important that r be judiciously chosen, the object being to minimize the function $M(r)r^{-n}$, where $M(r)$ is the maximum of $|f|$ on $|\zeta - a| = r$.

\subsection*{EXERCISES}
\originalnumber{1}{exercise:4.2.3.4.1}
1. Compute
\begin{equation*}
\int_ {| z | = 1} e ^ {z} z ^ {- n} d z, \quad \int_ {| z | = 2} z ^ {n} (1 - z) ^ {m} d z, \quad \int_ {| z | = \rho} | z - a | ^ {- 4} | d z | (| a | \neq \rho).
\end{equation*}
\originalnumber{2}{exercise:4.2.3.4.2}
2. Prove that a function which is analytic in the whole plane and satisfies an inequality $|f(z)| < |z|^{n}$ for some n and all sufficiently large $|z|$ reduces to a polynomial.

\originalnumber{3}{exercise:4.2.3.4.3}
3. If $f(z)$ is analytic and $|f(z)| \leq M$ for $|z| \leq R$, find an upper bound for $|f^{(n)}(z)|$ in $|z| \leq \rho < R$.

\originalnumber{4}{exercise:4.2.3.4.4}
4. If $f(z)$ is analytic for $|z| < 1$ and $|f(z)| \leq 1 / (1 - |z|)$, find the best estimate of $|f^{(n)}(0)|$ that Cauchy's inequality will yield.

\originalnumber{5}{exercise:4.2.3.4.5}
5. Show that the successive derivatives of an analytic function at a point can never satisfy $|f^{(n)}(z)| > n!n^{n}$. Formulate a sharper theorem of the same kind.

\originalnumber{6}{exercise:4.2.3.4.6}
*6. A more general form of Lemma~\hyperref[lem:4.3]{3} reads as follows:

Let the function $\varphi(z,t)$ be continuous as a function of both variables when $z$ lies in a region $\Omega$ and $\alpha \leq t \leq \beta$. Suppose further that $\varphi(z,t)$ is analytic as a function of $z \in \Omega$ for any fixed $t$. Then
\begin{equation*}
F (z) = \int_ {\alpha} ^ {\beta} \varphi (z, t) d t
\end{equation*}
is analytic in $z$ and
\begin{equation*}
F ^ {\prime} (z) = \int_ {\alpha} ^ {\beta} \frac {\partial \varphi (z , t)}{\partial z} d t.\tag{26}\label{eq:4.26}
\end{equation*}
To prove this represent $\varphi(z,t)$ as a Cauchy integral
\begin{equation*}
\varphi (z, t) = \frac {1}{2 \pi i} \int_ {C} \frac {\varphi (\zeta , t)}{\zeta - z} d \zeta .
\end{equation*}
Fill in the necessary details to obtain
\begin{equation*}
F (z) = \int_ {C} \left(\frac {1}{2 \pi i} \int_ {\alpha} ^ {\beta} \varphi (\zeta , t) d t\right) \frac {d \zeta}{\zeta - z}
\end{equation*}
and use Lemma~\hyperref[lem:4.3]{3} to prove \hyperref[eq:4.26]{(26)}.


% Original PDF physical page 140
\index{Cauchy's integral formula|)}\index{Cauchy's integral theorem|)}\sourcepage{140}
\setcounter{section}{2}\section{Local Properties Of Analytic Functions}\label{sec:4.3}
We have already proved that an analytic function has derivatives of all orders. In this section we will make a closer study of the local properties. It will include a classification of the isolated singularities\index{Isolated singularity} of analytic functions.

\setcounter{section}{3}\setcounter{subsection}{0}\subsection{Removable Singularities. Taylor's Theorem}\label{sec:4.3.1}
In Theorem~\hyperref[thm:4.3]{3} we introduced a weaker condition which could be substituted for analyticity at a finite number of points without affecting the end result. We showed moreover, in Theorem~\hyperref[thm:4.5]{5}, that Cauchy's theorem in a circular disk remains true under these weaker conditions. This was an essential point in our derivation of Cauchy's integral formula, for we were required to apply Cauchy's theorem to a function of the form $(f(z) - f(a))/(z - a)$.

Finally, it was pointed out that Cauchy's integral formula remains valid in the presence of a finite number of exceptional points, all satisfying the fundamental condition of Theorem~\hyperref[thm:4.3]{3}, provided that none of them coincides with $a$. This remark is more important than it may seem on the surface. Indeed, Cauchy's formula provides us with a representation of $f(z)$ through an integral which in its dependence on $z$ has the same character at the exceptional points as everywhere else. It follows that the exceptional points are such only by lack of information, and not by their intrinsic nature. Points with this character are called removable singularities\index{Removable singularity|(}\index{Removable singularity}. We shall prove the following precise theorem:

\begin{originalstatement}{Theorem}{7}\label{thm:4.7}
Suppose that $f(z)$ is analytic in the region $\Omega'$ obtained by omitting a point $a$ from a region $\Omega$. A necessary and sufficient condition that there exist an analytic function in $\Omega$ which coincides with $f(z)$ in $\Omega'$ is that $\lim_{z \to a} (z - a)f(z) = 0$. The extended function is uniquely determined.
\end{originalstatement}
The necessity and the uniqueness are trivial since the extended function must be continuous at a. To prove the sufficiency we draw a circle C about a so that C and its inside are contained in $\Omega$. Cauchy's formula is valid, and we can write
\begin{equation*}
f (z) = \frac {1}{2 \pi i} \int_ {C} \frac {f (\zeta) d \zeta}{\zeta - z}
\end{equation*}
for all $z \neq a$ inside of C. But the integral in the right-hand member represents an analytic function of z throughout the inside of C. Consequently, the function which is equal to $f(z)$ for $z \neq a$ and which has the value
\begin{equation*}
\frac {1}{2 \pi i} \int_ {C} \frac {f (\zeta) d \zeta}{\zeta - a}\tag{27}\label{eq:4.27}
\end{equation*}
% Original PDF physical page 141
\sourcepage{141}
for $z = a$ is analytic in $\Omega$. It is natural to denote the extended function by $f(z)$ and the value \hyperref[eq:4.27]{(27)} by $f(a)$.

We apply this result to the function
\begin{equation*}
F (z) = \frac {f (z) - f (a)}{z - a}
\end{equation*}
used in the proof of Cauchy's formula. It is not defined for $z = a$, but it satisfies the condition $\lim_{z \to a} (z - a) F(z) = 0$. The limit of $F(z)$ as $z$ tends to $a$ is $f'(a)$. Hence there exists an analytic function which is equal to $F(z)$ for $z \neq a$ and equal to $f'(a)$ for $z = a$. Let us denote this function by $f_1(z)$. Repeating the process we can define an analytic function $f_2(z)$ which equals $(f_1(z) - f_1(a)) / (z - a)$ for $z \neq a$ and $f_1'(a)$ for $z = a$, and so on.

The recursive scheme by which $f_{n}(z)$ is defined can be written in the form
\begin{equation*}
\begin{array}{l} f (z) = f (a) + (z - a) f _ {1} (z) \\ f _ {1} (z) = f _ {1} (a) + (z - a) f _ {2} (z) \\ \dots \dots \\ f _ {n - 1} (z) = f _ {n - 1} (a) + (z - a) f _ {n} (z). \end{array}
\end{equation*}
From these equations which are trivially valid also for $z = a$ we obtain
\begin{equation*}
\begin{aligned}
f(z)={}&f(a)+(z-a)f_1(a)+(z-a)^2f_2(a)+\dots\\
&+(z-a)^{n-1}f_{n-1}(a)+(z-a)^nf_n(z).
\end{aligned}
\end{equation*}
Differentiating $n$ times and setting $z = a$ we find
\begin{equation*}
f ^ {(n)} (a) = n! f _ {n} (a).
\end{equation*}
This determines the coefficients $f_{n}(a)$, and we obtain the following form of Taylor's theorem\index{Taylor's theorem}:

\begin{originalstatement}{Theorem}{8}\label{thm:4.8}
If $f(z)$ is analytic in a region $\Omega$, containing $a$, it is possible to write
\begin{equation*}
\begin{aligned}
f(z)={}&f(a)+\frac{f'(a)}{1!}(z-a)+\frac{f''(a)}{2!}(z-a)^2+\dots\\
&+\frac{f^{(n-1)}(a)}{(n-1)!}(z-a)^{n-1}+f_n(z)(z-a)^n,
\end{aligned}\tag{28}\label{eq:4.28}
\end{equation*}
where $f_{n}(z)$ is analytic in $\Omega$.

\end{originalstatement}

% Original PDF physical page 142
\sourcepage{142}
This finite development must be well distinguished from the infinite Taylor series which we will study later. It is, however, the finite development \hyperref[eq:4.28]{(28)} which is the most useful for the study of the local properties of $f(z)$. Its usefulness is enhanced by the fact that $f_{n}(z)$ has a simple explicit expression as a line integral.

Using the same circle C as before we have first
\begin{equation*}
f _ {n} (z) = \frac {1}{2 \pi i} \int_ {C} \frac {f _ {n} (\zeta) d \zeta}{\zeta - z}.
\end{equation*}
For $f_{n}(\zeta)$ we substitute the expression obtained from \hyperref[eq:4.28]{(28)}. There will be one main term containing $f(\zeta)$. The remaining terms are, except for constant factors, of the form
\begin{equation*}
F _ {\nu} (a) = \int_ {C} \frac {d \zeta}{(\zeta - a) ^ {\nu} (\zeta - z)}, \quad \nu \geq 1.
\end{equation*}
But
\begin{equation*}
F _ {1} (a) = \frac {1}{z - a} \int_ {C} \left(\frac {1}{\zeta - z} - \frac {1}{\zeta - a}\right) d \zeta = 0,
\end{equation*}
identically for all $a$ inside of $C$. By Lemma~\hyperref[lem:4.3]{3} we have $F_{\nu + 1}(a) = F_1^{(\nu)}(a) / \nu!$ and thus $F_{\nu}(a) = 0$ for all $\nu \geq 1$. Hence the expression for $f_n(z)$ reduces to
\begin{equation*}
f _ {n} (z) = \frac {1}{2 \pi i} \int_ {C} \frac {f (\zeta) d \zeta}{(\zeta - a) ^ {n} (\zeta - z)}.\tag{29}\label{eq:4.29}
\end{equation*}
The representation is valid inside of C.

\setcounter{section}{3}\setcounter{subsection}{1}\subsection{Zeros and Poles}\label{sec:4.3.2}
If $f(a)$ and all derivatives $f^{(\nu)}(a)$ vanish, we can write by \hyperref[eq:4.28]{(28)}
\begin{equation*}
f (z) = f _ {n} (z) (z - a) ^ {n}\tag{30}\label{eq:4.30}
\end{equation*}
for any n. An estimate for $f_{n}(z)$ can be obtained by \hyperref[eq:4.29]{(29)}. The disk with the circumference C has to be contained in the region $\Omega$ in which $f(z)$ is defined and analytic. The absolute value $|f(z)|$ has a maximum M on C; if the radius of C is denoted by R, we find
\begin{equation*}
\left| f _ {n} (z) \right| \leq \frac {M}{R ^ {n - 1} (R - | z - a |)}
\end{equation*}
for $|z - a| < R$. By \hyperref[eq:4.30]{(30)} we have thus
\begin{equation*}
| f (z) | \leq \left(\frac {| z - a |}{R}\right) ^ {n} \cdot \frac {M R}{R - | z - a |}.
\end{equation*}
But $(|z - a| / R)^n \to 0$ for $n \to \infty$, since $|z - a| < R$. Hence $f(z) = 0$ inside of $C$.


% Original PDF physical page 143
\index{Removable singularity|)}\sourcepage{143}
We\index{Function!regular|(}\index{Order, algebraic!of zero|(}\index{Regular function|(} show now that $f(z)$ is identically zero\index{Zero} in all of $\Omega$. Let $E_{1}$ be the set on which $f(z)$ and all derivatives vanish and $E_{2}$ the set on which the function or one of the derivatives is different from zero. $E_{1}$ is open by the above reasoning, and $E_{2}$ is open because the function and all derivatives are continuous. Therefore either $E_{1}$ or $E_{2}$ must be empty. If $E_{2}$ is empty, the function is identically zero. If $E_{1}$ is empty, $f(z)$ can never vanish together with all its derivatives.

Assume that $f(z)$ is not identically zero. Then, if $f(a) = 0$, there exists a first derivative $f^{(h)}(a)$ which is different from zero. We say then that $a$ is a zero of order $h$, and the result that we have just proved expresses that there are no zeros of infinite order. In this respect an analytic function has the same local behavior as a polynomial, and just as in the case of polynomials we find that it is possible to write $f(z) = (z - a)^{h}f_{h}(z)$ where $f_{h}(z)$ is analytic and $f_{h}(a) \neq 0$.

In the same situation, since $f_{h}(z)$ is continuous, $f_{h}(z) \neq 0$ in a neighborhood of a and z = a is the only zero of $f(z)$ in this neighborhood. In other words, the zeros of an analytic function which does not vanish identically are isolated. This property can also be formulated as a uniqueness theorem: If $f(z)$ and $g(z)$ are analytic in $\Omega$, and if $f(z) = g(z)$ on a set which has an accumulation point in $\Omega$, then $f(z)$ is identically equal to $g(z)$. The conclusion follows by consideration of the difference $f(z) - g(z)$.

Particular instances of this result which deserve to be quoted are the following: If $f(z)$ is identically zero in a subregion of $\Omega$, then it is identically zero in $\Omega$, and the same is true if $f(z)$ vanishes on an arc which does not reduce to a point. We can also say that an analytic function is uniquely determined by its values on any set with an accumulation point in the region of analyticity. This does not mean that we know of any way in which the values of the function can be computed.

We consider now a function $f(z)$ which is analytic in a neighborhood of a, except perhaps at a itself. In other words, $f(z)$ shall be analytic in a region $0 < |z - a| < \delta$. The point a is called an isolated singularity of $f(z)$. We have already treated the case of a removable singularity. Since we can then define $f(a)$ so that $f(z)$ becomes analytic in the disk $|z - a| < \delta$, it needs no further consideration.\hyperlink{fn:143:dagger}{\textsuperscript{\ensuremath{\dagger}}}

If $\lim_{z\to a}f(z)=\infty$, the point $a$ is said to be a \textit{pole}\index{Pole} of $f(z)$, and we set $f(a)=\infty$. There exists a $\delta' \leq \delta$ such that $f(z) \neq 0$ for $0 < |z - a| < \delta'$. In this region the function $g(z) = 1/f(z)$ is defined and analytic. But the singularity of $g(z)$ at $a$ is removable, and $g(z)$ has an analytic exten%
\begingroup\renewcommand{\thefootnote}{}\footnotetext{\hypertarget{fn:143:dagger}{}\textsuperscript{\ensuremath{\dagger}} If $a$ is a removable singularity, $f(z)$ is frequently said to be \emph{regular} at $a$; this term is sometimes used as a synonym for analytic.}\endgroup%
% Original PDF physical page 144
\index{Function!regular|)}\index{Order, algebraic!of zero|)}\index{Regular function|)}\sourcepage{144}%
sion\index{Order, algebraic!of a pole|(} with $g(a) = 0$. Since $g(z)$ does not vanish identically, the zero at a has a finite order, and we can write $g(z) = (z - a)^{h}g_{h}(z)$ with $g_{h}(a) \neq 0$. The number h is the order of the pole, and $f(z)$ has the representation $f(z) = (z - a)^{-h}f_{h}(z)$ where $f_{h}(z) = 1/g_{h}(z)$ is analytic and different from zero in a neighborhood of a. The nature of a pole is thus exactly the same as in the case of a rational function.

A function $f(z)$ which is analytic in a region $\Omega$, except for poles, is said to be meromorphic in $\Omega$. More precisely, to every $a \in \Omega$ there shall exist a neighborhood $|z - a| < \delta$, contained in $\Omega$, such that either $f(z)$ is analytic in the whole neighborhood, or else $f(z)$ is analytic for $0 < |z - a| < \delta$, and the isolated singularity is a pole. Observe that the poles of a meromorphic function\index{Meromorphic function} are isolated by definition. The quotient $f(z)/g(z)$ of two analytic functions in $\Omega$ is a meromorphic function in $\Omega$, provided that $g(z)$ is not identically zero. The only possible poles are the zeros of $g(z)$, but a common zero of $f(z)$ and $g(z)$ can also be a removable singularity. If this is the case, the value of the quotient must be determined by continuity. More generally, the sum, the product, and the quotient of two meromorphic functions are meromorphic. The case of an identically vanishing denominator must be excluded, unless we wish to consider the constant $\infty$ as a meromorphic function.

For a more detailed discussion of isolated singularities, we consider the conditions (1) $\lim_{z\to a}|z-a|^{\alpha}|f(z)|=0$, (2) $\lim_{z\to a}|z-a|^{\alpha}|f(z)|=\infty$, for real values of $\alpha$. If (1) holds for a certain $\alpha$, then it holds for all larger $\alpha$, and hence for some integer m. Then $(z-a)^{m}f(z)$ has a removable singularity and vanishes for z=a. Either $f(z)$ is identically zero, in which case (1) holds for all $\alpha$, or $(z-a)^{m}f(z)$ has a zero of finite order k. In the latter case it follows at once that (1) holds for all $\alpha>h=m-k$, while (2) holds for all $\alpha<h$. Assume now that (2) holds for some $\alpha$; then it holds for all smaller $\alpha$, and hence for some integer n. The function $(z-a)^{n}f(z)$ has a pole of finite order l, and setting $h=n+l$ we find again that (1) holds for $\alpha>h$ and (2) for $\alpha<h$. The discussion shows that there are three possibilities: (i) condition (1) holds for all $\alpha$, and $f(z)$ vanishes identically; (ii) there exists an integer h such that (1) holds for $\alpha>h$ and (2) for $\alpha<h$; (iii) neither (1) nor (2) holds for any $\alpha$.

Case (i) is uninteresting. In case (ii) h may be called the algebraic order\index{Order, algebraic} of $f(z)$ at a. It is positive in case of a pole, negative in case of a zero, and zero if $f(z)$ is analytic but $\neq 0$ at a. The remarkable thing is that the order is always an integer; there is no single-valued analytic function which tends to 0 or $\infty$ like a fractional power of $|z - a|$.

In the case of a pole of order h, let us apply Theorem~\hyperref[thm:4.8]{8} to the analytic function $(z - a)^{h}f(z)$. We obtain a development of the form
\begin{equation*}
(z - a) ^ {h} f (z) = B _ {h} + B _ {h - 1} (z - a) + \dots + B _ {1} (z - a) ^ {h - 1} + \varphi (z) (z - a) ^ {h}
\end{equation*}
% Original PDF physical page 145
\index{Order, algebraic!of a pole|)}\sourcepage{145}
where $\varphi(z)$ is analytic at $z = a$. For $z \neq a$ we can divide by $(z - a)^{h}$ and find
\begin{equation*}
f (z) = B _ {h} (z - a) ^ {- h} + B _ {h - 1} (z - a) ^ {- h + 1} + \dots + B _ {1} (z - a) ^ {- 1} + \varphi (z).
\end{equation*}
The part of this development which precedes $\varphi(z)$ is called the singular part of $f(z)$ at z = a. A pole has thus not only an order, but also a well-defined singular part. The difference of two functions with the same singular part is analytic at a.

In case (iii) the point a is an essential isolated singularity. In the neighborhood of an essential singularity\index{Essential singularity} $f(z)$ is at the same time unbounded and comes arbitrarily close to zero. As a characterization of the complicated behavior of a function in the neighborhood of an essential singularity, we prove the following classical theorem of Weierstrass\index{Weierstrass K@Weierstrass, K.}:

\begin{originalstatement}{Theorem}{9}\label{thm:4.9}
An analytic function comes arbitrarily close to any complex value in every neighborhood of an essential singularity.
\end{originalstatement}
If the assertion were not true, we could find a complex number $A$ and a $\delta > 0$ such that $|f(z) - A| > \delta$ in a neighborhood of $a$ (except for $z = a$). For any $\alpha < 0$ we have then $\lim_{z \to a} |z - a|^{\alpha}| f(z) - A| = \infty$. Hence $a$ would not be an essential singularity of $f(z) - A$. Accordingly, there exists a $\beta$ with $\lim_{z \to a} |z - a|^{\beta}| f(z) - A| = 0$, and we are free to choose $\beta > 0$. Since in that case $\lim_{z \to a} |z - a|^{\beta}| A| = 0$ it would follow that $\lim_{z \to a} |z - a|^{\beta}| f(z)| = 0$, and $a$ would not be an essential singularity of $f(z)$. The contradiction proves the theorem.

The notion of isolated singularity applies also to functions which are analytic in a neighborhood $|z| > R$ of $\infty$. Since $f(\infty)$ is not defined, we treat $\infty$ as an isolated singularity, and by convention it has the same character of removable singularity, pole, or essential singularity as the singularity of $g(z) = f(1/z)$ at z = 0. If the singularity is nonessential, $f(z)$ has an algebraic order h such that $\lim_{z \to \infty} z^{-h} f(z)$ is neither zero nor infinity, and for a pole the singular part is a polynomial in z. If $\infty$ is an essential singularity, the function has the property expressed by Theorem~\hyperref[thm:4.9]{9} in every neighborhood of infinity.

\subsection*{EXERCISES}
\originalnumber{1}{exercise:4.3.2.5.1}
1. If $f(z)$ and $g(z)$ have the algebraic orders h and k at z = a, show that fg has the order $h + k$, f/g the order h - k, and $f + g$ an order which does not exceed max $(h, k)$.


% Original PDF physical page 146
\sourcepage{146}
\originalnumber{2}{exercise:4.3.2.5.2}
2. Show that a function which is analytic in the whole plane and has a nonessential singularity at $\infty$ reduces to a polynomial.

\originalnumber{3}{exercise:4.3.2.5.3}
3. Show that the functions $e^{z}$, $\sin z$ and $\cos z$ have essential singularities at $\infty$.

\originalnumber{4}{exercise:4.3.2.5.4}
4. Show that any function which is meromorphic in the extended plane is rational.

\originalnumber{5}{exercise:4.3.2.5.5}
5. Prove that an isolated singularity of $f(z)$ is removable as soon as either $\operatorname{Re} f(z)$ or $\operatorname{Im} f(z)$ is bounded above or below. Hint: Apply a fractional linear transformation.

\originalnumber{6}{exercise:4.3.2.5.6}
6. Show that an isolated singularity of $f(z)$ cannot be a pole of $\exp f(z)$. Hint: $f$ and $e^f$ cannot have a common pole (why?). Now apply Theorem~\hyperref[thm:4.9]{9}.

\setcounter{section}{3}\setcounter{subsection}{2}\subsection{The Local Mapping}\index{Local mapping|(}\index{Mapping!local|(}\index{Local mapping}\index{Mapping!local}\label{sec:4.3.3}
We begin with the proof of a general formula which enables us to determine the number of zeros of an analytic function. We are considering a function $f(z)$ which is analytic and not identically zero in an open disk $\Delta$. Let $\gamma$ be a closed curve in $\Delta$ such that $f(z) \neq 0$ on $\gamma$. For the sake of simplicity we suppose first that $f(z)$ has only a finite number of zeros in $\Delta$, and we agree to denote them by $z_1, z_2, \ldots, z_n$ where each zero is repeated as many times as its order indicates.

By repeated applications of Theorem~\hyperref[thm:4.8]{8}, or rather its consequence \hyperref[eq:4.30]{(30)}, it is clear that we can write $f(z) = (z - z_1)(z - z_2) \cdots (z - z_n)g(z)$ where $g(z)$ is analytic and $\neq 0$ in $\Delta$. Forming the logarithmic derivative we obtain
\begin{equation*}
\frac {f ^ {\prime} (z)}{f (z)} = \frac {1}{z - z _ {1}} + \frac {1}{z - z _ {2}} + \dots + \frac {1}{z - z _ {n}} + \frac {g ^ {\prime} (z)}{g (z)}
\end{equation*}
for $z \neq z_{j}$, and particularly on $\gamma$. Since $g(z) \neq 0$ in $\Delta$, Cauchy's theorem yields
\begin{equation*}
\int_ {\gamma} \frac {g ^ {\prime} (z)}{g (z)} d z = 0.
\end{equation*}
Recalling the definition of $n(\gamma, z_j)$ we find
\begin{equation*}
n (\gamma , z _ {1}) + n (\gamma , z _ {2}) + \dots + n (\gamma , z _ {n}) = \frac {1}{2 \pi i} \int_ {\gamma} \frac {f ^ {\prime} (z)}{f (z)} d z.\tag{31}\label{eq:4.31}
\end{equation*}
This is still true if $f(z)$ has infinitely many zeros in $\Delta$. It is clear that $\gamma$ is contained in a concentric disk $\Delta'$ smaller than $\Delta$. Unless $f(z)$ is identically zero, a case which must obviously be excluded, it has only a finite number of zeros in $\Delta'$. This is an obvious consequence of the Bolzano-Weierstrass theorem, for if there were infinitely many zeros they would have an accumulation point in the closure of $\Delta'$, and this is impossible. We can now apply \hyperref[eq:4.31]{(31)} to the disk $\Delta'$. The zeros outside of $\Delta'$ satisfy $n(\gamma,z_{j})=0$ and hence do not contribute to the sum in \hyperref[eq:4.31]{(31)}. We have thus proved:


% Original PDF physical page 147
\sourcepage{147}
\begin{originalstatement}{Theorem}{10}\label{thm:4.10}
Let $z_{j}$ be the zeros of a function $f(z)$ which is analytic in a disk $\Delta$ and does not vanish identically, each zero being counted as many times as its order indicates. For every closed curve $\gamma$ in $\Delta$ which does not pass through a zero
\begin{equation*}
\sum_ {j} n (\gamma , z _ {j}) = \frac {1}{2 \pi i} \int_ {\gamma} \frac {f ^ {\prime} (z)}{f (z)} d z,\tag{32}\label{eq:4.32}
\end{equation*}
where the sum has only a finite number of terms $\neq 0$.

\end{originalstatement}
The function $w = f(z)$ maps $\gamma$ onto a closed curve $\Gamma$ in the $w$-plane, and we find
\begin{equation*}
\int_ {\Gamma} \frac {d w}{w} = \int_ {\gamma} \frac {f ^ {\prime} (z)}{f (z)} d z.
\end{equation*}
The formula \hyperref[eq:4.32]{(32)} has thus the following interpretation:
\begin{equation*}
n (\Gamma , 0) = \sum_ {j} n (\gamma , z _ {j}).\tag{33}\label{eq:4.33}
\end{equation*}
The simplest and most useful application is to the case where it is known beforehand that each $n(\gamma,z_{j})$ must be either 0 or 1. Then \hyperref[eq:4.32]{(32)} yields a formula for the total number of zeros enclosed by $\gamma$. This is evidently the case when $\gamma$ is a circle.

Let $a$ be an arbitrary complex value, and apply Theorem~\hyperref[thm:4.10]{10} to $f(z) - a$. The zeros of $f(z) - a$ are the roots of the equation $f(z) = a$, and we denote them by $z_j(a)$. In the place of \hyperref[eq:4.32]{(32)} we obtain the formula
\begin{equation*}
\sum_ {j} n (\gamma , z _ {j} (a)) = \frac {1}{2 \pi i} \int_ {\gamma} \frac {f ^ {\prime} (z)}{f (z) - a} d z
\end{equation*}
and \hyperref[eq:4.33]{(33)} takes the form
\begin{equation*}
n (\Gamma , a) = \sum_ {j} n (\gamma , z _ {j} (a)).
\end{equation*}
It is necessary to assume that $f(z) \neq a$ on $\gamma$.

If $a$ and $b$ are in the same region determined by $\Gamma$, we know that $n(\Gamma, a) = n(\Gamma, b)$, and hence we have also $\sum_{j} n(\gamma, z_j(a)) = \sum_{j} n(\gamma, z_j(b))$. If $\gamma$ is a circle, it follows that $f(z)$ takes the values $a$ and $b$ equally many times inside of $\gamma$. The following theorem on local correspondence is an immediate consequence of this result.

\begin{originalstatement}{Theorem}{11}\label{thm:4.11}
Suppose that $f(z)$ is analytic at $z_0$, $f(z_0) = w_0$, and that $f(z) - w_0$ has a zero of order $n$ at $z_0$. If $\varepsilon > 0$ is sufficiently small, there exists a corresponding $\delta > 0$ such that for all $a$ with $|a - w_0| < \delta$ the equation $f(z) = a$ has exactly $n$ roots in the disk $|z - z_0| < \varepsilon$.
\end{originalstatement}

% Original PDF physical page 148
\sourcepage{148}
We can choose $\varepsilon$ so that $f(z)$ is defined and analytic for $|z - z_0| \leq \varepsilon$ and so that $z_0$ is the only zero of $f(z) - w_0$ in this disk. Let $\gamma$ be the circle $|z - z_0| = \varepsilon$ and $\Gamma$ its image under the mapping $w = f(z)$. Since $w_0$ belongs to the complement of the closed set $\Gamma$, there exists a neighborhood $|w - w_0| < \delta$ which does not intersect $\Gamma$ (Fig.~\hyperref[fig:4-7]{4-7}). It follows immediately that all values $a$ in this neighborhood are taken the same number of times inside of $\gamma$. The equation $f(z) = w_0$ has exactly $n$ coinciding roots inside of $\gamma$, and hence every value $a$ is taken $n$ times. It is understood that multiple roots are counted according to their multiplicity, but if $\varepsilon$ is sufficiently small we can assert that all roots of the equation $f(z) = a$ are simple for $a \neq w_0$. Indeed, it is sufficient to choose $\varepsilon$ so that $f'(z)$ does not vanish for $0 < |z - z_0| < \varepsilon$.

\begin{originalstatement}{Corollary}{1}\label{cor:4.1}
A nonconstant analytic function maps open sets onto open sets.
\end{originalstatement}
This is merely another way of saying that the image of every sufficiently small disk $|z - z_{0}| < \varepsilon$ contains a neighborhood $|w - w_{0}| < \delta$.

In the case n = 1 there is one-to-one correspondence between the disk $|w - w_{0}| < \delta$ and an open subset $\Delta$ of $|z - z_{0}| < \varepsilon$. Since open sets in the z-plane correspond to open sets in the w-plane the inverse function of $f(z)$ is continuous, and the mapping is topological. The mapping can be restricted to a neighborhood of $z_{0}$ contained in $\Delta$, and we are able to state:

\begin{originalstatement}{Corollary}{2}\label{cor:4.2}
If $f(z)$ is analytic at $z_0$ with $f'(z_0) \neq 0$, it maps a neighborhood of $z_0$ conformally and topologically onto a region.
\end{originalstatement}
From the continuity of the inverse function it follows in the usual way that the inverse function is analytic, and hence the inverse mapping is

\bookfigure{148}{1}{FIG. 4-7. Local correspondence.}

% Original PDF physical page 149
\sourcepage{149}
\bookfigure{149}{1}{FIG. 4-8. Branch point: n = 3.}
likewise conformal. Conversely, if the local mapping\index{Local mapping}\index{Mapping!local} is one to one, Theorem~\hyperref[thm:4.11]{11} can hold only with n = 1, and hence $f'(z_0)$ must be different from zero.

For n > 1 the local correspondence can still be described in very precise terms. Under the assumption of Theorem~\hyperref[thm:4.11]{11} we can write
\begin{equation*}
f (z) - w _ {0} = (z - z _ {0}) ^ {n} g (z)
\end{equation*}
where $g(z)$ is analytic at $z_0$ and $g(z_0) \neq 0$. Choose $\varepsilon > 0$ so that $|g(z) - g(z_0)| < |g(z_0)|$ for $|z - z_0| < \varepsilon$. In this neighborhood it is possible to define a single-valued analytic branch of $\sqrt[n]{g(z)}$, which we denote by $h(z)$. We have thus
\begin{equation*}
\begin{array}{r l} f (z) - w _ {0} & = \zeta (z) ^ {n} \\ \zeta (z) & = (z - z _ {0}) h (z). \end{array}
\end{equation*}
Since $\zeta'(z_0) = h(z_0) \neq 0$ the mapping $\zeta = \zeta(z)$ is topological in a neighborhood of $z_0$. On the other hand, the mapping $w = w_0 + \zeta^n$ is of an elementary character and determines $n$ equally spaced values $\zeta$ for each value of $w$. By performing the mapping in two steps we obtain a very illuminating picture of the local correspondence. Figure~\hyperref[fig:4-8]{4-8} shows the inverse image of a small disk and the $n$ arcs which are mapped onto the positive radius.

\subsection*{EXERCISES}
\originalnumber{1}{exercise:4.3.3.6.1}
1. Determine explicitly the largest disk about the origin whose image under the mapping $w = z^{2} + z$ is one to one.

\originalnumber{2}{exercise:4.3.3.6.2}
2. Same problem for $w = e^{z}$.

\originalnumber{3}{exercise:4.3.3.6.3}
3. Apply the representation $f(z) = w_0 + \zeta(z)^n$ to $\cos z$ with $z_0 = 0$. Determine $\zeta(z)$ explicitly.

\originalnumber{4}{exercise:4.3.3.6.4}
4. If $f(z)$ is analytic at the origin and $f'(0) \neq 0$, prove the existence of an analytic $g(z)$ such that $f(z^{n}) = f(0) + g(z)^{n}$ in a neighborhood of 0.

\setcounter{section}{3}\setcounter{subsection}{3}\subsection{The Maximum Principle}\label{sec:4.3.4}
Corollary~\hyperref[cor:4.1]{1} of Theorem~\hyperref[thm:4.11]{11} has a very important analytical consequence known as the maximum principle\index{Maximum principle|(}\index{Maximum principle} for
% Original PDF physical page 150
\index{Local mapping|)}\index{Mapping!local|)}\sourcepage{150}
analytic functions. Because of its simple and explicit formulation it is one of the most useful general theorems in the theory of functions. As a rule all proofs based on the maximum principle\index{Maximum principle} are very straightforward, and preference is quite justly given to proofs of this kind.

\begin{originalstatement}{Theorem}{12}\label{thm:4.12}
(The maximum principle.) If $f(z)$ is analytic and non-constant in a region $\Omega$, then its absolute value $|f(z)|$ has no maximum in $\Omega$.
\end{originalstatement}
The proof is clear. If $w_{0} = f(z_{0})$ is any value taken in $\Omega$, there exists a neighborhood $|w - w_{0}| < \varepsilon$ contained in the image of $\Omega$. In this neighborhood there are points of modulus $> |w_{0}|$, and hence $|f(z_{0})|$ is not the maximum of $|f(z)|$.

In a positive formulation essentially the same theorem can be stated in the form:

\begin{originalstatement}{Theorem}{12'}\label{thm:4.12prime}
If $f(z)$ is defined and continuous on a closed bounded set $E$ and analytic on the interior of $E$, then the maximum of $|f(z)|$ on $E$ is assumed on the boundary of $E$.
\end{originalstatement}

Since E is compact, $|f(z)|$ has a maximum on E. Suppose that it is assumed at $z_{0}$. If $z_{0}$ is on the boundary, there is nothing to prove. If $z_{0}$ is an interior point, then $|f(z_{0})|$ is also the maximum of $|f(z)|$ in a disk $|z - z_{0}| < \delta$ contained in E. But this is not possible unless $f(z)$ is constant in the component of the interior of E which contains $z_{0}$. It follows by continuity that $|f(z)|$ is equal to its maximum on the whole boundary of that component. This boundary is not empty and it is contained in the boundary of E. Thus the maximum is always assumed at a boundary point.

The maximum principle can also be proved analytically, as a consequence of Cauchy's integral formula. If the formula \hyperref[eq:4.22]{(22)} is specialized to the case where $\gamma$ is a circle of center $z_{0}$ and radius r, we can write $\zeta = z_{0} + re^{i\theta}$, $d\zeta = ire^{i\theta}d\theta$ on $\gamma$ and obtain for $z = z_{0}$
\begin{equation*}
f (z _ {0}) = \frac {1}{2 \pi} \int_ {0} ^ {2 \pi} f (z _ {0} + r e ^ {i \theta}) d \theta .\tag{34}\label{eq:4.34}
\end{equation*}
This formula shows that the value of an analytic function at the center of a circle is equal to the arithmetic mean of its values on the circle, subject to the condition that the closed disk $|z - z_{0}| \leq r$ is contained in the region of analyticity.

From \hyperref[eq:4.34]{(34)} we derive the inequality
\begin{equation*}
\left| f (z _ {0}) \right| \leq \frac {1}{2 \pi} \int_ {0} ^ {2 \pi} \left| f (z _ {0} + r e ^ {i \theta}) \right| d \theta .\tag{35}\label{eq:4.35}
\end{equation*}
% Original PDF physical page 151
\index{Complex integration|)}\sourcepage{151}
Suppose\index{Schwarz H A@Schwarz, H. A.!lemma of|(} that $|f(z_0)|$ were a maximum. Then we would have $|f(z_0 + re^{i\theta})| \leq |f(z_0)|$, and if the strict inequality held for a single value of $\theta$ it would hold, by continuity, on a whole arc. But then the mean value of $|f(z_0 + re^{i\theta})|$ would be strictly less than $|f(z_0)|$, and \hyperref[eq:4.35]{(35)} would lead to the contradiction $|f(z_0)| < |f(z_0)|$. Thus $|f(z)|$ must be constantly equal to $|f(z_0)|$ on all sufficiently small circles $|z - z_0| = r$ and, hence, in a neighborhood of $z_0$. It follows easily that $f(z)$ must reduce to a constant. This reasoning provides a second proof of the maximum principle\index{Maximum principle}. We have given preference to the first proof because it shows that the maximum principle is a consequence of the topological properties of the mapping by an analytic function.

Consider now the case of a function $f(z)$ which is analytic in the open disk $|z| < R$ and continuous on the closed disk $|z| \leq R$. If it is known that $|f(z)| \leq M$ on $|z| = R$, then $|f(z)| \leq M$ in the whole disk. The equality can hold only if $f(z)$ is a constant of absolute value M. Therefore, if it is known that $f(z)$ takes some value of modulus < M, it may be expected that a better estimate can be given. Theorems to this effect are very useful. The following particular result is known as the lemma of Schwarz\index{Schwarz H A@Schwarz, H. A.}:

\begin{originalstatement}{Theorem}{13}\label{thm:4.13}
If $f(z)$ is analytic for $|z| < 1$ and satisfies the conditions $|f(z)| \leq 1$, $f(0) = 0$, then $|f(z)| \leq |z|$ and $|f'(0)| \leq 1$. If $|f(z)| = |z|$ for some $z \neq 0$, or if $|f'(0)| = 1$, then $f(z) = cz$ with a constant $c$ of absolute value 1.
\end{originalstatement}
We apply the maximum principle to the function $f_{1}(z)$ which is equal to $f(z)/z$ for $z \neq 0$ and to $f'(0)$ for z = 0. On the circle $|z| = r < 1$ it is of absolute value $\leq 1/r$, and hence $|f_{1}(z)| \leq 1/r$ for $|z| \leq r$. Letting r tend to 1 we find that $|f_{1}(z)| \leq 1$ for all z, and this is the assertion of the theorem. If the equality holds at a single point, it means that $|f_{1}(z)|$ attains its maximum and, hence, that $f_{1}(z)$ must reduce to a constant.

The rather specialized assumptions of Theorem~\hyperref[thm:4.13]{13} are not essential, but should be looked upon as the result of a normalization. For instance, if $f(z)$ is known to satisfy the conditions of the theorem in a disk of radius $R$, the original form of the theorem can be applied to the function $f(Rz)$. As a result we obtain $|f(Rz)| \leq |z|$, which can be rewritten as $|f(z)| \leq |z| / R$. Similarly, if the upper bound of the modulus is $M$ instead of 1, we apply the theorem to $f(z) / M$ or, in the more general case, to $f(Rz) / M$. The resulting inequality is $|f(z)| \leq M |z| / R$.

Still more generally, we may replace the condition $f(0) = 0$ by an arbitrary condition $f(z_0) = w_0$ where $|z_0| < R$ and $|w_0| < M$. Let $\zeta = Tz$ be a linear transformation which maps $|z| < R$ onto $|\zeta| < 1$ with $z_0$ going into the origin, and let $Sw$ be a linear transformation with $Sw_0 = 0$ which maps $|w| < M$ onto $|Sw| < 1$. It is clear that the function $Sf(T^{-1}\zeta)$


% Original PDF physical page 152
\index{Schwarz H A@Schwarz, H. A.!lemma of|)}\sourcepage{152}
satisfies the hypothesis of the original theorem. Hence we obtain $|Sf(T^{-1}\zeta)| \leq |\zeta|$, or $|Sf(z)| \leq |Tz|$. Explicitly, this inequality can be written in the form
\begin{equation*}
\left| \frac {M (f (z) - w _ {0})}{M ^ {2} - \bar {w} _ {0} f (z)} \right| \leq \left| \frac {R (z - z _ {0})}{R ^ {2} - \bar {z} _ {0} z} \right|.\tag{36}\label{eq:4.36}
\end{equation*}
\subsection*{EXERCISES}
\originalnumber{1}{exercise:4.3.4.7.1}
1. Show by use of \hyperref[eq:4.36]{(36)}, or directly, that $|f(z)| \leq 1$ for $|z| \leq 1$ implies
\begin{equation*}
\frac {\left| f ^ {\prime} (z) \right|}{(1 - | f (z) | ^ {2})} \leq \frac {1}{1 - | z | ^ {2}}.
\end{equation*}
\originalnumber{2}{exercise:4.3.4.7.2}
2. If $f(z)$ is analytic and $\operatorname{Im} f(z) \geq 0$ for $\operatorname{Im} z > 0$, show that
\begin{equation*}
\frac {\left| f (z) - f \left(z _ {0}\right) \right|}{\left| f (z) - \overline {{{f \left(z _ {0}\right)}}} \right|} \leq \frac {\left| z - z _ {0} \right|}{\left| z - \bar {z} _ {0} \right|}
\end{equation*}
and
\begin{equation*}
\frac {| f ^ {\prime} (z) |}{\operatorname{Im} f (z)} \leq \frac {1}{y} \quad (z = x + i y).
\end{equation*}
\originalnumber{3}{exercise:4.3.4.7.3}
3. In Ex.~\hyperref[exercise:4.3.4.7.1]{1} and 2, prove that equality implies that $f(z)$ is a linear transformation.

\originalnumber{4}{exercise:4.3.4.7.4}
4. Derive corresponding inequalities if $f(z)$ maps $|z| < 1$ into the upper half plane.

\originalnumber{5}{exercise:4.3.4.7.5}
5. Prove by use of Schwarz's lemma that every one-to-one conformal mapping of a disk onto another (or a half plane) is given by a linear transformation.

\originalnumber{6}{exercise:4.3.4.7.6}
*6. If $\gamma$ is a piecewise differentiable arc contained in $|z| < 1$ the integral
\begin{equation*}
\int_ {\gamma} \frac {| d z |}{1 - | z | ^ {2}}
\end{equation*}
is called the noneuclidean length (or hyperbolic length) of $\gamma$. Show that an analytic function $f(z)$ with $|f(z)| < 1$ for $|z| < 1$ maps every $\gamma$ on an arc with smaller or equal noneuclidean length.

Deduce that a linear transformation of the unit disk onto itself preserves noneuclidean lengths, and check the result by explicit computation.

\originalnumber{7}{exercise:4.3.4.7.7}
*7. Prove that the arc of smallest noneuclidean length that joins two given points in the unit disk is a circular arc which is orthogonal to the unit circle. (Make use of a linear transformation that carries one end point to the origin, the other to a point on the positive real axis.)

The shortest noneuclidean length is called the noneuclidean distance\index{Distance!noneuclidean}\index{Noneuclidean distance}
% Original PDF physical page 153
\sourcepage{153}
between\index{Cauchy's integral theorem|(} the end points. Derive a formula for the noneuclidean distance between $z_{1}$ and $z_{2}$. Answer:
\begin{equation*}
\frac {1}{2} \log \frac {1 + \left| \frac {z _ {1} - z _ {2}}{1 - \bar {z} _ {1} z _ {2}} \right|}{1 - \left| \frac {z _ {1} - z _ {2}}{1 - \bar {z} _ {1} z _ {2}} \right|}.
\end{equation*}
\originalnumber{8}{exercise:4.3.4.7.8}
*8. How should noneuclidean length in the upper half plane be defined?

\setcounter{section}{3}\section{The General Form Of Cauchy's Theorem}\label{sec:4.4}
In our preliminary treatment of Cauchy's theorem and the integral formula we considered only the case of a circular region. For the purpose of studying the local properties of analytic functions this was quite adequate, but from a more general point of view we cannot be satisfied with a result which is so obviously incomplete. The generalization can proceed in two directions. For one thing we can seek to characterize the regions in which Cauchy's theorem has universal validity. Secondly, we can consider an arbitrary region and look for the curves $\gamma$ for which the assertion of Cauchy's theorem is true.

\setcounter{section}{4}\setcounter{subsection}{0}\subsection{Chains and Cycles}\index{Cycle|(}\index{Cycle}\label{sec:4.4.1}
In the first place we must generalize the notion of line integral. To this end we examine the equation
\begin{equation*}
\int_ {\gamma_ {1} + \gamma_ {2} + \dots + \gamma_ {n}} f d z = \int_ {\gamma_ {1}} f d z + \int_ {\gamma_ {2}} f d z + \dots + \int_ {\gamma_ {n}} f d z\tag{37}\label{eq:4.37}
\end{equation*}
which is valid when $\gamma_{1}, \gamma_{2}, \ldots, \gamma_{n}$ form a subdivision of the arc $\gamma$. Since the right-hand member of \hyperref[eq:4.37]{(37)} has a meaning for any finite collection, nothing prevents us from considering an arbitrary formal sum $\gamma_{1} + \gamma_{2} + \cdots + \gamma_{n}$, which need not be an arc, and we define the corresponding integral by means of equation \hyperref[eq:4.37]{(37)}. Such formal sums of arcs are called chains\index{Chain|(}\index{Chain}. It is clear that nothing is lost and much may be gained by considering line integrals over arbitrary chains.

Just as there is nothing unique about the way in which an arc can be subdivided, it is clear that different formal sums can represent the same chain. The guiding principle is that two chains should be considered identical if they yield the same line integrals for all functions f. If this principle is analyzed, we find that the following operations do not change the identity of a chain: (1) permutation of two arcs, (2) subdivision of an arc, (3) fusion of subarcs to a single arc, (4) reparametrization of an arc, (5) cancellation of opposite arcs. On this basis it would be easy to
% Original PDF physical page 154
\index{Maximum principle|)}\sourcepage{154}
formulate a logical equivalence relation which defines the identity of chains\index{Chain} in a formal manner. Inasmuch as the situation does not involve any logical pitfalls, we shall dispense with this formalization.

The sum of two chains is defined in the obvious way by juxtaposition. It is clear that the additive property \hyperref[eq:4.37]{(37)} of line integrals remains valid for arbitrary chains. When identical chains are added, it is convenient to denote the sum as a multiple. With this notation every chain can be written in the form
\begin{equation*}
\gamma = a _ {1} \gamma_ {1} + a _ {2} \gamma_ {2} + \dots + a _ {n} \gamma_ {n}\tag{38}\label{eq:4.38}
\end{equation*}
where the $a_{j}$ are positive integers and the $\gamma_{j}$ are all different. For opposite arcs we are allowed to write $a(-\gamma) = -a\gamma$ and continue the reduction of \hyperref[eq:4.38]{(38)} until no two $\gamma_{j}$ are opposite. The coefficients will be arbitrary integers, and terms with zero coefficients can be added at will. The last device enables us to express any two chains in terms of the same arcs, and their sum is obtained by adding corresponding coefficients. The zero chain is either an empty sum or a sum with all coefficients equal to zero.

A chain is a cycle\index{Cycle} if it can be represented as a sum of closed curves. Very simple combinatorial considerations show that a chain is a cycle if and only if in any representation the initial and end points of the individual arcs are identical in pairs. Thus it is immediately possible to tell whether a chain is a cycle or not.

In the applications we shall consider chains which are contained in a given open set $\Omega$. By this we mean that the chains have a representation by arcs in $\Omega$ and that only such representations will be considered. It is clear that all theorems which we have heretofore formulated only for closed curves in a region are in fact valid for arbitrary cycles in a region. In particular, the integral of an exact differential over any cycle is zero.

The index of a point with respect to a cycle is defined in exactly the same way as in the case of a single closed curve. It has the same properties, and in addition we can formulate the obvious but important additive law expressed by the equation $n(\gamma_{1} + \gamma_{2}, a) = n(\gamma_{1}, a) + n(\gamma_{2}, a)$.

\setcounter{section}{4}\setcounter{subsection}{1}\subsection{Simple Connectivity}\label{sec:4.4.2}
There is little doubt that all readers will know what we mean if we speak about a region without holes. Such regions are said to be simply connected, and it is for simply connected regions that Cauchy's theorem is universally valid. The suggestive language we have used cannot take the place of a mathematical definition, but fortunately very little is needed to make the term precise. Indeed, a region without holes is obviously one whose complement consists of a single piece. We are thus led to the following definition:


% Original PDF physical page 155
\index{Chain|)}\index{Cycle|)}\sourcepage{155}
\begin{originalstatement}{Definition}{1}\label{def:4.1}
A region is simply connected if its complement with respect to the extended plane is connected.
\end{originalstatement}
At this point we warn the reader that this definition is not the one that is commonly accepted, the main reason being that our definition cannot be used in more than two real dimensions. In the course of our work we shall find, however, that the property expressed by Definition~\hyperref[def:4.1]{1} is equivalent to a number of other properties, more or less equally important. One of these states that any closed curve can be contracted to a point, and this condition is usually chosen as definition. Our choice has the advantage of leading very quickly to the essential results in complex integration\index{Integration} theory.

It is easy to see that a disk, a half plane, and a parallel strip are simply connected. The last example shows the importance of taking the complement with respect to the extended plane, for the complement of the strip in the finite plane is evidently not connected. The definition can be applied to regions on the Riemann sphere, and this is evidently the most symmetric situation. For our purposes it is nevertheless better to agree that all regions lie in the finite plane unless the contrary is explicitly stated. According to this convention the outside of a circle is not simply connected, for its complement consists of a closed disk and the point at infinity.

\begin{originalstatement}{Theorem}{14}\label{thm:4.14}
A region $\Omega$ is simply connected if and only if $n(\gamma, a) = 0$ for all cycles $\gamma$ in $\Omega$ and all points a which do not belong to $\Omega$.
\end{originalstatement}
This alternative condition is also very suggestive. It states that a closed curve in a simply connected region\index{Simply connected region|(}\index{Simply connected region} cannot wind around any point which does not belong to the region. It seems quite evident that this condition is not fulfilled in the case of a region with a hole.

The necessity of the condition is almost trivial. Let $\gamma$ be any cycle in $\Omega$. If the complement of $\Omega$ is connected, it must be contained in one of the regions determined by $\gamma$, and inasmuch as $\infty$ belongs to the complement this must be the unbounded region. Consequently $n(\gamma,a)=0$ for all finite points in the complement.

For the precise proof of the sufficiency an explicit construction is needed. We assume that the complement of $\Omega$ can be represented as the union $A \cup B$ of two disjoint closed sets. One of these sets contains $\infty$, and the other is consequently bounded; let A be the bounded set. The sets A and B have a shortest distance $\delta > 0$. Cover the whole plane with a net of squares Q of side $< \delta/\sqrt{2}$. We are free to choose the net so that a certain point $a \in A$ lies at the center of a square. The boundary


% Original PDF physical page 156
\sourcepage{156}
\bookfigure{156}{1}{FIG. 4-9. Curve with index 1.}
curve of Q is denoted by $\partial Q$; we assume explicitly that the squares Q are closed and that the interior of Q lies to the left of the directed line segments which make up $\partial Q$.

Consider now the cycle
\begin{equation*}
\gamma = \sum_ {j} \partial Q _ {j}\tag{39}\label{eq:4.39}
\end{equation*}
where the sum ranges over all squares $Q_{j}$ in the net which have a point in common with A (Fig.~\hyperref[fig:4-9]{4-9}). Because a is contained in one and only one of these squares, it is evident that $n(\gamma,a)=1$. Furthermore, it is clear that $\gamma$ does not meet B. But if the cancellations are carried out, it is equally clear that $\gamma$ does not meet A. Indeed, any side which meets A is a common side of two squares included in the sum \hyperref[eq:4.39]{(39)}, and since the directions are opposite the side does not appear in the reduced expression of $\gamma$. Hence $\gamma$ is contained in $\Omega$, and our theorem is proved.

We remark now that Cauchy's theorem is certainly not valid for regions which are not simply connected. In fact, if there is a cycle $\gamma$ in $\Omega$ such that $n(\gamma,a) \neq 0$ for some a outside of $\Omega$, then $1/(z - a)$ is analytic in


% Original PDF physical page 157
\sourcepage{157}
$\Omega$\index{Homologous, definition|(} while its integral
\begin{equation*}
\int_ {\gamma} \frac {d z}{z - a} = 2 \pi i n (\gamma , a) \neq 0.
\end{equation*}
\setcounter{section}{4}\setcounter{subsection}{2}\subsection{Homology}\label{sec:4.4.3}
The characterization of simple connectivity by Theorem~\hyperref[thm:4.14]{14} singles out a property that is common to all cycles in a simply connected region\index{Simply connected region}, but which a cycle in an arbitrary region or open set may or may not have. This property plays an important role in topology and therefore has a special name.

\begin{originalstatement}{Definition}{2}\label{def:4.2}
A cycle $\gamma$ in an open set $\Omega$ is said to be homologous to zero with respect to $\Omega$ if $n(\gamma, a) = 0$ for all points $a$ in the complement of $\Omega$.
\end{originalstatement}
In symbols we write $\gamma \sim 0$ (mod $\Omega$). When it is clear to what open set we are referring, $\Omega$ need not be mentioned. The notation $\gamma_1 \sim \gamma_2$ shall be equivalent to $\gamma_1 - \gamma_2 \sim 0$. Homologies can be added and subtracted, and $\gamma \sim 0$ (mod $\Omega$) implies $\gamma \sim 0$ (mod $\Omega'$) for all $\Omega' \supset \Omega$.

Again, our terminology does not quite agree with standard usage. It was Emil Artin\index{Artin E@Artin, E.} who discovered that the characterization of homology by vanishing winding numbers ties in precisely with what is needed for the general version of Cauchy's theorem. This idea has led to a remarkable simplification of earlier proofs.

\setcounter{section}{4}\setcounter{subsection}{3}\subsection{The General Statement of Cauchy's Theorem}\label{sec:4.4.4}
The definitive form of Cauchy's theorem is now very easy to state.

\begin{originalstatement}{Theorem}{15}\label{thm:4.15}
If $f(z)$ is analytic in $\Omega$, then
\begin{equation*}
\int_ {\gamma} f (z) d z = 0\tag{40}\label{eq:4.40}
\end{equation*}
for every cycle $\gamma$ which is homologous to zero in $\Omega$.

\end{originalstatement}
In a different formulation, we are claiming that if $\gamma$ is such that \hyperref[eq:4.40]{(40)} holds for a certain collection of analytic functions, namely those of the form $1/(z - a)$ with a not in $\Omega$, then it holds for all analytic functions in $\Omega$.

In combination with Theorem~\hyperref[thm:4.14]{14} we have the following corollary:

\begin{originalstatement}{Corollary}{1}\label{cor:4.1-p157-b12}
If $f(z)$ is analytic in a simply connected region $\Omega$, then \hyperref[eq:4.40]{(40)} holds for all cycles $\gamma$ in $\Omega$.
\end{originalstatement}
Before proving the theorem, we make an observation which ties up with the considerations in Section~\hyperref[sec:4.1.3]{1.3}. As pointed out in that connection,
% Original PDF physical page 158
\index{Homologous, definition|)}\sourcepage{158}
the validity of \hyperref[eq:4.40]{(40)} for all closed curves $\gamma$ in a region means that the line integral of f dz is independent of the path, or that f dz is an exact differential. By Theorem~\hyperref[thm:4.1]{1} there is then a single-valued analytic function $F(z)$ such that $F'(z) = f(z)$ (the pleonastic term “single-valued” is used for emphasis only). In a simply connected region\index{Simply connected region} every analytic function is thus a derivative.

A particular application of this fact occurs very frequently:

\begin{originalstatement}{Corollary}{2}\label{cor:4.2-p158-b2}
If $f(z)$ is analytic and $\neq 0$ in a simply connected region $\Omega$, then it is possible to define single-valued analytic branches of $\log f(z)$ and $\sqrt[n]{f(z)}$ in $\Omega$.
\end{originalstatement}
In fact, we know that there exists an analytic function $F(z)$ in $\Omega$ such that $F'(z) = f'(z)/f(z)$. The function $f(z)e^{-F(z)}$ has the derivative zero and is therefore a constant. Choosing a point $z_{0} \in \Omega$ and one of the infinitely many values $\log f(z_{0})$, we find that
\begin{equation*}
e ^ {F (z) - F (z _ {0}) + \log f (z _ {0})} = f (z),
\end{equation*}
and consequently we can set $\log f(z) = F(z) - F(z_0) + \log f(z_0)$. To define $\sqrt[n]{f(z)}$ we merely write it in the form $\exp((1/n) \log f(z))$.

\setcounter{section}{4}\setcounter{subsection}{4}\subsection{Proof of Cauchy's Theorem}\label{sec:4.4.5}
\hyperlink{fn:158:dagger}{\textsuperscript{\ensuremath{\dagger}}} We begin with a construction that parallels the one in the proof of Theorem~\hyperref[thm:4.14]{14}. Assume first that $\Omega$ is bounded, but otherwise arbitrary. Given $\delta > 0$, we cover the plane by a net of squares of side $\delta$, and we denote by $Q_j, j \in J$, the closed squares in the net which are contained in $\Omega$; because $\Omega$ is bounded the set $J$ is finite, and if $\delta$ is sufficiently small it is not empty. The union of the squares $Q_j, j \in J$, consists of closed regions whose oriented boundaries make up the cycle
\begin{equation*}
\Gamma_ {\delta} = \sum_ {j \in J} \partial Q _ {j}.
\end{equation*}
Clearly, $\Gamma_{\delta}$ is a sum of oriented line segments which are sides of exactly one $Q_{j}$. We denote by $\Omega_{\delta}$ the interior of the union $\cup Q_{j}$ (Fig.~\hyperref[fig:4-10]{4-10}).

Let $\gamma$ be a cycle which is homologous to zero in $\Omega$; we choose $\delta$ so small that $\gamma$ is contained in $\Omega_{\delta}$. Consider a point $\zeta \in \Omega - \Omega_{\delta}$. It belongs to at least one $Q$ which is not a $Q_{j}$. There is a point $\zeta_0 \in Q$ which is not in $\Omega$. It is possible to join $\zeta$ and $\zeta_0$ by a line segment which lies in $Q$ and therefore does not meet $\Omega_{\delta}$. Since $\gamma$, considered as a point set, is contained in $\Omega_{\delta}$ it follows that $n(\gamma, \zeta) = n(\gamma, \zeta_0) = 0$. In particular, $n(\gamma, \zeta) = 0$ for all points $\zeta$ on $\Gamma_{\delta}$.

\begingroup\renewcommand{\thefootnote}{}\footnotetext{\hypertarget{fn:158:dagger}{}\textsuperscript{\ensuremath{\dagger}} This proof follows a suggestion by A. F. Beardon\index{Beardon A F@Beardon, A. F.|hyperpagenote}, who has kindly permitted its use in this connection.}\endgroup

% Original PDF physical page 159
\sourcepage{159}
\bookfigure{159}{1}{FIG. 4-10}
Suppose now that $f$ is analytic in $\Omega$. If $z$ lies in the interior of $Q_{j_0}$, say, then
\begin{equation*}
\frac {1}{2 \pi i} \int_ {\partial Q _ {j}} \frac {f (\zeta) d \zeta}{\zeta - z} = \left\{ \begin{array}{l l} f (z) & \text {if} j = j _ {0} \\ 0 & \text {if} j \neq j _ {0} \end{array} \right.
\end{equation*}
and hence
\begin{equation*}
f (z) = \frac {1}{2 \pi i} \int_ {\Gamma_ {\delta}} \frac {f (\zeta) d \zeta}{\zeta - z}.\tag{41}\label{eq:4.41}
\end{equation*}
Since both sides are continuous functions of $z$, this equation will hold for all $z \in \Omega_{\delta}$.

As a consequence we obtain
\begin{equation*}
\int_ {\gamma} f (z) d z = \int_ {\gamma} \left(\frac {1}{2 \pi i} \int_ {\Gamma_ {\delta}} \frac {f (\zeta) d \zeta}{\zeta - z}\right) d z.\tag{42}\label{eq:4.42}
\end{equation*}
The integrand of the iterated integral is a continuous function of both integration\index{Integration} variables, namely the parameters of $\Gamma_{\delta}$ and $\gamma$. Therefore, the
% Original PDF physical page 160
\sourcepage{160}
order of integration\index{Integration} can be reversed. In other words,
\begin{equation*}
\int_ {\gamma} \left(\frac {1}{2 \pi i} \int_ {\Gamma_ {\delta}} \frac {f (\zeta) d \zeta}{\zeta - z}\right) d z = \int_ {\Gamma_ {\delta}} \left(\frac {1}{2 \pi i} \int_ {\gamma} \frac {d z}{\zeta - z}\right) f (\zeta) d \zeta .
\end{equation*}
On the right the inside integral is $-n(\gamma,\zeta)=0$. Hence the integral \hyperref[eq:4.42]{(42)} is zero, and we have proved the theorem for bounded $\Omega$.

If $\Omega$ is unbounded, we replace it by its intersection $\Omega'$ with a disk $|z| < R$ which is large enough to contain $\gamma$. Any point $a$ in the complement of $\Omega'$ is either in the complement of $\Omega$ or lies outside the disk. In either case $n(\gamma, a) = 0$, so that $\gamma \sim 0 \pmod{\Omega'}$. The proof is applicable to $\Omega'$, and we conclude that the theorem is valid for arbitrary $\Omega$.

\setcounter{section}{4}\setcounter{subsection}{5}\subsection{Locally Exact Differentials}\label{sec:4.4.6}
A differential $p \, dx + q \, dy$ is said to be locally exact in $\Omega$ if it is exact in some neighborhood of each point in $\Omega$. It is not difficult to see (Ex.~\hyperref[exercise:4.4.7.8.1]{1}, p.~\hyperlink{orig-p148}{148}) that this is so if and only if
\begin{equation*}
\int_ {\gamma} p d x + q d y = 0\tag{43}\label{eq:4.43}
\end{equation*}
for all $\gamma = \partial R$ where R is a rectangle contained in $\Omega$. This condition is certainly fulfilled if p dx + q dy = f(z) dz with f analytic in $\Omega$, and by Theorem~\hyperref[thm:4.15]{15}, \hyperref[eq:4.43]{(43)} is then true for any cycle $\gamma \sim 0 \pmod{\Omega}$.

\begin{originalstatement}{Theorem}{16}\label{thm:4.16}
If $p \, dx + q \, dy$ is locally exact in $\Omega$, then
\begin{equation*}
\int_ {\gamma} p d x + q d y = 0
\end{equation*}
for every cycle $\gamma \sim 0$ in $\Omega$.

\end{originalstatement}
There does not seem to be any direct way of modifying the proof of Theorem~\hyperref[thm:4.15]{15} so that it would cover this more general situation. We shall therefore end up by presenting two different proofs of Cauchy's general theorem. As in the earlier editions of this book, we shall follow Artin's proof of Theorem~\hyperref[thm:4.16]{16}. The separate proof of Cauchy's theorem has been included because of its special appeal.

For the proof of Theorem~\hyperref[thm:4.16]{16} we show first that $\gamma$ can be replaced by a polygon $\sigma$ with horizontal and vertical sides such that every locally exact differential\index{Locally exact differential|(}\index{Locally exact differential} has the same integral over $\sigma$ as over $\gamma$. This property implies, in particular, $n(\sigma,a) = n(\gamma,a)$ for a in the complement of $\Omega$, and hence $\sigma \sim 0$. It will thus be sufficient to prove the theorem for polygons with sides parallel to the axes.

We construct $\sigma$ as an approximation of $\gamma$. Let the distance from $\gamma$ to the complement of $\Omega$ be $\rho$. If $\gamma$ is given by $z = z(t)$, the function $z(t)$ is uniformly continuous on the closed interval $[a,b]$. We determine $\delta > 0$ so


% Original PDF physical page 161
\index{Simply connected region|)}\sourcepage{161}
\bookfigure{161}{1}{FIG. 4-11}
that $|z(t) - z(t')| < \rho$ for $|t - t'| < \delta$ and divide $[a, b]$ into subintervals of length $< \delta$. The corresponding aubarcs $\gamma_i$ of $\gamma$ have the property that each is contained in a disk of radius $\rho$ which lies entirely in $\Omega$. The end points of $\gamma_i$ can be joined within that disk by a polygon $\sigma_i$ consisting of a horizontal and a vertical segment. Since the differential is exact in the disk,
\begin{equation*}
\int_ {\sigma_ {i}} p d x + q d y = \int_ {\gamma_ {i}} p d x + q d y,
\end{equation*}
and if $\sigma = \Sigma \sigma_{i}$, we obtain
\begin{equation*}
\int_ {\sigma} p d x + q d y = \int_ {\gamma} p d x + q d y,
\end{equation*}
as desired.

To continue the proof we extend all segments that make up $\sigma$ to infinite lines (Fig.~\hyperref[fig:4-11]{4-11}). They divide the plane into some finite rectangles $R_{i}$ and some unbounded regions $R_{j}^{\prime}$ which may be regarded as infinite rectangles.

Choose a point $a_{i}$ from the interior of each $R_{i}$, and form the cycle
\begin{equation*}
\sigma_ {0} = \sum_ {i} n (\sigma , a _ {i}) \partial R _ {i}\tag{44}\label{eq:4.44}
\end{equation*}
where the sum ranges over all finite rectangles; the coefficients $n(\sigma, a_{i})$ are well determined, for no $a_{i}$ lies on $\sigma$. In the discussion that follows we shall also make use of points $a_{j}^{\prime}$ chosen from the interior of each $R_{j}^{\prime}$.

It is clear that $n(\partial R_i, a_k) = 1$ if $k = i$ and $0$ if $k \neq i$; similarly, $n(\partial R_i, a_j') = 0$ for all $j$. With this in mind it follows from \hyperref[eq:4.44]{(44)} that $n(\sigma_0, a_i) = n(\sigma, a_i)$ and $n(\sigma_0, a_j') = 0$. It is also true that $n(\sigma, a_j') = 0$, for
% Original PDF physical page 162
\sourcepage{162}
the interior of $R_j'$ belongs to the unbounded region determined by $\sigma$. We have thus shown that $n(\sigma - \sigma_0, a) = 0$ for all $a = a_i$ and $a = a_j'$.

From this property of $\sigma - \sigma_0$ we wish to conclude that $\sigma_0$ is identical with $\sigma$ up to segments that cancel against each other. Let $\sigma_{ik}$ be the common side of two adjacent rectangles $R_i, R_k$; we choose the orientation so that $R_i$ lies to the left of $\sigma_{ik}$. Suppose that the reduced expression of $\sigma - \sigma_0$ contains the multiple $c\sigma_{ik}$. Then the cycle $\sigma - \sigma_0 - c \partial R_i$ does not contain $\sigma_{ik}$, and it follows that $a_i$ and $a_k$ must have the same index with respect to this cycle. On the other hand, these indices are $-c$ and 0, respectively; we conclude that $c = 0$. The same reasoning applies if $\sigma_{ij}$ is the common side of a finite rectangle $R_i$ and an infinite rectangle $R_j'$. Thus every side of a finite rectangle occurs with coefficient zero in $\sigma - \sigma_0$, proving that
\begin{equation*}
\sigma = \sum_ {i} n (\sigma , a _ {i}) \partial R _ {i}.\tag{45}\label{eq:4.45}
\end{equation*}
We prove now that all the $R_{i}$ whose corresponding coefficient $n(\sigma, a_{i})$ is different from zero are actually contained in $\Omega$. Suppose that a point a in the closed rectangle $R_{i}$ were not in $\Omega$. Then $n(\sigma, a) = 0$ because $\sigma \sim 0 \pmod{\Omega}$. On the other hand, the line segment between a and $a_{i}$ does not intersect $\sigma$, and hence $n(\sigma, a_{i}) = n(\sigma, a) = 0$. We conclude by the local exactness that the integral of p dx + q dy over any $\partial R_{i}$ which occurs effectively in \hyperref[eq:4.45]{(45)} is zero. Consequently,
\begin{equation*}
\int_ {\sigma} p d x + q d y = 0,
\end{equation*}
and Theorem~\hyperref[thm:4.16]{16} is proved.

\setcounter{section}{4}\setcounter{subsection}{6}\subsection{Multiply Connected Regions}\index{Multiply connected region|(}\index{Multiply connected region}\label{sec:4.4.7}
A region which is not simply connected is called multiply connected. More precisely, $\Omega$ is said to have the finite connectivity\index{Connectivity|(}\index{Connectivity} n if the complement of $\Omega$ has exactly n components and infinite connectivity if the complement has infinitely many components. In a less precise but more suggestive language, a region of connectivity n arises by punching n holes in the Riemann sphere.

In the case of finite connectivity, let $A_1, A_2, \ldots, A_n$ be the components of the complement of $\Omega$, and assume that $\infty$ belongs to $A_n$. If $\gamma$ is an arbitrary cycle in $\Omega$, we can prove, just as in Theorem~\hyperref[thm:4.14]{14}, that $n(\gamma, a)$ is constant when $a$ varies over any one of the components $A_i$ and that $n(\gamma, a) = 0$ in $A_n$. Moreover, duplicating the construction used in the proof of the same theorem we can find cycles $\gamma_i, i = 1, \ldots, n - 1$, such that $n(\gamma_i, a) = 1$ for $a \in A_i$ and $n(\gamma_i, a) = 0$ for all other points outside of $\Omega$.


% Original PDF physical page 163
\index{Locally exact differential|)}\sourcepage{163}
For\index{Module|(} a given cycle $\gamma$ in $\Omega$, let $c_i$ be the constant value of $n(\gamma, a)$ for $a \in A_i$. We find that any point outside of $\Omega$ has the index zero with respect to the cycle $\gamma - c_1 \gamma_1 - c_2 \gamma_2 - \cdots - c_{n-1} \gamma_{n-1}$. In other words,
\begin{equation*}
\gamma \sim c _ {1} \gamma_ {1} + c _ {2} \gamma_ {2} + \dots + c _ {n - 1} \gamma_ {n - 1}.
\end{equation*}
Every cycle is thus homologous to a linear combination of the cycles $\gamma_{1}, \gamma_{2}, \ldots, \gamma_{n-1}$. This linear combination is uniquely determined, for if two linear combinations were homologous to the same cycle their difference would be a linear combination which is homologous to zero. But it is clear that the cycle $c_{1}\gamma_{1} + c_{2}\gamma_{2} + \cdots + c_{n-1}\gamma_{n-1}$ winds $c_{i}$ times around the points in $A_{i}$; hence it cannot be homologous to zero unless all the $c_{i}$ vanish.

In view of these circumstances the cycles $\gamma_{1}, \gamma_{2}, \ldots, \gamma_{n-1}$ are said to form a homology basis\index{Homology basis} for the region $\Omega$. It is not the only homology basis, but by an elementary theorem in linear algebra we may conclude that every homology basis has the same number of elements. We find that every region with a finite homology basis has finite connectivity\index{Connectivity}, and the number of basis elements is one less than the connectivity.

By Theorem~\hyperref[thm:4.18]{18} we obtain, for any analytic function $f(z)$ in $\Omega$,
\begin{equation*}
\int_ {\gamma} f d z = c _ {1} \int_ {\gamma_ {1}} f d z + c _ {2} \int_ {\gamma_ {2}} f d z + \dots + c _ {n - 1} \int_ {\gamma_ {n - 1}} f d z.
\end{equation*}
The numbers
\begin{equation*}
P _ {i} = \int_ {\gamma_ {i}} f d z
\end{equation*}
depend only on the function, and not on $\gamma$. They are called modules of periodicity of the differential f dz, or, with less accuracy, the periods of the indefinite integral. We have found that the integral of $f(z)$ over any cycle is a linear combination of the periods with integers as coefficients, and the integral along an arc from $z_{0}$ to z is determined up to additive multiples of the periods. The vanishing of the periods is a necessary and sufficient condition for the existence of a single-valued indefinite integral.

In order to illustrate, let us consider the extremely simple case of an annulus, defined by $r_{1} < |z| < r_{2}$. The complement has the components $|z| \leq r_{1}$ and $|z| \geq r_{2}$; we include the degenerate cases $r_{1} = 0$ and $r_{2} = \infty$. The annulus is doubly connected, and a homology basis is formed by any circle $|z| = r$, $r_{1} < r < r_{2}$. If this circle is denoted by C, any cycle in the annulus satisfies $\gamma \sim nC$ where $n = n(\gamma, 0)$. The integral of an analytic function over a cycle is a multiple of the single period
\begin{equation*}
P = \int_ {C} f d z
\end{equation*}
% Original PDF physical page 164
\index{Module|)}\sourcepage{164}
whose value is of course independent of the radius r.

\subsection*{EXERCISES}
\originalnumber{1}{exercise:4.4.7.8.1}
1. Prove without use of Theorem~\hyperref[thm:4.16]{16} that $p \, dx + q \, dy$ is locally exact in $\Omega$ if and only if
\begin{equation*}
\int_ {\partial R} p d x + q d y = 0
\end{equation*}
for every rectangle $R \subset \Omega$ with sides parallel to the axes.

\originalnumber{2}{exercise:4.4.7.8.2}
2. Prove that the region obtained from a simply connected region by removing m points has the connectivity\index{Connectivity} $m + 1$, and find a homology basis.

\originalnumber{3}{exercise:4.4.7.8.3}
3. Show that the bounded regions determined by a closed curve are simply connected, while the unbounded region is doubly connected.

\originalnumber{4}{exercise:4.4.7.8.4}
4. Show that single-valued analytic branches of $\log z$, $z^{a}$ and $z^{z}$ can be defined in any simply connected region which does not contain the origin.

\originalnumber{5}{exercise:4.4.7.8.5}
5. Show that a single-valued analytic branch of $\sqrt{1-z^{2}}$ can be defined in any region such that the points $\pm1$ are in the same component of the complement. What are the possible values of
\begin{equation*}
\int \frac {d z}{\sqrt {1 - z ^ {2}}}
\end{equation*}
over a closed curve in the region?

\setcounter{section}{4}\section{The Calculus Of Residues}\label{sec:4.5}
The results of the preceding section have shown that the determination of line integrals of analytic functions over closed curves can be reduced to the determination of periods. Under certain circumstances it turns out that the periods can be found without or with very little computation. We are thus in possession of a method which in many cases permits us to evaluate integrals without resorting to explicit calculation. This is of great value for practical purposes as well as for the further development of the theory.

In order to make this method more systematic a simple formalism, known as the calculus of residues\index{Calculus of residues|(}\index{Calculus of residues}, was introduced by Cauchy\index{Cauchy A@Cauchy, A.}, the founder of complex integration\index{Integration} theory. From the point of view adopted in this book the use of residues amounts essentially to an application of the results proved in Sec. 4 under particularly simple circumstances.

\setcounter{section}{5}\setcounter{subsection}{0}\subsection{The Residue Theorem}\index{Residue theorem|(}\index{Residue|(}\index{Residue theorem}\index{Residue}\label{sec:4.5.1}
Our first task is to review earlier results in the light of the more general theorems of Sec. 4. Clearly, all results which were derived as consequences of Cauchy's theorem for a disk remain valid in arbitrary regions for all cycles which are homologous
% Original PDF physical page 165
\index{Cauchy's integral theorem|)}\index{Connectivity|)}\index{Multiply connected region|)}\sourcepage{165}
to zero. For instance, and this application is typical, Cauchy's integral formula can now be expressed in the following form:

\begin{originalassertion}
If $f(z)$ is analytic in a region $\Omega$, then
\begin{equation*}
n (\gamma , a) f (a) = \frac {1}{2 \pi i} \int_ {\gamma} \frac {f (z)}{z - a} d z
\end{equation*}
for every cycle $\gamma$ which is homologous to zero in $\Omega$.
\end{originalassertion}

The proof is a repetition of the proof of Theorem~\hyperref[thm:4.6]{6}. In this connection we point out that there is of course no longer any need to give a separate proof of Theorem~\hyperref[thm:4.15]{15} in the presence of removable singularities. Indeed, our discussion of the local behavior has already shown that all removable singularities can simply be ignored.

We turn now to the discussion of a function $f(z)$ which is analytic in a region $\Omega$ except for isolated singularities. For a first orientation, let us assume that there are only a finite number of singular points, denoted by $a_{1}, a_{2}, \ldots, a_{n}$. The region obtained by excluding the points $a_{j}$ will be denoted by $\Omega'$.

To each $a_j$ there exists a $\delta_j > 0$ such that the doubly connected region $0 < |z - a_j| < \delta_j$ is contained in $\Omega'$. Draw a circle $C_j$ about $a_j$ of radius $< \delta_j$, and let
\begin{equation*}
P _ {j} = \int_ {C _ {j}} f (z) d z\tag{46}\label{eq:4.46}
\end{equation*}
be the corresponding period of $f(z)$. The particular function $1/(z - a_{j})$ has the period $2\pi i$. Therefore, if we set $R_{j} = P_{j}/2\pi i$, the combination
\begin{equation*}
f (z) - \frac {R _ {j}}{z - a _ {j}}
\end{equation*}
has a vanishing period. The constant $R_{j}$ which produces this result is called the residue\index{Residue} of $f(z)$ at the point $a_{j}$. We repeat the definition in the following form:

\begin{originalstatement}{Definition}{3}\label{def:4.3}
The residue of $f(z)$ at an isolated singularity $a$ is the unique complex number $R$ which makes $f(z) - R / (z - a)$ the derivative of a single-valued analytic function in an annulus $0 < |z - a| < \delta$.
\end{originalstatement}
It is helpful to use such self-explanatory notations as $R = \operatorname{Res}_{z=a} f(z)$

Let $\gamma$ be a cycle in $\Omega'$ which is homologous to zero with respect to $\Omega$. Then $\gamma$ satisfies the homology
\begin{equation*}
\gamma \sim \sum_ {j} n (\gamma , a _ {j}) C _ {j}
\end{equation*}
with respect to $\Omega'$; indeed, we can easily verify that the points $a_{j}$ as well as all points outside of $\Omega$ have the same order with respect to both cycles.


% Original PDF physical page 166
\sourcepage{166}
By virtue of the homology we obtain, with the notation \hyperref[eq:4.46]{(46)},
\begin{equation*}
\int_ {\gamma} f d z = \sum_ {j} n (\gamma , a _ {j}) P _ {j},
\end{equation*}
and since $P_{j} = 2\pi i \cdot R_{j}$ finally
\begin{equation*}
\frac {1}{2 \pi i} \int_ {\gamma} f d z = \sum_ {j} n (\gamma , a _ {j}) R _ {j}.
\end{equation*}
This is the residue\index{Residue} theorem\index{Residue theorem}, except for the restrictive assumption that there are only a finite number of singularities. In the general case we need only prove that $n(\gamma,a_{j})=0$ except for a finite number of points $a_{j}$, for then the same proof can be applied. The assertion follows by routine reasoning. The set of all points a with $n(\gamma,a)=0$ is open and contains all points outside of a large circle. The complement is consequently a compact set, and as such it cannot contain more than a finite number of the isolated points $a_{j}$. Therefore $n(\gamma,a_{j})\neq0$ only for a finite number of the singularities, and we have proved:

\begin{originalstatement}{Theorem}{17}\label{thm:4.17}
Let $f(z)$ be analytic except for isolated singularities $a_j$ in a region $\Omega$. Then
\begin{equation*}
\frac {1}{2 \pi i} \int_ {\gamma} f (z) d z = \sum_ {j} n (\gamma , a _ {j}) \operatorname{Res} _ {z = a _ {j}} f (z)\tag{47}\label{eq:4.47}
\end{equation*}
for any cycle $\gamma$ which is homologous to zero in $\Omega$ and does not pass through any of the points $a_{j}$.

\end{originalstatement}
In the applications it is frequently the case that each $n(\gamma, a_{j})$ is either 0 or 1. Then we have simply
\begin{equation*}
\frac {1}{2 \pi i} \int_ {\gamma} f (z) d z = \sum_ {j} \operatorname{Res} _ {z = a _ {j}} f (z)
\end{equation*}
where the sum is extended over all singularities enclosed by $\gamma$.

The residue theorem is of little value unless we have at our disposal a simple procedure to determine the residues. For essential singularities there is no such procedure of any practical value, and thus it is not surprising that the residue theorem is comparatively seldom used in the presence of essential singularities. With respect to poles the situation is entirely different. We need only look at the expansion
\begin{equation*}
f (z) = B _ {h} (z - a) ^ {- h} + \dots + B _ {1} (z - a) ^ {- 1} + \varphi (z)
\end{equation*}
to recognize that the residue equals the coefficient $B_{1}$. Indeed, when the term $B_{1}(z - a)^{-1}$ is omitted, the remainder is evidently a derivative.


% Original PDF physical page 167
\sourcepage{167}
Since the principal part at a pole is always either given or can be easily found, we have thus a very simple method for finding the residues.

For simple poles the method is even more immediate, for then the residue\index{Residue} equals the value of the function $(z - a)f(z)$ for z = a. For instance, let it be required to find the residues of the function
\begin{equation*}
\frac {e ^ {z}}{(z - a) (z - b)}
\end{equation*}
at the poles $a$ and $b \neq a$. The residue at $a$ is obviously $e^a / (a - b)$, and the residue at $b$ is $e^b / (b - a)$. If $b = a$, the situation is slightly more complicated. We must then expand $e^z$ by Taylor's theorem in the form $e^z = e^a + e^a (z - a) + f_2(z)(z - a)^2$. Dividing by $(z - a)^2$ we find that the residue of $e^z / (z - a)^2$ at $z = a$ is $e^a$.

Remark. In presentations of Cauchy's theorem, the integral formula and the residue theorem\index{Residue theorem} which follow more classical lines, there is no mention of homology, nor is the notion of index used explicitly. Instead, the curve $\gamma$ to which the theorems are applied is supposed to form the complete boundary of a subregion of $\Omega$, and the orientation is chosen so that the subregion lies to the left of $\Omega$. In rigorous texts considerable effort is spent on proving that these intuitive notions have a precise meaning. The main objection to this procedure is the necessity to allot time and attention to rather delicate questions which are peripheral in comparison with the main issues.

With the general point of view that we have adopted it is still possible, and indeed quite easy, to isolate the classical case. All that is needed is to accept the following definition:

\begin{originalstatement}{Definition}{4}\label{def:4.4}
A cycle $\gamma$ is said to bound the region $\Omega$ if and only if $n(\gamma, a)$ is defined and equal to 1 for all points $a \in \Omega$ and either undefined or equal to zero for all points $a$ not in $\Omega$.
\end{originalstatement}
If $\gamma$ bounds $\Omega$, and if $\Omega + \gamma$ is contained in a larger region $\Omega'$, then it is clear that $\gamma$ is homologous to zero with respect to $\Omega'$. The following statements are therefore trivial consequences of Theorems~\hyperref[thm:4.15]{15} and \hyperref[thm:4.17]{17}:

\begin{originalassertion}
If $\gamma$ bounds $\Omega$ and $f(z)$ is analytic on the set $\Omega +\gamma$, then
\begin{equation*}
\int_ {\gamma} f (z) d z = 0
\end{equation*}
and
\begin{equation*}
f (z) = \frac {1}{2 \pi i} \int_ {\gamma} \frac {f (\zeta) d \zeta}{\zeta - z}
\end{equation*}
for all $z \in \Omega$.
\end{originalassertion}


% Original PDF physical page 168
\index{Residue theorem|)}\sourcepage{168}
\begin{originalassertion}
If $f(z)$ is analytic on $\Omega + \gamma$ except for isolated singularities in $\Omega$, then
\begin{equation*}
\frac {1}{2 \pi i} \int_ {\gamma} f (z) d z = \sum_ {j} \operatorname{Res} _ {z = a _ {j}} f (z)
\end{equation*}
where the sum ranges over the singularities $a_{j} \in \Omega$.
\end{originalassertion}

We observe that a cycle $\gamma$ which bounds $\Omega$ must contain the set theoretic boundary of $\Omega$. Indeed, if $z_{0}$ lies on the boundary of $\Omega$, then every neighborhood of $z_{0}$ contains points from $\Omega$ and points not in $\Omega$. If such a neighborhood were free from points of $\gamma$, $n(\gamma,z)$ would be defined and constant in the neighborhood. This contradicts the definition, and hence every neighborhood of $z_{0}$ must meet $\gamma$; since $\gamma$ is closed, $z_{0}$ must lie on $\gamma$.

The converse of the preceding statement is not true, for a point on $\gamma$ may well have a neighborhood which does not meet $\Omega$. Normally, one would try to choose $\gamma$ so that it is identical with the boundary of $\Omega$, but for Cauchy's theorem and related considerations this assumption is not needed.

\setcounter{section}{5}\setcounter{subsection}{1}\subsection{The Argument Principle}\index{Argument principle|(}\index{Argument principle}\label{sec:4.5.2}
Cauchy's integral formula can be considered as a special case of the residue\index{Residue} theorem. Indeed, the function $f(z)/(z-a)$ has a simple pole at z=a with the residue $f(a)$, and when we apply \hyperref[eq:4.47]{(47)}, the integral formula results.

Another application of the residue theorem occurred in the proof of Theorem~\hyperref[thm:4.10]{10} which served to determine the number of zeros of an analytic function. For a zero of order h we can write $f(z) = (z - a)^{h}f_{h}(z)$, with $f_{h}(a) \neq 0$, and obtain $f'(z) = h(z - a)^{h-1}f_{h}(z) + (z - a)^{h}f_{h}'(z)$. Consequently $f'(z)/f(z) = h/(z - a) + f_{h}'(z)/f_{h}(z)$, and we see that $f'/f$ has a simple pole with the residue h. In the formula \hyperref[eq:4.32]{(32)} this residue is accounted for by a corresponding repetition of terms.

We can now generalize Theorem~\hyperref[thm:4.10]{10} to the case of a meromorphic function. If f has a pole of order h, we find by the same calculation as above, with -h replacing h, that $f'/f$ has the residue -h. The following theorem results:

\begin{originalstatement}{Theorem}{18}\label{thm:4.18}
If $f(z)$ is meromorphic in $\Omega$ with the zeros $a_{j}$ and the poles $b_{k}$, then
\begin{equation*}
\frac {1}{2 \pi i} \int_ {\gamma} \frac {f ^ {\prime} (z)}{f (z)} d z = \sum_ {j} n (\gamma , a _ {j}) - \sum_ {k} n (\gamma , b _ {k})\tag{48}\label{eq:4.48}
\end{equation*}
for every cycle $\gamma$ which is homologous to zero in $\Omega$ and does not pass through any of the zeros or poles.

\end{originalstatement}

% Original PDF physical page 169
\sourcepage{169}
It\index{Rouche's theorem|(} is understood that multiple zeros and poles have to be repeated as many times as their order indicates; the sums in \hyperref[eq:4.48]{(48)} are finite.

Theorem~\hyperref[thm:4.18]{18} is usually referred to as the argument principle\index{Argument principle}. The name refers to the interpretation of the left-hand member of \hyperref[eq:4.48]{(48)} as $n(\Gamma,0)$ where $\Gamma$ is the image cycle of $\gamma$. If $\Gamma$ lies in a disk which does not contain the origin, then $n(\Gamma,0) = 0$. This observation is the basis for the following corollary, known as Rouché's theorem:

\begin{originalstatement}{Corollary}{}\label{cor:4.unnumbered-p169-b2}
Let $\gamma$ be homologous to zero in $\Omega$ and such that $n(\gamma, z)$ is either 0 or 1 for any point $z$ not on $\gamma$. Suppose that $f(z)$ and $g(z)$ are analytic in $\Omega$ and satisfy the inequality $|f(z) - g(z)| < |f(z)|$ on $\gamma$. Then $f(z)$ and $g(z)$ have the same number of zeros enclosed by $\gamma$.
\end{originalstatement}
The assumption implies that $f(z)$ and $g(z)$ are zero-free on $\gamma$. Moreover, they satisfy the inequality
\begin{equation*}
\left| \frac {g (z)}{f (z)} - 1 \right| <   1
\end{equation*}
on $\gamma$. The values of $F(z) = g(z)/f(z)$ on $\gamma$ are thus contained in the open disk of center 1 and radius 1. When Theorem~\hyperref[thm:4.18]{18} is applied to $F(z)$, we have thus $n(\Gamma,0) = 0$, and the assertion follows.

A typical application of Rouché's theorem would be the following. Suppose that we wish to find the number of zeros of a function $f(z)$ in the disk $|z| \leq R$. Using Taylor's theorem we can write
\begin{equation*}
f (z) = P _ {n - 1} (z) + z ^ {n} f _ {n} (z)
\end{equation*}
where $P_{n-1}$ is a polynomial of degree $n - 1$. For a suitably chosen $n$ it may happen that we can prove the inequality $R^n |f_n(z)| < |P_{n-1}(z)|$ on $|z| = R$. Then $f(z)$ has the same number of zeros in $|z| \leq R$ as $P_{n-1}(z)$, and this number can be determined by approximate solution of the polynomial equation $P_{n-1}(z) = 0$.

Theorem~\hyperref[thm:4.18]{18} can be generalized in the following manner. If $g(z)$ is analytic in $\Omega$, then $g(z)\frac{f'(z)}{f(z)}$ has the residue\index{Residue} $hg(a)$ at a zero $a$ of order $h$ and the residue $-hg(a)$ at a pole. We obtain thus the formula
\begin{equation*}
\frac {1}{2 \pi i} \int_ {\gamma} g (z) \frac {f ^ {\prime} (z)}{f (z)} d z = \sum_ {j} n (\gamma , a _ {j}) g (a _ {j}) - \sum_ {k} n (\gamma , b _ {k}) g (b _ {k}).\tag{49}\label{eq:4.49}
\end{equation*}
This result is important for the study of the inverse function. With the notations of Theorem~\hyperref[thm:4.11]{11} we know that the equation $f(z) = w$, $|w - w_0| < \delta$ has $n$ roots $z_j(w)$ in the disk $|z - z_0| < \varepsilon$. If we apply


% Original PDF physical page 170
\index{Rouche's theorem|)}\sourcepage{170}
\hyperref[eq:4.49]{(49)} with $g(z) = z$, we obtain
\begin{equation*}
\sum_ {j = 1} ^ {n} z _ {j} (w) = \frac {1}{2 \pi i} \int_ {| z - z _ {0} | = \epsilon} \frac {f ^ {\prime} (z)}{f (z) - w} z d z.\tag{50}\label{eq:4.50}
\end{equation*}
For $n = 1$ the inverse function $f^{-1}(w)$ can thus be represented explicitly by
\begin{equation*}
f ^ {- 1} (w) = \frac {1}{2 \pi i} \int_ {| z - z _ {0} | = \epsilon} \frac {f ^ {\prime} (z)}{f (z) - w} z d z.
\end{equation*}
If \hyperref[eq:4.49]{(49)} is applied with $g(z) = z^m$, equation \hyperref[eq:4.50]{(50)} is replaced by
\begin{equation*}
\sum_ {j = 1} ^ {n} z _ {j} (w) ^ {m} = \frac {1}{2 \pi i} \int_ {| z - z _ {0} | = \epsilon} \frac {f ^ {\prime} (z)}{f (z) - w} z ^ {m} d z.
\end{equation*}
The right-hand member represents an analytic function of w for $|w - w_{0}| < \delta$. Thus the power sums of the roots $z_{j}(w)$ are single-valued analytic functions of w. But it is well known that the elementary symmetric functions can be expressed as polynomials in the power sums. Hence they are also analytic, and we find that the $z_{j}(w)$ are the roots of a polynomial equation
\begin{equation*}
z ^ {n} + a _ {1} (w) z ^ {n - 1} + \dots + a _ {n - 1} (w) z + a _ {n} (w) = 0
\end{equation*}
whose coefficients are analytic functions of w in $|w - w_{0}| < \delta$.

\subsection*{EXERCISES}
\originalnumber{1}{exercise:4.5.2.9.1}
1. How many roots does the equation $z^7 - 2z^5 + 6z^3 - z + 1 = 0$ have in the disk $|z| < 1$? Hint: Look for the biggest term when $|z| = 1$ and apply Rouché's theorem.

\originalnumber{2}{exercise:4.5.2.9.2}
2. How many roots of the equation $z^{4}-6z+3=0$ have their modulus between 1 and 2?

\originalnumber{3}{exercise:4.5.2.9.3}
3. How many roots of the equation $z^{4} + 8z^{3} + 3z^{2} + 8z + 3 = 0$ lie in the right half plane? Hint: Sketch the image of the imaginary axis and apply the argument principle\index{Argument principle} to a large half disk.

\setcounter{section}{5}\setcounter{subsection}{2}\subsection{Evaluation of Definite Integrals}\label{sec:4.5.3}
The calculus of residues\index{Calculus of residues} provides a very efficient tool for the evaluation of definite integrals. It is particularly important when it is impossible to find the indefinite integral explicitly, but even if the ordinary methods of calculus can be applied the use of residues is frequently a laborsaving device. The fact that the calculus of residues yields complex rather than real integrals is no disadvantage, for clearly the evaluation of a complex integral is equivalent to the evaluation of two definite integrals.


% Original PDF physical page 171
\index{Argument principle|)}\sourcepage{171}
There are, however, some serious limitations, and the method is far from infallible. In the first place, the integrand must be closely connected with some analytic function. This is not very serious, for usually we are only required to integrate elementary functions, and they can all be extended to the complex domain. It is much more serious that the technique of complex integration\index{Integration} applies only to closed curves, while a real integral is always extended over an interval. A special device must be used in order to reduce the problem to one which concerns integration over a closed curve. There are a number of ways in which this can be accomplished, but they all apply under rather special circumstances. The technique can be learned at the hand of typical examples, but even complete mastery does not guarantee success.

1. All integrals of the form
\begin{equation*}
\int_ {0} ^ {2 \pi} R (\cos \theta , \sin \theta) d \theta\tag{51}\label{eq:4.51}
\end{equation*}
where the integrand is a rational function of $\cos\theta$ and $\sin\theta$ can be easily evaluated by means of residues. Of course these integrals can also be computed by explicit integration, but this technique is usually very laborious. It is very natural to make the substitution $z=e^{i\theta}$ which immediately transforms \hyperref[eq:4.51]{(51)} into the line integral
\begin{equation*}
- i \int_ {| z | = 1} R \left[ \frac {1}{2} \left(z + \frac {1}{z}\right), \frac {1}{2 i} \left(z - \frac {1}{z}\right) \right] \frac {d z}{z}.
\end{equation*}
It remains only to determine the residues which correspond to the poles of the integrand inside the unit circle.

As an example, let us compute
\begin{equation*}
\int_ {0} ^ {\pi} \frac {d \theta}{a + \cos \theta}, \quad a > 1.
\end{equation*}
This integral is not extended over $(0,2\pi)$, but since $\cos\theta$ takes the same values in the intervals $(0,\pi)$ and $(\pi,2\pi)$ it is clear that the integral from 0 to $\pi$ is one-half of the integral from 0 to $2\pi$. Taking this into account we find that the integral equals
\begin{equation*}
- i \int_ {| z | = 1} \frac {d z}{z ^ {2} + 2 a z + 1}
\end{equation*}
The denominator can be factored into $(z - \alpha)(z - \beta)$ with
\begin{equation*}
\alpha = - a + \sqrt {a ^ {2} - 1}, \quad \beta = - a - \sqrt {a ^ {2} - 1}.
\end{equation*}
Evidently $|\alpha| < 1$, $|\beta| > 1$, and the residue\index{Residue} at $\alpha$ is $1 / (\alpha - \beta)$. The value


% Original PDF physical page 172
\sourcepage{172}
of the integral is found to be $\pi/\sqrt{a^{2}-1}$.

2. An integral of the form
\begin{equation*}
\int_ {- \infty} ^ {\infty} R (x) d x
\end{equation*}
converges if and only if in the rational function $R(x)$ the degree of the denominator is at least two units higher than the degree of the numerator, and if no pole lies on the real axis. The standard procedure is to integrate the complex function $R(z)$ over a closed curve consisting of a line segment $(- \rho, \rho)$ and the semicircle from $\rho$ to $-\rho$ in the upper half plane. If $\rho$ is large enough this curve encloses all poles in the upper half plane, and the corresponding integral is equal to $2\pi i$ times the sum of the residues in the upper half plane. As $\rho \to \infty$ obvious estimates show that the integral over the semicircle tends to 0, and we obtain
\begin{equation*}
\int_ {- \infty} ^ {\infty} R (x) d x = 2 \pi i \sum_ {y > 0} \operatorname{Res} R (z).
\end{equation*}
3. The same method can be applied to an integral of the form
\begin{equation*}
\int_ {- \infty} ^ {\infty} R (x) e ^ {i x} d x\tag{52}\label{eq:4.52}
\end{equation*}
whose real and imaginary parts determine the important integrals
\begin{equation*}
\int_ {- \infty} ^ {\infty} R (x) \cos x d x, \quad \int_ {- \infty} ^ {\infty} R (x) \sin x d x.\tag{53}\label{eq:4.53}
\end{equation*}
Since $|e^{iz}| = e^{-y}$ is bounded in the upper half plane, we can again conclude that the integral over the semicircle tends to zero, provided that the rational function $R(z)$ has a zero of at least order two at infinity. We obtain
\begin{equation*}
\int_ {- \infty} ^ {\infty} R (x) e ^ {i x} d x = 2 \pi i \sum_ {y > 0} \operatorname{Res} R (z) e ^ {i z}.
\end{equation*}
It is less obvious that the same result holds when $R(z)$ has only a simple zero at infinity. In this case it is not convenient to use semicircles. For one thing, it is not so easy to estimate the integral over the semicircle, and secondly, even if we were successful we would only have proved that the integral
\begin{equation*}
\int_ {- \rho} ^ {\rho} R (x) e ^ {i x} d x
\end{equation*}
over a symmetric interval has the desired limit for $\rho\to\infty$. In reality we are of course required to prove that
\begin{equation*}
\int_ {- X _ {1}} ^ {X _ {2}} R (x) e ^ {i x} d x
\end{equation*}
% Original PDF physical page 173
\sourcepage{173}
has a limit when $X_{1}$ and $X_{2}$ tend independently to $\infty$. In the earlier examples this question did not arise because the convergence of the integral was assured beforehand.

For the proof we integrate over the perimeter of a rectangle with the vertices $X_{2}$, $X_{2} + iY$, $-X_{1} + iY$, $-X_{1}$ where $Y > 0$. As soon as $X_{1}$, $X_{2}$ and $Y$ are sufficiently large, this rectangle contains all the poles in the upper half plane. Under the hypothesis $|zR(z)|$ is bounded. Hence the integral over the right vertical side is, except for a constant factor,

less than
\begin{equation*}
\int_ {0} ^ {Y} e ^ {- \nu} \frac {d y}{| z |} <   \frac {1}{X _ {2}} \int_ {0} ^ {Y} e ^ {- \nu} d y.
\end{equation*}
The last integral can be evaluated explicitly and is found to be < 1. Hence the integral over the right vertical side is less than a constant times $1/X_{2}$, and a corresponding result is found for the left vertical side. The integral over the upper horizontal side is evidently less than $e^{-Y}(X_{1} + X_{2})/Y$ multiplied with a constant. For fixed $X_{1}, X_{2}$ it tends to 0 as $Y \to \infty$, and we conclude that
\begin{equation*}
\left| \int_ {- X _ {1}} ^ {X _ {2}} R (x) e ^ {i x} d x - 2 \pi i \sum_ {y > 0} \operatorname{Res} R (z) e ^ {i z} \right| <   A \left(\frac {1}{X _ {1}} + \frac {1}{X _ {2}}\right)
\end{equation*}
where A denotes a constant. This inequality proves that
\begin{equation*}
\int_ {- \infty} ^ {\infty} R (x) e ^ {i x} d x = 2 \pi i \sum_ {y > 0} \operatorname{Res} R (z) e ^ {i z}
\end{equation*}
under the sole condition that $R(\infty) = 0$.

In the discussion we have assumed, tacitly, that $R(z)$ has no poles on the real axis since otherwise the integral \hyperref[eq:4.52]{(52)} has no meaning. However, one of the integrals \hyperref[eq:4.53]{(53)} may well exist, namely, if $R(z)$ has simple poles which coincide with zeros of $\sin x$ or $\cos x$. Let us suppose, for instance, that $R(z)$ has a simple pole at $z = 0$. Then the second integral \hyperref[eq:4.53]{(53)} has a meaning and calls for evaluation.

We use the same method as before, but we use a path which avoids the origin by following a small semicircle of radius $\delta$ in the lower half plane (Fig.~\hyperref[fig:4-12]{4-12}). It is easy to see that this closed curve encloses the poles in the upper half plane, the pole at the origin, and no others, as soon as $X_{1}, X_{2}, Y$ are sufficiently large and $\delta$ is sufficiently small. Suppose that the residue\index{Residue} at 0 is B, so that we can write $R(z)e^{iz} = B/z + R_{0}(z)$ where $R_{0}(z)$ is analytic at the origin. The integral of the first term over the semicircle is $\pi i B$, while the integral of the second term tends to 0 with $\delta$.


% Original PDF physical page 174
\sourcepage{174}
\bookfigure{174}{1}{FIG. 4-12}
It is clear that we are led to the result
\begin{equation*}
\lim _ {\delta \rightarrow 0} \int_ {- \infty} ^ {- \delta} + \int_ {\delta} ^ {\infty} R (x) e ^ {i x} d x = 2 \pi i \left[ \sum_ {y > 0} \operatorname{Res} R (z) e ^ {i z} + \frac {1}{2} B \right] \cdot
\end{equation*}
The limit on the left is called the Cauchy principal value\index{Cauchy principal value} of the integral; it exists although the integral itself has no meaning. On the right-hand side we observe that one-half of the residue\index{Residue} at 0 has been included; this is as if one-half of the pole were counted as lying in the upper half plane.

In the general case where several poles lie on the real axis we obtain
\begin{equation*}
\operatorname{pr.v.} \int_ {- \infty} ^ {\infty} R (x) e ^ {i x} d x = 2 \pi i \sum_ {y > 0} \operatorname{Res} R (z) e ^ {i z} + \pi i \sum_ {y = 0} \operatorname{Res} R (z) e ^ {i z}
\end{equation*}
where the notations are self-explanatory. It is an essential hypothesis that all the poles on the real axis be simple, and as before we must assume that $R(\infty) = 0$.

As the simplest example we have
\begin{equation*}
\operatorname{pr.v.} \int_ {- \infty} ^ {\infty} \frac {e ^ {i x}}{x} d x = \pi i.
\end{equation*}
Separating the real and imaginary part we observe that the real part of the equation is trivial by the fact that the integrand is odd. In the imaginary part it is not necessary to take the principal value, and since the integrand is even we find
\begin{equation*}
\int_ {0} ^ {\infty} \frac {\sin x}{x} d x = \frac {\pi}{2}.
\end{equation*}
We remark that integrals containing a factor $\cos^{n}x$ or $\sin^{n}x$ can be evaluated by the same technique. Indeed, these factors can be written as linear combinations of terms $\cos mx$ and $\sin mx$, and the corresponding integrals can be reduced to the form \hyperref[eq:4.52]{(52)} by a change of variable:


% Original PDF physical page 175
\sourcepage{175}
\begin{equation*}
\int_ {- \infty} ^ {\infty} R (x) e ^ {i m x} d x = \frac {1}{m} \int_ {- \infty} ^ {\infty} R \left(\frac {x}{m}\right) e ^ {i x} d x.
\end{equation*}
4. The next category of integrals have the form
\begin{equation*}
\int_ {0} ^ {\infty} x ^ {\alpha} R (x) d x
\end{equation*}
where the exponent $\alpha$ is real and may be supposed to lie in the interval $0 < \alpha < 1$. For convergence $R(z)$ must have a zero of at least order two at $\infty$ and at most a simple pole at the origin.

The new feature is the fact that $R(z)z^{\alpha}$ is not single-valued. This, however, is just the circumstance which makes it possible to find the integral from 0 to $\infty$.

The simplest procedure is to start with the substitution $x = t^{2}$ which transforms the integral into
\begin{equation*}
2 \int_ {0} ^ {\infty} t ^ {2 \alpha + 1} R (t ^ {2}) d t.
\end{equation*}
For the function $z^{2\alpha}$ we may choose the branch whose argument lies between $-\pi\alpha$ and $3\pi\alpha$; it is well defined and analytic in the region obtained by omitting the negative imaginary axis. As long as we avoid the negative imaginary axis, we can apply the residue\index{Residue} theorem to the function $z^{2\alpha+1}R(z^{2})$. We use a closed curve consisting of two line segments along the positive and negative axis and two semicircles in the upper half plane, one very large and one very small (Fig.~\hyperref[fig:4-13]{4-13}). Under our assumptions it is easy to show that the integrals over the semicircles tend to zero. Hence the residue theorem yields the value of the integral
\begin{equation*}
\int_ {- \infty} ^ {\infty} z ^ {2 \alpha + 1} R (z ^ {2}) d z = \int_ {0} ^ {\infty} (z ^ {2 \alpha + 1} + (- z) ^ {2 \alpha + 1}) R (z ^ {2}) d z.
\end{equation*}
However, $(-z)^{2\alpha} = e^{2\pi i\alpha}z^{2\alpha}$, and the integral equals
\begin{equation*}
(1 - e ^ {2 \pi i \alpha}) \int_ {0} ^ {\infty} z ^ {2 \alpha + 1} R (z ^ {2}) d z.
\end{equation*}
Since the factor in front is $\neq0$, we are finally able to determine the value of the desired integral.

The evaluation calls for determination of the residues of $z^{2\alpha+1}R(z^{2})$ in the upper half plane. These are the same as the residues of $z^{\alpha}R(z)$ in the whole plane. For practical purposes it may be preferable not to use any preliminary substitution and integrate the function $z^{\alpha}R(z)$ over the closed curve shown in Fig.~\hyperref[fig:4-14]{4-14}. We have then to use the branch of $z^{\alpha}$ whose argument lies between 0 and $2\pi\alpha$. This method needs some justification, for it does not conform to the hypotheses of the residue theorem. The justification is trivial.


% Original PDF physical page 176
\sourcepage{176}
\bookfigure{176}{1}{}
5. As a final example we compute the special integral

$\int_0^\pi \log \sin \theta d\theta .$

Consider the function $1 - e^{2iz} = -2ie^{iz} \sin z$. From the representation $1 - e^{2iz} = 1 - e^{-2y} (\cos 2x + i \sin 2x)$, we find that this function is real and negative only for $x = n\pi, y \leq 0$. In the region obtained by omitting these half lines the principal branch of $\log (1 - e^{2iz})$ is hence single-valued and analytic. We apply Cauchy's theorem to the rectangle whose vertices are $0, \pi, \pi + iY$, and $iY$; however, the points 0 and $\pi$ have to be avoided, and we do this by following small circular quadrants of radius $\delta$.

Because of the periodicity the integrals over the vertical sides cancel against each other. The integral over the upper horizontal side tends to 0 as $Y \to \infty$. Finally, the integrals over the quadrants can also be seen to approach zero as $\delta \to 0$. Indeed, since the imaginary part of the logarithm is bounded we need only consider the real part. From the fact that $|1 - e^{2iz}| / |z| \to 2$ for $z \to 0$ we see that $\log |1 - e^{2iz}|$ becomes infinite like $\log \delta$, and since $\delta \log \delta \to 0$ the integral over the quadrant near the origin will tend to zero.

The same proof applies near the vertex $\pi$, and we obtain
\begin{equation*}
\int_ {0} ^ {\pi} \log (- 2 i e ^ {i x} \sin x) d x = 0
\end{equation*}
If we choose $\log e^{ix} = ix$, the imaginary part lies between 0 and $\pi$. Therefore, in order to obtain the principal branch with an imaginary part between $-\pi$ and $\pi$, we must choose $\log (-i) = -\pi i/2$. The equation can now be written in the form


% Original PDF physical page 177
\sourcepage{177}
\begin{equation*}
\pi \log 2 - \left(\frac {\pi^ {2}}{2}\right) i + \int_ {0} ^ {\pi} \log \sin x d x + \left(\frac {\pi^ {2}}{2}\right) i = 0,
\end{equation*}
and we find
\begin{equation*}
\int_ {0} ^ {\pi} \log \sin x d x = - \pi \log 2.
\end{equation*}
\subsection*{EXERCISES}
\originalnumber{1}{exercise:4.5.3.10.1}
1. Find the poles and residues of the following functions:
\begin{equation*}
\begin{aligned}
\text{(a)}\;&\frac{1}{z^{2}+5z+6}, & \text{(b)}\;&\frac{1}{(z^{2}-1)^{2}},\\
\text{(c)}\;&\frac{1}{\sin z}, & \text{(d)}\;&\cot z,\\
\text{(e)}\;&\frac{1}{\sin^{2}z}, & \text{(f)}\;&\frac{1}{z^{m}(1-z)^{n}}.
\end{aligned}
\end{equation*}
($m,n$ positive integers).

\originalnumber{2}{exercise:4.5.3.10.2}
2. Show that in Sec.~\hyperref[sec:4.5.3]{5.3}, Example 3, the integral may be extended over a right-angled isosceles triangle. (Suggested by a student.)

\originalnumber{3}{exercise:4.5.3.10.3}
3. Evaluate the following integrals by the method of residues:
\begin{align*}
\text{(a)}\quad &\int_0^{\pi/2}\frac{dx}{a+\sin^{2}x},\quad |a|>1,\\
\text{(b)}\quad &\int_0^{\infty}\frac{x^{2}\,dx}{x^{4}+5x^{2}+6},\\
\text{(c)}\quad &\int_{-\infty}^{\infty}\frac{x^{2}-x+2}{x^{4}+10x^{2}+9}\,dx,\\
\text{(d)}\quad &\int_0^{\infty}\frac{x^{2}\,dx}{(x^{2}+a^{2})^{3}},\quad a\text{ real},\\
\text{(e)}\quad &\int_0^{\infty}\frac{\cos x}{x^{2}+a^{2}}\,dx,\quad a\text{ real},\\
\text{(f)}\quad &\int_0^{\infty}\frac{x\sin x}{x^{2}+a^{2}}\,dx,\quad a\text{ real},\\
\text{(g)}\quad &\int_0^{\infty}\frac{x^{1/3}}{1+x^{2}}\,dx,\\
\text{(h)}\quad &\int_0^{\infty}(1+x^{2})^{-1}\log x\,dx,\\
\text{(i)}\quad &\int_0^{\infty}\log(1+x^{2})\frac{dx}{x^{1+\alpha}},\quad 0<\alpha<2.
\end{align*}
(Try integration\index{Integration} by parts.)

\originalnumber{4}{exercise:4.5.3.10.4}
4. Compute
\begin{equation*}
\int_ {| z | = \rho} \frac {| d z |}{| z - a | ^ {2}}, \quad | a | \neq \rho .
\end{equation*}
Hint: Use $z\bar{z} = \rho^{2}$ to convert the integral to a line integral of a rational function.

\originalnumber{5}{exercise:4.5.3.10.5}
*5. Complex integration can sometimes be used to evaluate area integrals. As an illustration, show that if $f(z)$ is analytic and bounded for $|z| < 1$ and if $|\zeta| < 1$, then
\begin{equation*}
f (\zeta) = \frac {1}{\pi} \iint_ {| z | <   1} \frac {f (z) d x d y}{(1 - \bar {z} \zeta) ^ {2}}.
\end{equation*}
Remark. This is known as Bergman\index{Bergman S@Bergman, S.}'s kernel\index{Kernel} formula. To prove it, express the area integral in polar coordinates, then transform the inside integral to a line integral which can be evaluated by residues.


% Original PDF physical page 178
\index{Calculus of residues|)}\index{Residue|)}\sourcepage{178}
\setcounter{section}{5}\section{Harmonic Functions}\label{sec:4.6}
The real and imaginary parts of an analytic function are conjugate harmonic functions\index{Function!harmonic|(}\index{Harmonic function|(}\index{Function!harmonic}\index{Harmonic function}. Therefore, all theorems on analytic functions are also theorems on pairs of conjugate harmonic functions. However, harmonic functions are important in their own right, and their treatment is not always simplified by the use of complex methods. This is particularly true when the conjugate harmonic function is not single-valued.

We assemble in this section some facts about harmonic functions that are intimately connected with Cauchy's theorem. The more delicate properties of harmonic functions are postponed to a later chapter.

\setcounter{section}{6}\setcounter{subsection}{0}\subsection{Definition and Basic Properties}\label{sec:4.6.1}
A real-valued function $u(z)$ or $u(x,y)$, defined and single-valued in a region $\Omega$, is said to be harmonic in $\Omega$, or a potential function, if it is continuous together with its partial derivatives of the first two orders and satisfies Laplace's equation\index{Laplace's equation}
\begin{equation*}
\Delta u = \frac {\partial^ {2} u}{\partial x ^ {2}} + \frac {\partial^ {2} u}{\partial y ^ {2}} = 0.\tag{54}\label{eq:4.54}
\end{equation*}
We shall see later that the regularity conditions can be weakened, but this is a point of relatively minor importance.

The sum of two harmonic functions and a constant multiple of a harmonic function are again harmonic; this is due to the linear character of Laplace's equation. The simplest harmonic functions are the linear functions $ax + by$. In polar coordinates $(r, \theta)$ equation \hyperref[eq:4.54]{(54)} takes the form
\begin{equation*}
r \frac {\partial}{\partial r} \left(r \frac {\partial u}{\partial r}\right) + \frac {\partial^ {2} u}{\partial \theta^ {2}} = 0. \text{\hyperlink{fn:178:dagger}{\textsuperscript{\ensuremath{\dagger}}}}
\end{equation*}
This shows that $\log r$ is a harmonic function and that any harmonic function which depends only on r must be of the form $a \log r + b$. The argument $\theta$ is harmonic whenever it can be uniquely defined.

If $u$ is harmonic in $\Omega$, then
\begin{equation*}
f (z) = \frac {\partial u}{\partial x} - i \frac {\partial u}{\partial y}\tag{55}\label{eq:4.55}
\end{equation*}
is analytic, for writing $U = \frac{\partial u}{\partial x}$, $V = -\frac{\partial u}{\partial y}$ we have
\begin{equation*}
\frac {\partial U}{\partial x} = \frac {\partial^ {2} u}{\partial x ^ {2}} = - \frac {\partial^ {2} u}{\partial y ^ {2}} = \frac {\partial V}{\partial y}
\end{equation*}
\begin{equation*}
\frac {\partial U}{\partial y} = \frac {\partial^ {2} u}{\partial x \partial y} = - \frac {\partial V}{\partial x}.
\end{equation*}
\begingroup\renewcommand{\thefootnote}{}\footnotetext{\hypertarget{fn:178:dagger}{}\textsuperscript{\ensuremath{\dagger}} This form cannot be used for $r = C$.}\endgroup

% Original PDF physical page 179
\sourcepage{179}
This, it should be remembered, is the most natural way of passing from harmonic to analytic functions.

From \hyperref[eq:4.55]{(55)} we pass to the differential
\begin{equation*}
f d z = \left(\frac {\partial u}{\partial x} d x + \frac {\partial u}{\partial y} d y\right) + i \left(- \frac {\partial u}{\partial y} d x + \frac {\partial u}{\partial x} d y\right).\tag{56}\label{eq:4.56}
\end{equation*}
In this expression the real part is the differential of u,
\begin{equation*}
d u = \frac {\partial u}{\partial x} d x + \frac {\partial u}{\partial y} d y.
\end{equation*}
If $u$ has a conjugate harmonic function\index{Function!harmonic}\index{Harmonic function} $v$, then the imaginary part can be written as
\begin{equation*}
d v = \frac {\partial v}{\partial x} d x + \frac {\partial v}{\partial y} d y = - \frac {\partial u}{\partial y} d x + \frac {\partial u}{\partial x} d y.
\end{equation*}
In general, however, there is no single-valued conjugate function, and in these circumstances it is better not to use the notation dv. Instead we write
\begin{equation*}
^ {*} d u = - \frac {\partial u}{\partial y} d x + \frac {\partial u}{\partial x} d y
\end{equation*}
and call *du the conjugate differential\index{Conjugate differential} of du. We have by \hyperref[eq:4.56]{(56)}
\begin{equation*}
f d z = d u + i ^ {*} d u.\tag{57}\label{eq:4.57}
\end{equation*}
By Cauchy's theorem the integral of $f \, dz$ vanishes along any cycle which is homologous to zero in $\Omega$. On the other hand, the integral of the exact differential $du$ vanishes along all cycles. It follows by \hyperref[eq:4.57]{(57)} that
\begin{equation*}
\int_ {\gamma} ^ {*} d u = \int_ {\gamma} - \frac {\partial u}{\partial y} d x + \frac {\partial u}{\partial x} d y = 0\tag{58}\label{eq:4.58}
\end{equation*}
for all cycles $\gamma$ which are homologous to zero in $\Omega$.

The integral in \hyperref[eq:4.58]{(58)} has an important interpretation which cannot be left unmentioned. If $\gamma$ is a regular curve with the equation $z = z(t)$, the direction of the tangent is determined by the angle $\alpha = \arg z'(t)$, and we can write $dx = |dz| \cos \alpha$, $dy = |dz| \sin \alpha$. The normal which points to the right of the tangent has the direction $\beta = \alpha - \pi/2$, and thus $\cos \alpha = -\sin \beta$, $\sin \alpha = \cos \beta$. The expression
\begin{equation*}
\frac {\partial u}{\partial n} = \frac {\partial u}{\partial x} \cos \beta + \frac {\partial u}{\partial y} \sin \beta
\end{equation*}
is a directional derivative of u, the right-hand normal derivative\index{Normal derivative} with respect to the curve $\gamma$. We obtain $*du = (\partial u/\partial n) |dz|$, and \hyperref[eq:4.58]{(58)} can be


% Original PDF physical page 180
\sourcepage{180}
written in the form
\begin{equation*}
\int_ {\gamma} \frac {\partial u}{\partial n} | d z | = 0.\tag{59}\label{eq:4.59}
\end{equation*}
This is the classical notation. Its main advantage is that $\partial u/\partial n$ actually represents a rate of change in the direction perpendicular to $\gamma$. For instance, if $\gamma$ is the circle $|z| = r$, described in the positive sense, $\partial u/\partial n$ can be replaced by the partial derivative $\partial u/\partial r$. It has the disadvantage that \hyperref[eq:4.59]{(59)} is not expressed as an ordinary line integral, but as an integral with respect to arc length. For this reason the classical notation is less natural in connection with homology theory, and we prefer to use the notation $*du$.

In a simply connected region the integral of $*du$ vanishes over all cycles, and u has a single-valued conjugate function v which is determined up to an additive constant. In the multiply connected case the conjugate function has periods
\begin{equation*}
\int_ {\gamma^ {i}} ^ {*} d u = \int_ {\gamma^ {i}} \frac {\partial u}{\partial n} | d z |
\end{equation*}
corresponding to the cycles in a homology basis.

There is an important generalization of \hyperref[eq:4.58]{(58)} which deals with a pair of harmonic functions\index{Function!harmonic}\index{Harmonic function}. If $u_{1}$ and $u_{2}$ are harmonic in $\Omega$, we claim that
\begin{equation*}
\int_ {\gamma} u _ {1} ^ {*} d u _ {2} - u _ {2} ^ {*} d u _ {1} = 0\tag{60}\label{eq:4.60}
\end{equation*}
for every cycle $\gamma$ which is homologous to zero in $\Omega$. According to Theorem~\hyperref[thm:4.16]{16}, Sec.~\hyperref[sec:4.4.6]{4.6}, it is sufficient to prove \hyperref[eq:4.60]{(60)} for $\gamma = \partial R$, where R is a rectangle contained in $\Omega$. In R, $u_{1}$ and $u_{2}$ have single-valued conjugate functions $v_{1}, v_{2}$ and we can write
\begin{equation*}
u _ {1} ^ {*} d u _ {2} - u _ {2} ^ {*} d u _ {1} = u _ {1} d v _ {2} - u _ {2} d v _ {1} = u _ {1} d v _ {2} + v _ {1} d u _ {2} - d \left(u _ {2} v _ {1}\right).
\end{equation*}
Here $d(u_{2}v_{1})$ is an exact differential, and $u_{1}dv_{2} + v_{1}du_{2}$ is the imaginary part of
\begin{equation*}
(u _ {1} + i v _ {1}) (d u _ {2} + i d v _ {2}).
\end{equation*}
The last differential can be written in the form $F_{1}f_{2} dz$ where $F_{1}(z)$ and $f_{2}(z)$ are analytic on R. The integral of $F_{1}f_{2} dz$ vanishes by Cauchy's theorem, and so does therefore the integral of its imaginary part. We conclude that \hyperref[eq:4.60]{(60)} holds for $\gamma = \partial R$, and we have proved:

\begin{originalstatement}{Theorem}{19}\label{thm:4.19}
If $u_{1}$ and $u_{2}$ are harmonic in a region $\Omega$, then
\begin{equation*}
\int_ {\gamma} u _ {1} ^ {*} d u _ {2} - u _ {2} ^ {*} d u _ {1} = 0\tag{60}\label{eq:4.60-p180-b14}
\end{equation*}
% Original PDF physical page 181
\sourcepage{181}
for every cycle $\gamma$ which is homologous to zero in $\Omega$.

\end{originalstatement}
For $u_{1} = 1$, $u_{2} = u$ the formula reduces to \hyperref[eq:4.58]{(58)}. In the classical notation \hyperref[eq:4.60]{(60)} would be written as
\begin{equation*}
\int_ {\gamma} \left(u _ {1} \frac {\partial u _ {2}}{\partial n} - u _ {2} \frac {\partial u _ {1}}{\partial n}\right) | d z | = 0.
\end{equation*}
\setcounter{section}{6}\setcounter{subsection}{1}\subsection{The Mean-value Property}\index{Mean-value property|(}\index{Mean-value property}\label{sec:4.6.2}
Let us apply Theorem~\hyperref[thm:4.19]{19} with $u_{1} = \log r$ and $u_{2}$ equal to a function $u$, harmonic in $|z| < \rho$. For $\Omega$ we choose the punctured disk $0 < |z| < \rho$, and for $\gamma$ we take the cycle $C_{1} - C_{2}$ where $C_{i}$ is a circle $|z| = r_{i} < \rho$ described in the positive sense. On a circle $|z| = r$ we have $^{*}du = r(\partial u / \partial r) d\theta$ and hence \hyperref[eq:4.60]{(60)} yields
\begin{equation*}
\log r _ {1} \int_ {C _ {1}} r _ {1} \frac {\partial u}{\partial r} d \theta - \int_ {C _ {1}} u d \theta = \log r _ {2} \int_ {C _ {2}} r _ {2} \frac {\partial u}{\partial r} d \theta - \int_ {C _ {2}} u d \theta .
\end{equation*}
In other words, the expression
\begin{equation*}
\int_ {| z | = r} u d \theta - \log r \int_ {| z | = r} r \frac {\partial u}{\partial r} d \theta
\end{equation*}
is constant, and this is true even if u is only known to be harmonic in an annulus. By \hyperref[eq:4.58]{(58)} we find in the same way that
\begin{equation*}
\int_ {| z | = r} r \frac {\partial u}{\partial r} d \theta
\end{equation*}
is constant in the case of an annulus and zero if u is harmonic in the whole disk. Combining these results we obtain:

\begin{originalstatement}{Theorem}{20}\label{thm:4.20}
The arithmetic mean of a harmonic function\index{Function!harmonic}\index{Harmonic function} over concentric circles $|z| = r$ is a linear function of $\log r$,
\begin{equation*}
\frac {1}{2 \pi} \int_ {| z | = r} u d \theta = \alpha \log r + \beta ,\tag{61}\label{eq:4.61}
\end{equation*}
and if $u$ is harmonic in a disk $\alpha = 0$ and the arithmetic mean is constant.

\end{originalstatement}
In the latter case $\beta = u(0)$, by continuity, and changing to a new origin we find
\begin{equation*}
u (z _ {0}) = \frac {1}{2 \pi} \int_ {0} ^ {2 \pi} u (z _ {0} + r e ^ {i \theta}) d \theta .\tag{62}\label{eq:4.62}
\end{equation*}
It is clear that \hyperref[eq:4.62]{(62)} could also have been derived from the corre%
% Original PDF physical page 182
\sourcepage{182}%
sponding\index{Poisson formula|(} formula for analytic functions, Sec.~\hyperref[sec:4.3.4]{3.4}, \hyperref[eq:4.34]{(34)}. It leads directly to the maximum principle\index{Maximum principle} for harmonic functions\index{Function!harmonic}\index{Harmonic function}:

\begin{originalstatement}{Theorem}{21}\label{thm:4.21}
A nonconstant harmonic function has neither a maximum nor a minimum in its region of definition. Consequently, the maximum and the minimum on a closed bounded set E are taken on the boundary of E.
\end{originalstatement}
The proof is the same as for the maximum principle of analytic functions and will not be repeated. It applies also to the minimum for the reason that -u is harmonic together with u. In the case of analytic functions the corresponding procedure would have been to apply the maximum principle to $1/f(z)$ which is illegitimate unless $f(z) \neq 0$. Observe that the maximum principle for analytic functions follows from the maximum principle for harmonic functions by applying the latter to $\log |f(z)|$ which is harmonic when $f(z) \neq 0$.

\subsection*{EXERCISES}
\originalnumber{1}{exercise:4.6.2.11.1}
1. If u is harmonic and bounded in $0 < |z| < \rho$, show that the origin is a removable singularity in the sense that u becomes harmonic in $|z| < \rho$ when $u(0)$ is properly defined.

\originalnumber{2}{exercise:4.6.2.11.2}
2. Suppose that $f(z)$ is analytic in the annulus $r_{1} < |z| < r_{2}$ and continuous on the closed annulus. If $M(r)$ denotes the maximum of $|f(z)|$ for $|z| = r$, show that
\begin{equation*}
M (r) \leq M (r _ {1}) ^ {\alpha} M (r _ {2}) ^ {1 - \alpha}
\end{equation*}
where $\alpha = \log (r_{2}/r)$: $\log (r_{2}/r_{1})$ (Hadamard's three-circle theorem\index{Hadamard's three-circle theorem}). Discuss cases of equality. Hint: Apply the maximum principle to a linear combination of $\log |f(z)|$ and $\log |z|$.

\setcounter{section}{6}\setcounter{subsection}{2}\subsection{Poisson's Formula}\label{sec:4.6.3}
The maximum principle has the following important consequence: If $u(z)$ is continuous on a closed bounded set $E$ and harmonic on the interior of $E$, then it is uniquely determined by its values on the boundary of $E$. Indeed, if $u_1$ and $u_2$ are two such functions with the same boundary values, then $u_1 - u_2$ is harmonic with the boundary values 0. By the maximum and minimum principle the difference $u_1 - u_2$ must then be identically zero on $E$.

There arises the problem of finding u when its boundary values are given. At this point we shall solve the problem only in the simplest case, namely for a closed disk.

Formula \hyperref[eq:4.62]{(62)} determines the value of u at the center of the disk. But this is all we need, for there exists a linear transformation which carries
% Original PDF physical page 183
\index{Mean-value property|)}\sourcepage{183}
any point to the center. To be explicit, suppose that $u(z)$ is harmonic in the closed disk $|z| \leq R$. The linear transformation
\begin{equation*}
z = S (\zeta) = \frac {R (R \zeta + a)}{R + \bar {a} \zeta}
\end{equation*}
maps $|\zeta| \leq 1$ onto $|z| \leq R$ with $\zeta = 0$ corresponding to $z = a$. The function $u(S(\zeta))$ is harmonic in $|\zeta| \leq 1$, and by \hyperref[eq:4.62]{(62)} we obtain
\begin{equation*}
u (a) = \frac {1}{2 \pi} \int_ {| \zeta | = 1} u (S (\zeta)) d \arg \zeta .
\end{equation*}
From
\begin{equation*}
\zeta = \frac {R (z - a)}{R ^ {2} - \bar {a} z}
\end{equation*}
we compute
\begin{equation*}
\begin{aligned}
d\arg\zeta&=-i\frac{d\zeta}{\zeta}
=-i\left(\frac1{z-a}+\frac{\bar a}{R^2-\bar az}\right)dz\\
&=\left(\frac z{z-a}+\frac{\bar az}{R^2-\bar az}\right)d\theta.
\end{aligned}
\end{equation*}
On substituting $R^{2} = z\bar{z}$ the coefficient of $d\theta$ in the last expression can be rewritten as
\begin{equation*}
\frac {z}{z - a} + \frac {\bar {a}}{\bar {z} - \bar {a}} = \frac {R ^ {2} - | a | ^ {2}}{| z - a | ^ {2}}
\end{equation*}
or, equivalently, as
\begin{equation*}
\frac {1}{2} \left(\frac {z + a}{z - a} + \frac {\bar {z} + \bar {a}}{\bar {z} - \bar {a}}\right) = \mathrm{Re} \frac {z + a}{z - a}.
\end{equation*}
We obtain the two forms
\begin{equation*}
u (a) = \frac {1}{2 \pi} \int_ {| z | = R} \frac {R ^ {2} - | a | ^ {2}}{| z - a | ^ {2}} u (z) d \theta = \frac {1}{2 \pi} \int_ {| z | = R} \operatorname{Re} \frac {z + a}{z - a} u (z) d \theta\tag{63}\label{eq:4.63}
\end{equation*}
of Poisson's formula. In polar coordinates,
\begin{equation*}
u (r e ^ {i \varphi}) = \frac {1}{2 \pi} \int_ {0} ^ {2 \pi} \frac {R ^ {2} - r ^ {2}}{R ^ {2} - 2 r R \cos (\theta - \varphi) + r ^ {2}} u (R e ^ {i \theta}) d \theta .
\end{equation*}
In the derivation we have assumed that $u(z)$ is harmonic in the closed disk. However, the result remains true under the weaker condition that $u(z)$ is harmonic in the open disk and continuous in the closed disk. Indeed, if 0 < r < 1, then $u(rz)$ is harmonic in the closed disk, and we obtain
\begin{equation*}
u (r a) = \frac {1}{2 \pi} \int_ {| z | = R} \frac {R ^ {2} - | a | ^ {2}}{| z - a | ^ {2}} u (r z) d \theta .
\end{equation*}
% Original PDF physical page 184
\sourcepage{184}
Now all we need to do is to let $r$ tend to 1. Because $u(z)$ is uniformly continuous on $|z| \leq R$ it is true that $u(rz) \to u(z)$ uniformly for $|z| = R$, and we conclude that \hyperref[eq:4.63]{(63)} remains valid.

We shall formulate the result as a theorem:

\begin{originalstatement}{Theorem}{22}\label{thm:4.22}
Suppose that $u(z)$ is harmonic for $|z| < R$, continuous for $|z| \leq R$. Then
\begin{equation*}
u (a) = \frac {1}{2 \pi} \int_ {| z | = R} \frac {R ^ {2} - | a | ^ {2}}{| z - a | ^ {2}} u (z) d \theta\tag{64}\label{eq:4.64}
\end{equation*}
for all $|a| < R$.

\end{originalstatement}
The theorem leads at once to an explicit expression for the conjugate function of u. Indeed, formula \hyperref[eq:4.63]{(63)} gives
\begin{equation*}
u (z) = \operatorname{Re} \left[ \frac {1}{2 \pi i} \int_ {| \zeta | = R} \frac {\zeta + z}{\zeta - z} u (\zeta) \frac {d \zeta}{\zeta} \right].\tag{65}\label{eq:4.65}
\end{equation*}
The bracketed expression is an analytic function of z for $|z| < R$. It follows that $u(z)$ is the real part of
\begin{equation*}
f (z) = \frac {1}{2 \pi i} \int_ {| \zeta | = R} \frac {\zeta + z}{\zeta - z} u (\zeta) \frac {d \zeta}{\zeta} + i C\tag{66}\label{eq:4.66}
\end{equation*}
where $C$ is an arbitrary real constant. This formula is known as Schwarz's formula.

As a special case of \hyperref[eq:4.64]{(64)}, note that $u = 1$ yields
\begin{equation*}
\int_ {| z | = R} \frac {R ^ {2} - | z | ^ {2}}{| z - a | ^ {2}} d \theta = 2 \pi\tag{67}\label{eq:4.67}
\end{equation*}
for all $|a| < R$.

\setcounter{section}{6}\setcounter{subsection}{3}\subsection{Schwarz's Theorem}\label{sec:4.6.4}
Theorem~\hyperref[thm:4.22]{22} serves to express a given harmonic function\index{Function!harmonic}\index{Harmonic function} through its values on a circle. But the right-hand side of formula \hyperref[eq:4.64]{(64)} has a meaning as soon as u is defined on $|z| = R$, provided it is sufficiently regular, for instance piecewise continuous. As in \hyperref[eq:4.65]{(65)} the integral can again be written as the real part of an analytic function, and consequently it is a harmonic function. The question is, does it have the boundary values $u(z)$ on $|z| = R$?

There is reason to clarify the notations. Choosing $R = 1$ we define, for any piecewise continuous function $U(\theta)$ in $0 \leq \theta \leq 2\pi$,


% Original PDF physical page 185
\index{Poisson formula|)}\sourcepage{185}
\begin{equation*}
P _ {U} (z) = \frac {1}{2 \pi} \int_ {0} ^ {2 \pi} \mathrm{Re} \frac {e ^ {i \theta} + z}{e ^ {i \theta} - z} U (\theta) d \theta
\end{equation*}\index{Schwarz H A@Schwarz, H. A.!theorem proved by|(}
and call this the Poisson integral of U. Observe that $P_{U}(z)$ is not only a function of z, but also a function of the function U; as such it is called a functional\index{Functional}. The functional is linear inasmuch as
\begin{equation*}
P _ {U + V} = P _ {U} + P _ {V}
\end{equation*}
and
\begin{equation*}
P _ {c U} = c P _ {U}
\end{equation*}
for constant c. Moreover, $U \geq 0$ implies $P_{U}(z) \geq 0$; because of this property $P_{U}$ is said to be a positive linear functional.

We deduce from \hyperref[eq:4.67]{(67)} that $P_{c} = c$. From this property, together with the linear and positive character of the functional, it follows that any inequality $m \leq U \leq M$ implies $m \leq P_{U} \leq M$.

The question of boundary values is settled by the following fundamental theorem that was first proved by H. A. Schwarz:

\begin{originalstatement}{Theorem}{23}\label{thm:4.23}
The function $P_U(z)$ is harmonic for $|z| < 1$, and
\begin{equation*}
\lim _ {z \rightarrow e ^ {i \theta_ {0}}} P _ {U} (z) = U (\theta_ {0})\tag{68}\label{eq:4.68}
\end{equation*}
provided that U is continuous at $\theta_{0}$.

\end{originalstatement}
We have already remarked that $P_{U}$ is harmonic. To study the boundary behavior, let $C_{1}$ and $C_{2}$ be complementary arcs of the unit circle, and denote by $U_{1}$ the function which coincides with U on $C_{1}$ and vanishes on $C_{2}$, by $U_{2}$ the corresponding function for $C_{2}$. Clearly, $P_{U} = P_{U_{1}} + P_{U_{2}}$.

Since $P_{U_{1}}$ can be regarded as a line integral over $C_{1}$ it is, by the same reasoning as before, harmonic everywhere except on the closed arc $C_{1}$. The expression
\begin{equation*}
\mathrm{Re} \frac {e ^ {i \theta} + z}{e ^ {i \theta} - z} = \frac {1 - | z | ^ {2}}{| e ^ {i \theta} - z | ^ {2}}
\end{equation*}
vanishes on $|z| = 1$ for $z \neq e^{i\theta}$. It follows that $P_{U_1}$ is zero on the open arc $C_2$, and since it is continuous $P_{U_1}(z) \to 0$ as $z \to e^{i\theta} \in C_2$.

In proving \hyperref[eq:4.68]{(68)} we may suppose that $U(\theta_0) = 0$, for if this is not the case we need only replace $U$ by $U - U(\theta_0)$. Given $\varepsilon > 0$ we can find $C_1$ and $C_2$ such that $e^{i\theta_0}$ is an interior point of $C_2$ and $|U(\theta)| < \varepsilon / 2$ for $e^{i\theta} \in C_2$. Under this condition $|U_2(\theta)| < \varepsilon / 2$ for all $\theta$, and hence $|P_{U_2}(z)| < \varepsilon / 2$ for all $|z| < 1$. On the other hand, since $U_1$ is continuous and vanishes


% Original PDF physical page 186
\index{Schwarz H A@Schwarz, H. A.!theorem proved by|)}\sourcepage{186}
\bookfigure{186}{1}{FIG. 4-15}
at $e^{i\theta_0}$ there exists a $\delta$ such that $|P_{U_1}(z)| < \varepsilon / 2$ for $|z - e^{i\theta_0}| < \delta$. It follows that $|P_U(z)| \leq |P_{U_1}| + |P_{U_2}| < \varepsilon$ as soon as $|z| < 1$ and $|z - e^{i\theta_0}| < \delta$, which is precisely what we had to prove.

There is an interesting geometric interpretation of Poisson's formula, also due to Schwarz. Given a fixed $z$ inside the unit circle we determine for each $e^{i\theta}$ the point $e^{i\theta^*}$ which is such that $e^{i\theta}, z$ and $e^{i\theta^*}$ are in a straight line (Fig.~\hyperref[fig:4-15]{4-15}). It is clear geometrically, or by simple calculation, that
\begin{equation*}
1 - | z | ^ {2} = \left| e ^ {i \theta} - z \right| \left| e ^ {i \theta *} - z \right|.\tag{69}\label{eq:4.69}
\end{equation*}
But the ratio $(e^{i\theta} - z) / (e^{i\theta^*} - z)$ is negative, so we must have
\begin{equation*}
1 - | z | ^ {2} = - (e ^ {i \theta} - z) (e ^ {- i \theta^ {*}} - \bar {z}).
\end{equation*}
We regard $\theta^{*}$ as a function of $\theta$ and differentiate. Since $z$ is constant we obtain
\begin{equation*}
\frac {e ^ {i \theta} d \theta}{e ^ {i \theta} - z} = \frac {e ^ {- i \theta *} d \theta^ {*}}{e ^ {- i \theta *} - \bar {z}}
\end{equation*}
and, on taking absolute values,
\begin{equation*}
\frac {d \theta^ {*}}{d \theta} = \left| \frac {e ^ {i \theta^ {*}} - z}{e ^ {i \theta} - z} \right|.\tag{70}\label{eq:4.70}
\end{equation*}
It follows by \hyperref[eq:4.69]{(69)} and \hyperref[eq:4.70]{(70)} that
\begin{equation*}
\frac {1 - | z | ^ {2}}{| e ^ {i \theta} - z | ^ {2}} = \frac {d \theta^ {*}}{d \theta}
\end{equation*}
and hence
\begin{equation*}
P _ {U} (z) = \frac {1}{2 \pi} \int_ {0} ^ {2 \pi} U (\theta) d \theta^ {*} = \frac {1}{2 \pi} \int_ {0} ^ {2 \pi} U (\theta^ {*}) d \theta .
\end{equation*}
% Original PDF physical page 187
\sourcepage{187}
In other words, to find $P_{U}(z)$, replace each value of $U(\theta)$ by the value at the point opposite to z, and take the average over the circle.

\subsection*{EXERCISES}
\originalnumber{1}{exercise:4.6.4.12.1}
1. Assume that $U(\xi)$ is piecewise continuous and bounded for all real $\xi$. Show that
\begin{equation*}
P _ {U} (z) = \frac {1}{\pi} \int_ {- \infty} ^ {\infty} \frac {y}{(x - \xi) ^ {2} + y ^ {2}} U (\xi) d \xi
\end{equation*}
represents a harmonic function\index{Function!harmonic}\index{Harmonic function} in the upper half plane with boundary values $U(\xi)$ at points of continuity (Poisson's integral for the half plane).

\originalnumber{2}{exercise:4.6.4.12.2}
2. Prove that a function which is harmonic and bounded in the upper half plane, continuous on the real axis, can be represented as a Poisson integral (Ex.~\hyperref[exercise:4.6.4.12.1]{1}).

Remark. The point at $\infty$ presents an added difficulty, for we cannot immediately apply the maximum and minimum principle to $u - P_u$. A good try would be to apply the maximum principle to $u - P_u - \varepsilon y$ for $\varepsilon > 0$, with the idea of letting $\varepsilon$ tend to 0. This almost works, for the function tends to 0 for $y \to 0$ and to $-\infty$ for $y \to \infty$, but we lack control when $|x| \to \infty$. Show that the reasoning can be carried out successfully by application to $u - P_u - \varepsilon\operatorname{Im}(\sqrt{iz})$.

\originalnumber{3}{exercise:4.6.4.12.3}
3. In Ex.~\hyperref[exercise:4.6.4.12.1]{1}, assume that $U$ has a jump at 0, for instance $U(+0) = 0$, $U(-0) = 1$. Show that $P_U(z) - \frac{1}{\pi} \arg z$ tends to 0 as $z \to 0$. Generalize to arbitrary jumps and to the case of the circle.

\originalnumber{4}{exercise:4.6.4.12.4}
4. If $C_{1}$ and $C_{2}$ are complementary arcs on the unit circle, set U = 1 on $C_{1}$, U = 0 on $C_{2}$. Find $P_{U}(z)$ explicitly and show that $2\pi P_{U}(z)$ equals the length of the arc, opposite to $C_{1}$, cut off by the straight lines through z and the end points of $C_{1}$.

\originalnumber{5}{exercise:4.6.4.12.5}
5. Show that the mean-value formula \hyperref[eq:4.62]{(62)} remains valid for $u = \log |1 + z|$, $z_0 = 0$, $r = 1$, and use this fact to compute
\begin{equation*}
\int_ {0} ^ {\pi} \log \sin \theta d \theta .
\end{equation*}
\originalnumber{6}{exercise:4.6.4.12.6}
6. If $f(z)$ is analytic in the whole plane and if $z^{-1} \operatorname{Re} f(z) \to 0$ when $z \to \infty$, show that $f$ is a constant. Hint: Use \hyperref[eq:4.66]{(66)}.

\originalnumber{7}{exercise:4.6.4.12.7}
7. If $f(z)$ is analytic in a neighborhood of $\infty$ and if $z^{-1} \operatorname{Re} f(z) \to 0$ when $z \to \infty$, show that $\lim_{z \to \infty} f(z)$ exists. (In other words, the isolated singularity at $\infty$ is removable.)

Hint: Show first, by use of Cauchy's integral formula, that $f = f_{1} + f_{2}$ where $f_{1}(z) \to 0$ for $z \to \infty$ and $f_{2}(z)$ is analytic in the whole plane.


% Original PDF physical page 188
\sourcepage{188}
\originalnumber{8}{exercise:4.6.4.12.8}
*8.\index{Reflection principal|(} If $u(z)$ is harmonic for $0 < |z| < \rho$ and $\lim_{z \to 0} zu(z) = 0$, prove that $u$ can be written in the form $u(z) = \alpha \log |z| + u_0(z)$ where $\alpha$ is a constant and $u_0$ is harmonic in $|z| < \rho$.

Hint: Choose $\alpha$ as in \hyperref[eq:4.61]{(61)}. Then show that $u_{0}$ is the real part of an analytic function $f_{0}(z)$ and use the preceding exercise to conclude that the singularity at 0 is removable.

\setcounter{section}{6}\setcounter{subsection}{4}\subsection{The Reflection Principle}\label{sec:4.6.5}
An elementary aspect of the symmetry principle\index{Symmetry principle}, or reflection principle, has been discussed already in connection with linear transformations (Chap.~\hyperref[chap:3]{3}, Sec.~\hyperref[sec:3.3.3]{3.3}). There are many more general variants first formulated by H. A. Schwarz.

The principle of reflection is based on the observation that if $u(z)$ is a harmonic function\index{Function!harmonic}\index{Harmonic function}, then $u(\bar{z})$ is likewise harmonic, and if $f(z)$ is an analytic function, then $\overline{f(\bar{z})}$ is also analytic. More precisely, if $u(z)$ is harmonic and $f(z)$ analytic in a region then $u(\bar{z})$ is harmonic and $\overline{f(\bar{z})}$ analytic as functions of z in the region $\Omega^{*}$ obtained by reflecting $\Omega$ in the real axis; that is, $z \in \Omega^{*}$ if and only if $\bar{z} \in \Omega$. The proofs of these statements consist in trivial verifications.

Consider the case of a symmetric region: $\Omega^{*} = \Omega$. Because $\Omega$ is connected it must intersect the real axis along at least one open interval. Assume now that $f(z)$ is analytic in $\Omega$ and real on at least one interval of the real axis. Since $f(z) - \overline{f(\bar{z})}$ is analytic and vanishes on an interval it must be identically zero, and we conclude that $f(z) = \overline{f(\bar{z})}$ in $\Omega$. With the notation $f = u + iv$ we have thus $u(z) = u(\bar{z})$, $v(z) = -v(\bar{z})$.

This is important, but it is a rather weak result, for we are assuming that $f(z)$ is already known to be analytic in all of $\Omega$. Let us denote the intersection of $\Omega$ with the upper half plane by $\Omega^{+}$, and the intersection of $\Omega$ with the real axis by $\sigma$. Suppose that $f(z)$ is defined on $\Omega^{+} \cup \sigma$, analytic in $\Omega^{+}$, continuous and real on $\sigma$. Under these conditions we want to show that $f(z)$ is the restriction to $\Omega^{+}$ of a function which is analytic in all of $\Omega$ and satisfies the symmetry condition $f(z) = \overline{f(\bar{z})}$. In other words, part of our theorem asserts that $f(z)$ has an analytic continuation\index{Analytic continuation} to $\Omega$.

Even in this formulation the assumptions are too strong. Indeed, the main thing is that the imaginary part $v(z)$ vanishes on $\sigma$, and nothing at all need to be assumed about the real part. In the definitive statement of the reflection principle the emphasis should therefore be on harmonic functions.

\begin{originalstatement}{Theorem}{24}\label{thm:4.24}
Let $\Omega^{+}$ be the part in the upper half plane of a symmetric region $\Omega$, and let $\sigma$ be the part of the real axis in $\Omega$. Suppose that $v(x)$ is continuous in $\Omega^{+} \cup \sigma$, harmonic in $\Omega^{+}$, and zero on $\sigma$. Then $v$ has a harmonic extension to $\Omega$ which satisfies the symmetry relation $v(\bar{z}) = -v(z)$.
\par
% Original PDF physical page 189
\sourcepage{189}
In the same situation, if $v$ is the imaginary part of an analytic function $f(z)$ in $\Omega^{+}$, then $f(z)$ has an analytic extension which satisfies $f(z) = \overline{f(\overline{z})}$.
\end{originalstatement}

For the proof we construct the function $V(z)$ which is equal to $v(z)$ in $\Omega^{+}$, 0 on $\sigma$, and equal to $-v(\bar{z})$ in the mirror image of $\Omega^{+}$. We have to show that V is harmonic on $\sigma$. For a point $x_{0} \in \sigma$ consider a disk with center $x_{0}$ contained in $\Omega$, and let $P_{V}$ denote the Poisson integral with respect to this disk formed with the boundary values V. The difference $V - P_{V}$ is harmonic in the upper half of the disk. It vanishes on the half circle, by Theorem~\hyperref[thm:4.23]{23}, and also on the diameter, because V tends to zero by definition and $P_{V}$ vanishes by obvious symmetry. The maximum and minimum principle implies that $V = P_{V}$ in the upper half disk, and the same proof can be repeated for the lower half. We conclude that V is harmonic in the whole disk, and in particular at $x_{0}$.

For the remaining part of the theorem, let us again consider a disk with center on $\sigma$. We have already extended v to the whole disk, and v has a conjugate harmonic function\index{Function!harmonic}\index{Harmonic function} $-u_{0}$ in the same disk which we may normalize so that $u_{0} = \operatorname{Re} f(z)$ in the upper half. Consider
\begin{equation*}
U _ {0} (z) = u _ {0} (z) - u _ {0} (\bar {z}).
\end{equation*}
On the real diameter it is clear that $\partial U_0 / \partial x = 0$ and also
\begin{equation*}
\frac {\partial U _ {0}}{\partial y} = 2 \frac {\partial u _ {0}}{\partial y} = - 2 \frac {\partial v}{\partial x} = 0.
\end{equation*}
It follows that the analytic function $\partial U_{0}/\partial x - i \partial U_{0}/\partial y$ vanishes on the real axis, and hence identically. Therefore $U_{0}$ is a constant, and this constant is evidently zero. We have proved that $u_{0}(z) = u_{0}(\bar{z})$.

The construction can be repeated for arbitrary disks. It is clear that the $u_{0}$ coincide in overlapping disks. The definition can be extended to all of $\Omega$, and the theorem follows.

The theorem has obvious generalizations. The domain $\Omega$ can be taken to be symmetric with respect to a circle $C$ rather than with respect to a straight line, and when $z$ tends to $C$ it may be assumed that $f(z)$ approaches another circle $C'$. Under such conditions $f(z)$ has an analytic continuation which maps symmetric points with respect to $C$ onto symmetric points with respect to $C'$.

\subsection*{EXERCISES}
\originalnumber{1}{exercise:4.6.5.13.1}
1. If $f(z)$ is analytic in the whole plane and real on the real axis, purely imaginary on the imaginary axis, show that $f(z)$ is odd.

\originalnumber{2}{exercise:4.6.5.13.2}
2. Show that every function $f$ which is analytic in a symmetric region $\Omega$ can be written in the form $f_1 + if_2$ where $f_1, f_2$ are analytic in $\Omega$ and


% Original PDF physical page 190
\index{Integration|)}\sourcepage{190}
real on the real axis.

\originalnumber{3}{exercise:4.6.5.13.3}
3. If $f(z)$ is analytic in $|z| \leq 1$ and satisfies $|f| = 1$ on $|z| = 1$, show that $f(z)$ is rational.

\originalnumber{4}{exercise:4.6.5.13.4}
4. Use \hyperref[eq:4.66]{(66)} to derive a formula for $f'(z)$ in terms of $u(z)$.

\originalnumber{5}{exercise:4.6.5.13.5}
5. If $u(z)$ is harmonic and $0 \leq u(z) \leq Ky$ for $y > 0$, prove that $u = ky$ with $0 \leq k \leq K$. [Reflect over the real axis, complete to an analytic function $f(z) = u + iv$, and use Ex.~\hyperref[exercise:4.6.5.13.4]{4} to show that $f'(z)$ is bounded.]

\index{Function!harmonic|)}\index{Harmonic function|)}\index{Reflection principal|)}